% Options for packages loaded elsewhere
\PassOptionsToPackage{unicode}{hyperref}
\PassOptionsToPackage{hyphens}{url}
\PassOptionsToPackage{dvipsnames,svgnames,x11names}{xcolor}
\documentclass[
  english,
  oneside]{book}
\usepackage{xcolor}
\usepackage[margin=1in]{geometry}
\usepackage{amsmath,amssymb}
\setcounter{secnumdepth}{-\maxdimen} % remove section numbering
\usepackage{iftex}
\ifPDFTeX
  \usepackage[T1]{fontenc}
  \usepackage[utf8]{inputenc}
  \usepackage{textcomp} % provide euro and other symbols
\else % if luatex or xetex
  \usepackage{unicode-math} % this also loads fontspec
  \defaultfontfeatures{Scale=MatchLowercase}
  \defaultfontfeatures[\rmfamily]{Ligatures=TeX,Scale=1}
\fi
\usepackage{lmodern}
\ifPDFTeX\else
  % xetex/luatex font selection
\fi
% Use upquote if available, for straight quotes in verbatim environments
\IfFileExists{upquote.sty}{\usepackage{upquote}}{}
\IfFileExists{microtype.sty}{% use microtype if available
  \usepackage[]{microtype}
  \UseMicrotypeSet[protrusion]{basicmath} % disable protrusion for tt fonts
}{}
\makeatletter
\@ifundefined{KOMAClassName}{% if non-KOMA class
  \IfFileExists{parskip.sty}{%
    \usepackage{parskip}
  }{% else
    \setlength{\parindent}{0pt}
    \setlength{\parskip}{6pt plus 2pt minus 1pt}}
}{% if KOMA class
  \KOMAoptions{parskip=half}}
\makeatother
\usepackage{color}
\usepackage{fancyvrb}
\newcommand{\VerbBar}{|}
\newcommand{\VERB}{\Verb[commandchars=\\\{\}]}
\DefineVerbatimEnvironment{Highlighting}{Verbatim}{commandchars=\\\{\}}
% Add ',fontsize=\small' for more characters per line
\usepackage{framed}
\definecolor{shadecolor}{RGB}{30,30,30}
\newenvironment{Shaded}{\begin{snugshade}}{\end{snugshade}}
\newcommand{\AlertTok}[1]{\textcolor[rgb]{0.83,0.83,0.83}{#1}}
\newcommand{\AnnotationTok}[1]{\textcolor[rgb]{0.83,0.83,0.83}{#1}}
\newcommand{\AttributeTok}[1]{\textcolor[rgb]{0.83,0.83,0.83}{#1}}
\newcommand{\BaseNTok}[1]{\textcolor[rgb]{0.71,0.81,0.66}{#1}}
\newcommand{\BuiltInTok}[1]{\textcolor[rgb]{0.83,0.83,0.83}{#1}}
\newcommand{\CharTok}[1]{\textcolor[rgb]{0.83,0.83,0.83}{#1}}
\newcommand{\CommentTok}[1]{\textcolor[rgb]{0.42,0.60,0.33}{\textit{#1}}}
\newcommand{\CommentVarTok}[1]{\textcolor[rgb]{0.83,0.83,0.83}{#1}}
\newcommand{\ConstantTok}[1]{\textcolor[rgb]{0.31,0.76,1.00}{#1}}
\newcommand{\ControlFlowTok}[1]{\textcolor[rgb]{0.77,0.53,0.75}{\textbf{#1}}}
\newcommand{\DataTypeTok}[1]{\textcolor[rgb]{0.31,0.79,0.69}{#1}}
\newcommand{\DecValTok}[1]{\textcolor[rgb]{0.71,0.81,0.66}{#1}}
\newcommand{\DocumentationTok}[1]{\textcolor[rgb]{0.83,0.83,0.83}{#1}}
\newcommand{\ErrorTok}[1]{\textcolor[rgb]{0.83,0.83,0.83}{#1}}
\newcommand{\ExtensionTok}[1]{\textcolor[rgb]{0.83,0.83,0.83}{#1}}
\newcommand{\FloatTok}[1]{\textcolor[rgb]{0.71,0.81,0.66}{#1}}
\newcommand{\FunctionTok}[1]{\textcolor[rgb]{0.86,0.86,0.67}{#1}}
\newcommand{\ImportTok}[1]{\textcolor[rgb]{0.83,0.83,0.83}{#1}}
\newcommand{\InformationTok}[1]{\textcolor[rgb]{0.83,0.83,0.83}{#1}}
\newcommand{\KeywordTok}[1]{\textcolor[rgb]{0.34,0.61,0.84}{\textbf{#1}}}
\newcommand{\NormalTok}[1]{\textcolor[rgb]{0.83,0.83,0.83}{#1}}
\newcommand{\OperatorTok}[1]{\textcolor[rgb]{0.83,0.83,0.83}{#1}}
\newcommand{\OtherTok}[1]{\textcolor[rgb]{0.83,0.83,0.83}{#1}}
\newcommand{\PreprocessorTok}[1]{\textcolor[rgb]{0.83,0.83,0.83}{#1}}
\newcommand{\RegionMarkerTok}[1]{\textcolor[rgb]{0.83,0.83,0.83}{#1}}
\newcommand{\SpecialCharTok}[1]{\textcolor[rgb]{0.84,0.73,0.49}{#1}}
\newcommand{\SpecialStringTok}[1]{\textcolor[rgb]{0.83,0.83,0.83}{#1}}
\newcommand{\StringTok}[1]{\textcolor[rgb]{0.81,0.57,0.47}{#1}}
\newcommand{\VariableTok}[1]{\textcolor[rgb]{0.83,0.83,0.83}{#1}}
\newcommand{\VerbatimStringTok}[1]{\textcolor[rgb]{0.83,0.83,0.83}{#1}}
\newcommand{\WarningTok}[1]{\textcolor[rgb]{0.83,0.83,0.83}{#1}}
\usepackage{longtable,booktabs,array}
\usepackage{caption}
\captionsetup[table]{skip=6pt}
\newcounter{none} % for unnumbered tables
\usepackage{calc} % for calculating minipage widths
% Correct order of tables after \paragraph or \subparagraph
\usepackage{etoolbox}
\makeatletter
\patchcmd\longtable{\par}{\if@noskipsec\mbox{}\fi\par}{}{}
\makeatother
\ifLuaTeX
\usepackage[bidi=basic,shorthands=off]{babel}
\else
\usepackage[bidi=default,shorthands=off]{babel}
\fi
\ifLuaTeX
  \usepackage{selnolig} % disable illegal ligatures
\fi
\setlength{\emergencystretch}{3em} % prevent overfull lines
\providecommand{\tightlist}{%
  \setlength{\itemsep}{0pt}\setlength{\parskip}{0pt}}
\usepackage{fontspec}
\usepackage{graphicx}
\setmonofont[Scale=0.86]{Cascadia Mono}
\graphicspath{{logo/}{../logo/}}
\setcounter{secnumdepth}{-2}
\IfFileExists{fvextra.sty}{\usepackage{fvextra}\DefineVerbatimEnvironment{Highlighting}{Verbatim}{breaklines,commandchars=\\\{\},fontsize=\small}}{}
\makeatletter
\renewcommand{\maketitle}{\begin{titlepage}\centering\vspace*{1.4cm}\includegraphics[width=0.30\textwidth]{iris-logo.png}\par\vspace{1.1cm}{\Huge\bfseries\@title\par}\vspace{0.5cm}{\large\color[HTML]{374151}Intermediate Representation for Intelligent Systems --- Version 1.0.0-rc2\par}\vspace{0.9cm}{\color[HTML]{005EC3}\rule{0.45\textwidth}{1.6pt}}\par\vspace{1.0cm}{\Large\@author\par}\vspace{0.35cm}{\normalsize\color[HTML]{4B5563}\@date\par}\vfill{\small\color[HTML]{6B7280}Complete Language Guide \& Reference\par}\vspace{0.3cm}{\small\color[HTML]{6B7280}\url{https://github.com/Eclipse-Softworks/IRIS}\par}\end{titlepage}}
\makeatother
\usepackage{bookmark}
\IfFileExists{xurl.sty}{\usepackage{xurl}}{} % add URL line breaks if available
\urlstyle{same}
% fallback for those not using the hyperref driver hyperxmp:
\makeatletter
\@ifundefined{xmpquote}{\newcommand{\xmpquote}[1]{#1}}{}
\makeatother
\hypersetup{
  pdftitle={The IRIS Programming Language},
  pdfauthor={Moon9t},
  pdflang={en},
  colorlinks=true,
  linkcolor={blue},
  filecolor={Maroon},
  citecolor={Blue},
  urlcolor={blue},
  pdfcreator={LaTeX via pandoc}}

\title{The IRIS Programming Language}
\author{Moon9t}
\date{November 2026}

\begin{document}
\frontmatter
\maketitle

{
\hypersetup{linkcolor=blue}
\setcounter{tocdepth}{1}
\tableofcontents
}
\mainmatter
\chapter{Copyright and License}\label{copyright-and-license}

\textbf{The IRIS Programming Language}

Copyright © 2024--2026 Moon9t. All rights reserved.

This book is licensed under the terms of the GNU General Public License
as published by the Free Software Foundation; either version 2 of the
License, or (at your option) any later version.

This documentation is distributed in the hope that it will be useful,
but WITHOUT ANY WARRANTY; without even the implied warranty of
MERCHANTABILITY or FITNESS FOR A PARTICULAR PURPOSE. See the GNU General
Public License for more details.

You should have received a copy of the GNU General Public License along
with this book. If not, see \url{https://www.gnu.org/licenses/}.

\newpage

\chapter{Foreword}\label{foreword}

Modern intelligent software rarely lives inside a single tier of
abstractions. A single system routinely combines numerical optimization
or automatic differentiation, asynchronous networking, SQLite
persistence, hardware sensor loops or live camera feeds, a responsive
native graphical interface, and strict deployment safety guarantees.
Historically, teams built each layer in a separate language with fragile
glue between them.

\textbf{IRIS} (\emph{Intermediate Representation for Intelligent
Systems}) unifies those layers behind one statically typed language, one
block-parameter SSA intermediate representation, one deterministic
optimization pipeline, and one unified toolchain:

\begin{itemize}
\tightlist
\item
  \textbf{Native execution and JIT}: compiles through LLVM
  (\texttt{clang} / \texttt{ld.lld} and LLVM-C/ORC) to standalone
  machine-code executables, shared libraries, and desktop applications.
\item
  \textbf{Bounded interpreter and time-travel debugging}: executes the
  same checked SSA IR directly for rapid prototyping, trace-based
  step-back debugging (\texttt{iris\ dap}), and safe host-governed
  candidate evaluation.
\item
  \textbf{Embedded C-free cross-compilation}: lowers a verified
  allocation-free scalar subset directly to \texttt{.o} object bundles
  for \texttt{cortex-m4f}, \texttt{cortex-m33}, \texttt{esp32-c3},
  \texttt{esp32}, and \texttt{arduino-uno}.
\item
  \textbf{Differentiable computing and autonomy}: provides
  language-level reverse-mode autodiff (\texttt{tape},
  \texttt{backward}, \texttt{grad}), differentiable tensors, a
  declarative \texttt{model} DSL with ONNX/TensorRT export, ROS 2
  robotics integration, and 7-gate governed self-evolution.
\item
  \textbf{Native desktop applications}: ships an embedded CPU-rendered
  cross-platform UI stack (\texttt{std.ui*} paired with
  \texttt{iris\_ui\_host}) with responsive flexbox/grid layout, 6
  built-in themes, 99 vector SVG icons, spring animations, overlays,
  custom vector painting, live native camera capture, and single-command
  application bundling (\texttt{iris\ bundle}).
\end{itemize}

\section{Who This Book Is For}\label{who-this-book-is-for}

This book is written for programmers comfortable with general concepts
in at least one systems, scientific, or application language (such as
Rust, Python, C, C++, Go, or TypeScript). You do not need prior compiler
or machine-learning experience.

\begin{itemize}
\tightlist
\item
  \textbf{If you are new to IRIS}, read \textbf{Part I} and \textbf{Part
  II} in order (Chapters 1--8). They establish the type system,
  algebraic data modeling, pattern matching, traits, error handling,
  ownership, and memory discipline.
\item
  \textbf{If you are building services, embedded firmware, or concurrent
  pipelines}, read \textbf{Part III} (Chapters 9--12).
\item
  \textbf{If you are working on optimization, neural models, robotics,
  or live evolution}, read \textbf{Part IV} (Chapters 13--15).
\item
  \textbf{If you are building standalone desktop GUI applications or
  vision tools}, read \textbf{Part V} (Chapters 16--17).
\item
  \textbf{For day-to-day workflows and real system designs}, use
  \textbf{Part VI} (Chapters 18--19) and \textbf{Appendices A--E} as
  your workbench reference.
\end{itemize}

Throughout this book, maturity boundaries and hardware-specific
requirements are stated explicitly inline so you always know what runs
identically across platforms, what uses host-specific OS integrations,
and where physical hardware validation or optional third-party SDKs
apply.

\begin{center}\rule{0.5\linewidth}{0.5pt}\end{center}

\part{Part I: Foundations}\label{part-i-foundations}

\chapter{Chapter 1: Getting Started}\label{chapter-1-getting-started}

\section{1.1 What Is IRIS?}\label{what-is-iris}

IRIS is a statically typed, compiled programming language whose front
end, IR optimizer, native LLVM backend, package manager, formatter,
documentation generator, LSP language server, DAP debugger, and
application bundler all live inside one \texttt{iris} executable.

Every \texttt{.iris} program travels through the same front-end and
mid-level IR pipeline before branching to a target backend:

\begin{Shaded}
\begin{Highlighting}[]
\NormalTok{Source (.iris) + brought std.* / local modules}
\NormalTok{       │}
\NormalTok{       ▼  Lexer ─► Parser ─► AST Lowerer ─► Type \& Effect Checks}
\NormalTok{       │}
\NormalTok{       ▼  20{-}Pass SSA Optimizer}
\NormalTok{          (CFG NLL Borrow Check, Drop Elaboration, Shape Check,}
\NormalTok{           Inline, ConstFold, TCO, ML Fusion, GC \& Frame{-}Pool Scope, Async)}
\NormalTok{       │}
\NormalTok{       ├──► Native AOT Executable / Library   (iris build / iris run / iris bundle)}
\NormalTok{       ├──► In{-}Process LLVM ORC JIT           ({-}{-}emit jit / governed hot{-}swap)}
\NormalTok{       ├──► Bounded SSA IR Interpreter        ({-}{-}emit eval / iris dap / tests)}
\NormalTok{       ├──► Bare{-}Metal Microcontroller Bundle (iris embedded {-}{-}target ...)}
\NormalTok{       └──► Graph / Model Exports             ({-}{-}emit onnx, onnx{-}binary, tensorrt, cuda{-}ptx)}
\end{Highlighting}
\end{Shaded}

Because both native codegen and the SSA interpreter consume the
post-lowered, pass-verified IR, compile-time type errors, borrow-checker
rejections (\texttt{E0382}), checked arithmetic traps (\texttt{E0409}),
and out-of-bounds array diagnostics (\texttt{E0403}) behave consistently
across execution modes.

\begin{center}\rule{0.5\linewidth}{0.5pt}\end{center}

\section{1.2 Installation and Toolchain
Health}\label{installation-and-toolchain-health}

\subsection{Building from Source}\label{building-from-source}

To build IRIS from the Eclipse Softworks repository with Cargo (requires
Rust 1.75+):

\begin{Shaded}
\begin{Highlighting}[]
\NormalTok{git clone https}\OperatorTok{://}\NormalTok{github}\OperatorTok{.}\FunctionTok{com}\OperatorTok{/}\NormalTok{Eclipse{-}Softworks}\OperatorTok{/}\NormalTok{IRIS}\OperatorTok{.}\FunctionTok{git}
\FunctionTok{cd}\NormalTok{ IRIS}
\NormalTok{cargo build }\OperatorTok{{-}{-}}\NormalTok{release }\OperatorTok{{-}{-}}\NormalTok{locked}
\end{Highlighting}
\end{Shaded}

On Windows, this produces
\texttt{target\textbackslash{}release\textbackslash{}iris.exe} (and on
Linux/macOS, \texttt{target/release/iris}). Place the binary in a stable
directory on your \texttt{PATH} (for example
\texttt{\$env:USERPROFILE\textbackslash{}.iris\textbackslash{}bin\textbackslash{}iris.exe}
on Windows, or \texttt{\textasciitilde{}/.local/bin/iris} on
Linux/macOS).

If you intend to build standalone desktop GUI applications
(\texttt{std.ui*}), also compile the native UI host library in the same
workspace:

\begin{Shaded}
\begin{Highlighting}[]
\NormalTok{cargo build }\OperatorTok{{-}{-}}\NormalTok{release }\OperatorTok{{-}}\NormalTok{p iris\_ui\_host}
\end{Highlighting}
\end{Shaded}

\subsection{\texorpdfstring{Verifying Prerequisites with
\texttt{iris\ doctor}}{Verifying Prerequisites with iris doctor}}\label{verifying-prerequisites-with-iris-doctor}

IRIS parsing, \texttt{iris\ check}, \texttt{iris\ fmt},
\texttt{iris\ repl}, \texttt{iris\ lsp}, \texttt{iris\ pkg}, and
interpreter execution are completely self-contained. Native builds
(\texttt{iris\ build}, \texttt{iris\ run}, \texttt{iris\ bundle},
\texttt{-\/-emit\ jit}) additionally require:

\begin{itemize}
\tightlist
\item
  A dynamically loadable \textbf{LLVM-C} shared library (release bundles
  ship LLVM \textbf{23.1.1}).
\item
  A linker (\texttt{lld} / \texttt{ld.lld}, or system \texttt{clang}
  fallback).
\item
  Target C system libraries (MinGW \textbf{UCRT64} sysroot on Windows,
  \texttt{glibc} + \texttt{libm}/\texttt{pthreads} on Linux, Xcode
  Command Line Tools / Apple SDK on macOS).
\end{itemize}

Run the built-in health check right after installation:

\begin{Shaded}
\begin{Highlighting}[]
\NormalTok{iris {-}{-}version}
\NormalTok{iris doctor}
\end{Highlighting}
\end{Shaded}

\texttt{iris\ doctor} compiles and runs a test program natively and in
the interpreter, verifies LLVM-C, the linker, debugger adapters, and
optional ROS 2 / toolchain discovery, and prints exact remediation
commands if any dependency is missing. For automated toolchain
provisioning on supported environments, run \texttt{iris\ setup}.

\begin{quote}
\textbf{Environment Note on \texttt{-\/-emit\ eval} vs
\texttt{iris\ run}}: \texttt{iris\ run\ file.iris} compiles to a native
binary and executes it. \texttt{iris\ file.iris\ -\/-emit\ eval} runs
natively when available unless \texttt{IRIS\_FORCE\_INTERP=1} is set (or
\texttt{-\/-emit\ eval} triggers explicit fallback under
\texttt{IRIS\_INTERP\_FALLBACK=1}).
\end{quote}

\begin{center}\rule{0.5\linewidth}{0.5pt}\end{center}

\section{1.3 Hello, World}\label{hello-world}

Create a file named \texttt{hello.iris}:

\begin{Shaded}
\begin{Highlighting}[]
\CommentTok{/// Entry point returning 0 on success.}
\KeywordTok{def}\NormalTok{ main() }\OperatorTok{{-}\textgreater{}} \DataTypeTok{i64}\NormalTok{ \{}
    \KeywordTok{val}\NormalTok{ language }\OperatorTok{=} \StringTok{"IRIS"}
    \KeywordTok{val}\NormalTok{ version }\OperatorTok{=} \DecValTok{1}
\NormalTok{    println(}\StringTok{f"Hello from }\ControlFlowTok{\{}\NormalTok{language}\ControlFlowTok{\}}\StringTok{ v}\ControlFlowTok{\{}\NormalTok{version}\ControlFlowTok{\}}\StringTok{!"}\NormalTok{)}
    \FunctionTok{assert}\NormalTok{(version }\OperatorTok{==} \DecValTok{1}\NormalTok{)}
    \ControlFlowTok{return} \DecValTok{0}
\NormalTok{\}}
\end{Highlighting}
\end{Shaded}

Compile and run it natively:

\begin{Shaded}
\begin{Highlighting}[]
\NormalTok{iris run hello.iris}
\end{Highlighting}
\end{Shaded}

Output:

\begin{Shaded}
\begin{Highlighting}[]
\NormalTok{Hello from IRIS v1!}
\end{Highlighting}
\end{Shaded}

Or compile a standalone executable explicitly:

\begin{Shaded}
\begin{Highlighting}[]
\NormalTok{iris build hello.iris {-}o hello}
\NormalTok{./hello        \# or .\textbackslash{}hello.exe on Windows}
\end{Highlighting}
\end{Shaded}

Three conventions appear immediately in that six-line program:

\begin{enumerate}
\def\labelenumi{\arabic{enumi}.}
\tightlist
\item
  Programs start at \texttt{def\ main()\ -\textgreater{}\ i64} and
  return an exit status (\texttt{0} indicates success).
\item
  Immutable local bindings use \texttt{val}; the compiler infers
  \texttt{str} and \texttt{i64}.
\item
  Public and top-level functions end control paths with explicit
  \texttt{return}. Semicolons terminate statements when multiple
  statements share a block, though newline-separated single statements
  in clean blocks also parse cleanly across modern IRIS code. Inside
  book examples we consistently use clear \texttt{val}/\texttt{var}
  declarations and explicit \texttt{return} statements.
\end{enumerate}

\begin{center}\rule{0.5\linewidth}{0.5pt}\end{center}

\section{1.4 Project Templates}\label{project-templates}

For anything beyond a single script, scaffold a project using
\texttt{iris\ new\ {[}TEMPLATE{]}\ NAME}:

\begin{Shaded}
\begin{Highlighting}[]
\NormalTok{iris new app my\_service}
\NormalTok{cd my\_service}
\NormalTok{iris run main.iris}
\NormalTok{iris test}
\end{Highlighting}
\end{Shaded}

IRIS ships six built-in templates:

{\def\LTcaptype{none} % do not increment counter
\begin{longtable}[]{@{}
  >{\raggedright\arraybackslash}p{(\linewidth - 4\tabcolsep) * \real{0.3333}}
  >{\raggedright\arraybackslash}p{(\linewidth - 4\tabcolsep) * \real{0.3333}}
  >{\raggedright\arraybackslash}p{(\linewidth - 4\tabcolsep) * \real{0.3333}}@{}}
\toprule\noalign{}
\begin{minipage}[b]{\linewidth}\raggedright
Template
\end{minipage} & \begin{minipage}[b]{\linewidth}\raggedright
Command
\end{minipage} & \begin{minipage}[b]{\linewidth}\raggedright
Purpose
\end{minipage} \\
\midrule\noalign{}
\endhead
\bottomrule\noalign{}
\endlastfoot
\texttt{app} & \texttt{iris\ new\ app\ hello} & General executable
project with \texttt{iris.toml}, \texttt{main.iris}, and
\texttt{test\_*} assertions \\
\texttt{cli} & \texttt{iris\ new\ cli\ mytool} & Command-line tool
reading \texttt{args()} and formatted stdout/stderr \\
\texttt{server} & \texttt{iris\ new\ server\ api} & HTTP/TCP service
skeleton built on \texttt{std.http\_server} and \texttt{std.net} \\
\texttt{ros2} & \texttt{iris\ new\ ros2\ rover} & ROS 2 node project
importing \texttt{std.ros2} \\
\texttt{controller} & \texttt{iris\ new\ controller\ pid} &
Allocation-free control loop structured for embedded / evolution
targets \\
\texttt{ui} & \texttt{iris\ new\ ui\ studio} & Native desktop
application with \texttt{{[}app{]}} metadata, \texttt{std.ui} loop, and
bundling setup \\
\end{longtable}
}

Every template includes a working \texttt{iris.toml} manifest and
passing starter tests.

\begin{center}\rule{0.5\linewidth}{0.5pt}\end{center}

\section{1.5 Interactive REPL, Watch Mode, Testing, and IDE
Setup}\label{interactive-repl-watch-mode-testing-and-ide-setup}

A typical IRIS inner development loop relies on four fast commands:

\begin{Shaded}
\begin{Highlighting}[]
\NormalTok{iris repl                   \# Interactive expression evaluator (supports \textasciigrave{}bring std.math\textasciigrave{}, \textasciigrave{}:help\textasciigrave{}, \textasciigrave{}:quit\textasciigrave{})}
\NormalTok{iris check                  \# Fast static check across all .iris files in the workspace (types, borrows, effects)}
\NormalTok{iris watch main.iris        \# Re{-}run on every save (use \textasciigrave{}{-}{-}test\textasciigrave{} or \textasciigrave{}{-}{-}check\textasciigrave{} flags)}
\NormalTok{iris test                   \# Run every zero{-}argument \textasciigrave{}test\_*\textasciigrave{} function in the workspace}
\end{Highlighting}
\end{Shaded}

\subsection{VS Code and Editor
Support}\label{vs-code-and-editor-support}

The official VS Code extension (\texttt{vscode-iris/}) communicates with
the built-in compiler servers:

\begin{itemize}
\tightlist
\item
  \textbf{Language Server (\texttt{iris\ lsp})}: real-time diagnostics
  (powered by incremental \texttt{QueryEngine} caching), hover
  types/docstrings/effects, rich completions (including all embedded
  \texttt{std.*} modules), go-to-definition, find-references, rename,
  semantic tokens, folding ranges, inlay hints, code actions, and
  document formatting via \texttt{iris\ fmt}.
\item
  \textbf{Debug Adapter (\texttt{iris\ dap})}: supports both trace-based
  interpreter debugging (including time-travel \texttt{stepBack} and
  \texttt{reverseContinue}) and native DWARF debugging
  (\texttt{gdb\ -i\ dap} / \texttt{lldb-dap}).
\end{itemize}

Set \texttt{iris.executablePath} in VS Code to a stable installed copy
of \texttt{iris} (avoid pointing an active LSP server directly at a
Cargo build directory you intend to overwrite during
\texttt{cargo\ build}).

\subsection{Common Mistakes}\label{common-mistakes}

\begin{itemize}
\tightlist
\item
  \textbf{Omitting \texttt{-\textgreater{}\ i64} or \texttt{return\ 0}
  on \texttt{main()}}: IRIS entry points must declare
  \texttt{-\textgreater{}\ i64} and explicitly return an integer exit
  code.
\item
  \textbf{Using \texttt{print} vs \texttt{println}}: \texttt{print(msg)}
  writes without a trailing newline; \texttt{println(msg)} writes a
  complete line. Passing non-strings without \texttt{to\_str(x)} or an
  f-string fails type checking (\texttt{E0110}).
\item
  \textbf{Passing multiple files to \texttt{iris\ fmt}}:
  \texttt{iris\ fmt} formats one specified file at a time, or all
  \texttt{*.iris} files in the current directory when given no filename.
\end{itemize}

\subsection{Try It Yourself}\label{try-it-yourself}

\begin{enumerate}
\def\labelenumi{\arabic{enumi}.}
\tightlist
\item
  Run \texttt{iris\ doctor} in your terminal and inspect which native
  compilation and debugger backends are active on your machine.
\item
  Create a new project with \texttt{iris\ new\ cli\ greeter}, edit
  \texttt{main.iris} to read command-line arguments with
  \texttt{args()}, and test it via
  \texttt{iris\ run\ main.iris\ -\/-\ Alice}.
\item
  Launch \texttt{iris\ watch\ main.iris\ -\/-test} and verify that
  saving a broken \texttt{assert(...)} immediately reports the failing
  test line.
\end{enumerate}

\begin{center}\rule{0.5\linewidth}{0.5pt}\end{center}

\chapter{Chapter 2: Values, Bindings, and
Types}\label{chapter-2-values-bindings-and-types}

\section{2.1 Scalar Types}\label{scalar-types}

Every value in IRIS has a statically known type. Scalar primitives cover
signed integers, unsigned integers, IEEE 754 floating-point numbers,
booleans, and immutable UTF-8 strings:

{\def\LTcaptype{none} % do not increment counter
\begin{longtable}[]{@{}
  >{\raggedright\arraybackslash}p{(\linewidth - 6\tabcolsep) * \real{0.2500}}
  >{\raggedright\arraybackslash}p{(\linewidth - 6\tabcolsep) * \real{0.2500}}
  >{\raggedright\arraybackslash}p{(\linewidth - 6\tabcolsep) * \real{0.2500}}
  >{\raggedright\arraybackslash}p{(\linewidth - 6\tabcolsep) * \real{0.2500}}@{}}
\toprule\noalign{}
\begin{minipage}[b]{\linewidth}\raggedright
Type
\end{minipage} & \begin{minipage}[b]{\linewidth}\raggedright
Bit Width
\end{minipage} & \begin{minipage}[b]{\linewidth}\raggedright
Description
\end{minipage} & \begin{minipage}[b]{\linewidth}\raggedright
Example Literal / Cast
\end{minipage} \\
\midrule\noalign{}
\endhead
\bottomrule\noalign{}
\endlastfoot
\texttt{i64} & 64-bit signed & Default integer type & \texttt{42},
\texttt{-17}, \texttt{1000000} \\
\texttt{i32} & 32-bit signed & Compact/FFI signed integer &
\texttt{42\ to\ i32} \\
\texttt{i8} & 8-bit signed & Byte-sized signed integer &
\texttt{7\ to\ i8} \\
\texttt{u8} & 8-bit unsigned & Raw byte / pixel channel &
\texttt{255\ to\ u8} \\
\texttt{u32} & 32-bit unsigned & Packed color / flags / FFI &
\texttt{1024\ to\ u32} \\
\texttt{u64} & 64-bit unsigned & Large unsigned counts / hashes &
\texttt{99\ to\ u64} \\
\texttt{usize} & Pointer width & Unsigned indexing / const sizes &
\texttt{16\ to\ usize} \\
\texttt{f64} & 64-bit float & Default floating-point type &
\texttt{3.14159}, \texttt{1.0e-6} \\
\texttt{f32} & 32-bit float & Neural weights / compact buffers &
\texttt{0.5\ to\ f32} \\
\texttt{bool} & 1-bit logical & Boolean truth value & \texttt{true},
\texttt{false} \\
\texttt{str} & UTF-8 sequence & Immutable heap-allocated string &
\texttt{"IRIS"} \\
\end{longtable}
}

Unannotated integer literals default to \texttt{i64}, and unannotated
decimal literals default to \texttt{f64}. IRIS performs \textbf{no
implicit numeric conversions}: adding an \texttt{i64} and an
\texttt{f64} is a compile-time type error (\texttt{E0101}). Always make
precision and signedness boundaries explicit.

\begin{center}\rule{0.5\linewidth}{0.5pt}\end{center}

\section{2.2 Bindings, Constants, and Type
Inference}\label{bindings-constants-and-type-inference}

IRIS separates immutable value bindings (\texttt{val}), mutable
variables (\texttt{var}), and module-level constants (\texttt{const}):

\begin{Shaded}
\begin{Highlighting}[]
\KeywordTok{pub} \KeywordTok{const}\NormalTok{ MAX\_RETRIES: }\DataTypeTok{i64} \OperatorTok{=} \DecValTok{5}
\KeywordTok{pub} \KeywordTok{const}\NormalTok{ PI\_APPROX: }\DataTypeTok{f64} \OperatorTok{=} \FloatTok{3.141592653589793}

\KeywordTok{def}\NormalTok{ main() }\OperatorTok{{-}\textgreater{}} \DataTypeTok{i64}\NormalTok{ \{}
    \KeywordTok{val}\NormalTok{ threshold }\OperatorTok{=} \DecValTok{42}          \CommentTok{// immutable, inferred i64}
    \KeywordTok{var}\NormalTok{ attempts: }\DataTypeTok{i64} \OperatorTok{=} \DecValTok{0}       \CommentTok{// mutable, explicit i64}

    \ControlFlowTok{while}\NormalTok{ attempts }\OperatorTok{\textless{}}\NormalTok{ MAX\_RETRIES \{}
\NormalTok{        attempts }\OperatorTok{=}\NormalTok{ attempts }\OperatorTok{+} \DecValTok{1}
\NormalTok{    \}}

    \KeywordTok{val}\NormalTok{ ok\_flag }\OperatorTok{=}\NormalTok{ attempts }\OperatorTok{==}\NormalTok{ MAX\_RETRIES }\OperatorTok{\&\&}\NormalTok{ threshold }\OperatorTok{\textgreater{}} \DecValTok{0}
    \FunctionTok{assert}\NormalTok{(ok\_flag)}
    \FunctionTok{assert}\NormalTok{(PI\_APPROX }\OperatorTok{\textgreater{}} \FloatTok{3.14}\NormalTok{)}
\NormalTok{    println(}\StringTok{"bindings and constants: ok"}\NormalTok{)}
    \ControlFlowTok{return} \DecValTok{0}
\NormalTok{\}}
\end{Highlighting}
\end{Shaded}

Use \textbf{\texttt{val} by default} and reserve \texttt{var} for
accumulators, loop counters, and genuinely re-assigned local state.
Attempting to reassign a \texttt{val} raises compile error
\texttt{E0112} (\emph{Assignment to an Immutable Binding}).

\begin{quote}
\textbf{When Type Annotations Are Required}: Type inference works
locally from initializer expressions, literal types, and function
parameter signatures. However, two situations always require explicit
annotations:

\begin{enumerate}
\def\labelenumi{\arabic{enumi}.}
\tightlist
\item
  \textbf{All function parameters and return types} (IRIS requires
  explicit function boundaries).
\item
  \textbf{Generic constructors whose type parameter appears only in the
  return type}, such as empty collections:
  \texttt{val\ scores:\ map\textless{}str,\ i64\textgreater{}\ =\ map()},
  \texttt{val\ items:\ list\textless{}i64\textgreater{}\ =\ list()},
  \texttt{val\ s:\ Set\textless{}str\textgreater{}\ =\ set\_new()}.
\end{enumerate}
\end{quote}

\begin{center}\rule{0.5\linewidth}{0.5pt}\end{center}

\section{2.3 Literals, Escapes, F-Strings, and Explicit
Casts}\label{literals-escapes-f-strings-and-explicit-casts}

\subsection{Numeric and Scientific
Notation}\label{numeric-and-scientific-notation}

Floating-point literals in IRIS \textbf{must contain a decimal
point}---including scientific notation! Write \texttt{1.0e-6},
\texttt{2.5E3}, or \texttt{0.000001}, rather than \texttt{1e-6} (which
parses as an invalid token sequence). Underscore digit separators are
not part of scalar literal syntax; write plain digits such as
\texttt{1000000}.

\subsection{String Escapes and
F-Strings}\label{string-escapes-and-f-strings}

Double-quoted string literals support five escape sequences:

{\def\LTcaptype{none} % do not increment counter
\begin{longtable}[]{@{}ll@{}}
\toprule\noalign{}
Escape & Meaning \\
\midrule\noalign{}
\endhead
\bottomrule\noalign{}
\endlastfoot
\texttt{\textbackslash{}n} & Newline (LF) \\
\texttt{\textbackslash{}t} & Horizontal tab \\
\texttt{\textbackslash{}r} & Carriage return \\
\texttt{\textbackslash{}\textbackslash{}} & Literal backslash \\
\texttt{\textbackslash{}"} & Literal double quote \\
\end{longtable}
}

No other escape codes (such as
\texttt{\textbackslash{}\textquotesingle{}} or
\texttt{\textbackslash{}0}) are accepted inside string literals; an
unrecognized escape triggers compile error \texttt{E0003}.

String interpolation uses \texttt{f"..."} strings. At compile time,
\texttt{\{ident\}} lowers to \texttt{to\_str(ident)} concatenated
cleanly with surrounding segments:

\begin{Shaded}
\begin{Highlighting}[]
\KeywordTok{def}\NormalTok{ main() }\OperatorTok{{-}\textgreater{}} \DataTypeTok{i64}\NormalTok{ \{}
    \KeywordTok{val}\NormalTok{ label }\OperatorTok{=} \StringTok{"orbit"}
    \KeywordTok{val}\NormalTok{ step }\OperatorTok{=} \DecValTok{12}
    \KeywordTok{val}\NormalTok{ msg }\OperatorTok{=} \StringTok{f"[}\ControlFlowTok{\{}\NormalTok{label}\ControlFlowTok{\}}\StringTok{] step=}\ControlFlowTok{\{}\NormalTok{step}\ControlFlowTok{\}}\StringTok{\textbackslash{}nnext=\textbackslash{}"}\NormalTok{ready\textbackslash{}}\StringTok{""}
    \FunctionTok{assert}\NormalTok{(}\FunctionTok{contains}\NormalTok{(msg, }\StringTok{"[orbit] step=12"}\NormalTok{))}
\NormalTok{    println(msg)}
    \ControlFlowTok{return} \DecValTok{0}
\NormalTok{\}}
\end{Highlighting}
\end{Shaded}

\begin{quote}
\textbf{F-String Rule}: \texttt{\{...\}} placeholders accept
\textbf{simple identifiers only}. Expressions with dots or calls such as
\texttt{f"\{point.x\}"} or \texttt{f"\{len(xs)\}"} are rejected by the
lexer/parser. Extract them into a local \texttt{val} first:

\begin{Shaded}
\begin{Highlighting}[]
\KeywordTok{val}\NormalTok{ px }\OperatorTok{=}\NormalTok{ point}\OperatorTok{.}\NormalTok{x}
\KeywordTok{val}\NormalTok{ text }\OperatorTok{=} \StringTok{f"x = }\ControlFlowTok{\{}\NormalTok{px}\ControlFlowTok{\}}\StringTok{"}
\end{Highlighting}
\end{Shaded}
\end{quote}

\subsection{\texorpdfstring{Explicit Type Conversions (\texttt{to},
\texttt{to\_i64}, \texttt{to\_f64},
\texttt{to\_str})}{Explicit Type Conversions (to, to\_i64, to\_f64, to\_str)}}\label{explicit-type-conversions-to-to_i64-to_f64-to_str}

Convert between numeric widths with \texttt{expr\ to\ target\_type}, or
via the core conversion builtins \texttt{to\_i64(x)},
\texttt{to\_f64(x)}, and \texttt{to\_str(x)}:

\begin{Shaded}
\begin{Highlighting}[]
\KeywordTok{def}\NormalTok{ main() }\OperatorTok{{-}\textgreater{}} \DataTypeTok{i64}\NormalTok{ \{}
    \KeywordTok{val}\NormalTok{ count: }\DataTypeTok{i64} \OperatorTok{=} \DecValTok{7}
    \KeywordTok{val}\NormalTok{ total: }\DataTypeTok{f64} \OperatorTok{=} \FloatTok{100.0}
    \KeywordTok{val}\NormalTok{ ratio }\OperatorTok{=}\NormalTok{ total }\OperatorTok{/}\NormalTok{ to\_f64(count)}
    \KeywordTok{val}\NormalTok{ byte\_val }\OperatorTok{=} \DecValTok{200} \KeywordTok{to} \DataTypeTok{u8}
    \KeywordTok{val}\NormalTok{ back\_i64 }\OperatorTok{=}\NormalTok{ byte\_val }\KeywordTok{to} \DataTypeTok{i64}
    \KeywordTok{val}\NormalTok{ half\_i32 }\OperatorTok{=}\NormalTok{ (count }\KeywordTok{to} \DataTypeTok{i32}\NormalTok{) }\OperatorTok{+}\NormalTok{ (}\DecValTok{3} \KeywordTok{to} \DataTypeTok{i32}\NormalTok{)}

    \FunctionTok{assert}\NormalTok{(ratio }\OperatorTok{\textgreater{}} \FloatTok{14.28} \OperatorTok{\&\&}\NormalTok{ ratio }\OperatorTok{\textless{}} \FloatTok{14.29}\NormalTok{)}
    \FunctionTok{assert}\NormalTok{(back\_i64 }\OperatorTok{==} \DecValTok{200}\NormalTok{)}
    \FunctionTok{assert}\NormalTok{((half\_i32 }\KeywordTok{to} \DataTypeTok{i64}\NormalTok{) }\OperatorTok{==} \DecValTok{10}\NormalTok{)}
\NormalTok{    println(}\StringTok{"casts: ok"}\NormalTok{)}
    \ControlFlowTok{return} \DecValTok{0}
\NormalTok{\}}
\end{Highlighting}
\end{Shaded}

\begin{center}\rule{0.5\linewidth}{0.5pt}\end{center}

\section{2.4 Checked Integer Arithmetic and Wrapping
Operations}\label{checked-integer-arithmetic-and-wrapping-operations}

In IRIS, \texttt{+}, \texttt{-}, \texttt{*}, and \texttt{/} on
\texttt{i64} and \texttt{i32} are \textbf{checked at runtime} on both
native LLVM builds and the interpreter. If an operation overflows or
underflows its signed range (including \texttt{i64::MIN\ /\ -1}),
execution halts cleanly with diagnostic \texttt{E0409} naming the exact
operation and operands instead of silently wrapping around. Division or
remainder by zero halts with \texttt{E0402}.

Constant folding obeys the exact same rule: expressions that would
overflow are \textbf{not} silently folded to wrapped values.

When wrap-around two's-complement arithmetic is genuinely desired---such
as hash mixing, linear congruential PRNGs, or cyclic counters---use the
explicit wrapping intrinsics \texttt{wrapping\_add(a,\ b)},
\texttt{wrapping\_sub(a,\ b)}, and \texttt{wrapping\_mul(a,\ b)}:

\begin{Shaded}
\begin{Highlighting}[]
\KeywordTok{def}\NormalTok{ hash\_mix(state: }\DataTypeTok{i64}\NormalTok{, value: }\DataTypeTok{i64}\NormalTok{) }\OperatorTok{{-}\textgreater{}} \DataTypeTok{i64}\NormalTok{ \{}
    \KeywordTok{val}\NormalTok{ stepped }\OperatorTok{=}\NormalTok{ wrapping\_mul(state, }\DecValTok{6364136223846793005}\NormalTok{)}
    \ControlFlowTok{return}\NormalTok{ wrapping\_add(stepped, value)}
\NormalTok{\}}

\KeywordTok{def}\NormalTok{ main() }\OperatorTok{{-}\textgreater{}} \DataTypeTok{i64}\NormalTok{ \{}
    \KeywordTok{val}\NormalTok{ near\_max }\OperatorTok{=} \DecValTok{9223372036854775800}
    \CommentTok{// Regular addition would trap with E0409 if it exceeded i64 max;}
    \CommentTok{// wrapping\_add wraps cleanly in two\textquotesingle{}s complement:}
    \KeywordTok{val}\NormalTok{ wrapped }\OperatorTok{=}\NormalTok{ wrapping\_add(near\_max, }\DecValTok{100}\NormalTok{)}
    \KeywordTok{val}\NormalTok{ mixed }\OperatorTok{=}\NormalTok{ hash\_mix(}\DecValTok{17}\NormalTok{, }\DecValTok{42}\NormalTok{)}
    \FunctionTok{assert}\NormalTok{(wrapped }\OperatorTok{\textless{}} \DecValTok{0}\NormalTok{)}
    \FunctionTok{assert}\NormalTok{(mixed }\OperatorTok{!=} \DecValTok{0}\NormalTok{)}
\NormalTok{    println(}\StringTok{"wrapping arithmetic: ok"}\NormalTok{)}
    \ControlFlowTok{return} \DecValTok{0}
\NormalTok{\}}
\end{Highlighting}
\end{Shaded}

\begin{center}\rule{0.5\linewidth}{0.5pt}\end{center}

\section{2.5 Type Aliases}\label{type-aliases}

Use \texttt{type\ Name\ =\ ExistingType} to give domain meaning to
function pointers, matrices, keys, or compound tuples without runtime
overhead:

\begin{Shaded}
\begin{Highlighting}[]
\KeywordTok{pub} \KeywordTok{type}\NormalTok{ NodeId }\OperatorTok{=} \DataTypeTok{i64}
\KeywordTok{pub} \KeywordTok{type}\NormalTok{ Weight }\OperatorTok{=} \DataTypeTok{f64}
\KeywordTok{pub} \KeywordTok{type}\NormalTok{ Predicate }\OperatorTok{=}\NormalTok{ (}\DataTypeTok{i64}\NormalTok{) }\OperatorTok{{-}\textgreater{}} \DataTypeTok{bool}

\KeywordTok{def}\NormalTok{ count\_matching(items: }\DataTypeTok{list}\OperatorTok{\textless{}}\NormalTok{NodeId}\OperatorTok{\textgreater{}}\NormalTok{, pred: Predicate) }\OperatorTok{{-}\textgreater{}} \DataTypeTok{i64}\NormalTok{ \{}
    \KeywordTok{var}\NormalTok{ hits }\OperatorTok{=} \DecValTok{0}
    \ControlFlowTok{for}\NormalTok{ id }\ControlFlowTok{in}\NormalTok{ items \{}
        \ControlFlowTok{if}\NormalTok{ pred(id) \{}
\NormalTok{            hits }\OperatorTok{=}\NormalTok{ hits }\OperatorTok{+} \DecValTok{1}
\NormalTok{        \}}
\NormalTok{    \}}
    \ControlFlowTok{return}\NormalTok{ hits}
\NormalTok{\}}

\KeywordTok{def}\NormalTok{ main() }\OperatorTok{{-}\textgreater{}} \DataTypeTok{i64}\NormalTok{ \{}
    \KeywordTok{val}\NormalTok{ ids: }\DataTypeTok{list}\OperatorTok{\textless{}}\NormalTok{NodeId}\OperatorTok{\textgreater{}} \OperatorTok{=} \DataTypeTok{list}\NormalTok{(}\DecValTok{1}\NormalTok{, }\DecValTok{4}\NormalTok{, }\DecValTok{7}\NormalTok{, }\DecValTok{10}\NormalTok{, }\DecValTok{13}\NormalTok{)}
    \KeywordTok{val}\NormalTok{ is\_even: Predicate }\OperatorTok{=} \OperatorTok{|}\NormalTok{n: NodeId}\OperatorTok{|}\NormalTok{ n }\OperatorTok{\%} \DecValTok{2} \OperatorTok{==} \DecValTok{0}
    \KeywordTok{val}\NormalTok{ threshold: Weight }\OperatorTok{=} \FloatTok{2.0}
    \FunctionTok{assert}\NormalTok{(count\_matching(ids, is\_even) }\OperatorTok{==} \DecValTok{2}\NormalTok{)}
    \FunctionTok{assert}\NormalTok{(threshold }\OperatorTok{\textgreater{}} \FloatTok{1.0}\NormalTok{)}
\NormalTok{    println(}\StringTok{"type aliases: ok"}\NormalTok{)}
    \ControlFlowTok{return} \DecValTok{0}
\NormalTok{\}}
\end{Highlighting}
\end{Shaded}

\subsection{Common Mistakes}\label{common-mistakes-1}

\begin{itemize}
\tightlist
\item
  \textbf{Writing \texttt{1e-6} without a decimal point}: write
  \texttt{1.0e-6}.
\item
  \textbf{Interpolating expressions inside f-strings}: bind
  \texttt{val\ field\ =\ obj.field} before \texttt{f"\{field\}"}.
\item
  \textbf{Mixing \texttt{i64} and \texttt{f64} in arithmetic or
  comparisons}: convert one side explicitly with \texttt{to\_f64(n)} or
  \texttt{to\_i64(x)}.
\item
  \textbf{Using \texttt{\textbackslash{}\textquotesingle{}} in string
  literals}: single quotes need no escape inside \texttt{"..."};
  unsupported escapes fail with \texttt{E0003}.
\end{itemize}

\subsection{Try It Yourself}\label{try-it-yourself-1}

\begin{enumerate}
\def\labelenumi{\arabic{enumi}.}
\tightlist
\item
  Write a function
  \texttt{celsius\_to\_fahrenheit(c:\ f64)\ -\textgreater{}\ f64} and
  verify \texttt{0.0\ -\textgreater{}\ 32.0} and
  \texttt{100.0\ -\textgreater{}\ 212.0} with \texttt{assert}.
\item
  Write a tiny PRNG function
  \texttt{next\_rand(seed:\ i64)\ -\textgreater{}\ i64} using
  \texttt{wrapping\_mul} and \texttt{wrapping\_add} that executes safely
  for 1,000 iterations without overflowing.
\item
  Experiment with forcing an intentional overflow in a small scratch
  file
  (\texttt{val\ x\ =\ 9223372036854775807;\ println(to\_str(x\ +\ 1))})
  and inspect the compiler's \texttt{E0409} message and
  \texttt{iris\ explain\ E0409}.
\end{enumerate}

\begin{center}\rule{0.5\linewidth}{0.5pt}\end{center}

\chapter{Chapter 3: Functions and
Modules}\label{chapter-3-functions-and-modules}

\section{3.1 Defining Functions and Explicit
Returns}\label{defining-functions-and-explicit-returns}

Functions are declared with \texttt{def} (or \texttt{pub\ def} to export
from a module). Every parameter declares \texttt{name:\ type}, followed
by \texttt{-\textgreater{}\ ReturnType} before the function body:

\begin{Shaded}
\begin{Highlighting}[]
\KeywordTok{def}\NormalTok{ clamp\_score(raw: }\DataTypeTok{i64}\NormalTok{, low: }\DataTypeTok{i64}\NormalTok{, high: }\DataTypeTok{i64}\NormalTok{) }\OperatorTok{{-}\textgreater{}} \DataTypeTok{i64}\NormalTok{ \{}
    \ControlFlowTok{if}\NormalTok{ raw }\OperatorTok{\textless{}}\NormalTok{ low \{}
        \ControlFlowTok{return}\NormalTok{ low}
\NormalTok{    \}}
    \ControlFlowTok{if}\NormalTok{ raw }\OperatorTok{\textgreater{}}\NormalTok{ high \{}
        \ControlFlowTok{return}\NormalTok{ high}
\NormalTok{    \}}
    \ControlFlowTok{return}\NormalTok{ raw}
\NormalTok{\}}

\KeywordTok{def}\NormalTok{ main() }\OperatorTok{{-}\textgreater{}} \DataTypeTok{i64}\NormalTok{ \{}
    \FunctionTok{assert}\NormalTok{(clamp\_score(}\OperatorTok{{-}}\DecValTok{5}\NormalTok{, }\DecValTok{0}\NormalTok{, }\DecValTok{100}\NormalTok{) }\OperatorTok{==} \DecValTok{0}\NormalTok{)}
    \FunctionTok{assert}\NormalTok{(clamp\_score(}\DecValTok{150}\NormalTok{, }\DecValTok{0}\NormalTok{, }\DecValTok{100}\NormalTok{) }\OperatorTok{==} \DecValTok{100}\NormalTok{)}
    \FunctionTok{assert}\NormalTok{(clamp\_score(}\DecValTok{74}\NormalTok{, }\DecValTok{0}\NormalTok{, }\DecValTok{100}\NormalTok{) }\OperatorTok{==} \DecValTok{74}\NormalTok{)}
\NormalTok{    println(}\StringTok{"functions: ok"}\NormalTok{)}
    \ControlFlowTok{return} \DecValTok{0}
\NormalTok{\}}
\end{Highlighting}
\end{Shaded}

Named functions terminate control-flow paths with explicit
\texttt{return} statements (\texttt{return} without a value exits
\texttt{unit}-returning procedures). Inside expression blocks
(\texttt{if}/\texttt{else}, \texttt{when}, and lambdas), tail
expressions without \texttt{return} remain valid and produce the block's
value.

\begin{center}\rule{0.5\linewidth}{0.5pt}\end{center}

\section{3.2 Default Parameters, Named Arguments, and Extension
Calls}\label{default-parameters-named-arguments-and-extension-calls}

IRIS supports \textbf{default parameter values} and \textbf{named
arguments} on user-defined functions, making multi-option APIs readable
without builder boilerplate. In addition, any function whose first
parameter has type \texttt{T} can be called with \textbf{UFCS dot
notation} (\texttt{x.func(...)}):

\begin{Shaded}
\begin{Highlighting}[]
\KeywordTok{def}\NormalTok{ scale\_and\_bias(value: }\DataTypeTok{f64}\NormalTok{, gain: }\DataTypeTok{f64} \OperatorTok{=} \FloatTok{1.0}\NormalTok{, bias: }\DataTypeTok{f64} \OperatorTok{=} \FloatTok{0.0}\NormalTok{) }\OperatorTok{{-}\textgreater{}} \DataTypeTok{f64}\NormalTok{ \{}
    \ControlFlowTok{return}\NormalTok{ value }\OperatorTok{*}\NormalTok{ gain }\OperatorTok{+}\NormalTok{ bias}
\NormalTok{\}}

\KeywordTok{def}\NormalTok{ difference(left: }\DataTypeTok{i64}\NormalTok{, right: }\DataTypeTok{i64}\NormalTok{) }\OperatorTok{{-}\textgreater{}} \DataTypeTok{i64}\NormalTok{ \{}
    \ControlFlowTok{return}\NormalTok{ left }\OperatorTok{{-}}\NormalTok{ right}
\NormalTok{\}}

\KeywordTok{def}\NormalTok{ doubled(n: }\DataTypeTok{i64}\NormalTok{) }\OperatorTok{{-}\textgreater{}} \DataTypeTok{i64}\NormalTok{ \{}
    \ControlFlowTok{return}\NormalTok{ n }\OperatorTok{*} \DecValTok{2}
\NormalTok{\}}

\KeywordTok{def}\NormalTok{ main() }\OperatorTok{{-}\textgreater{}} \DataTypeTok{i64}\NormalTok{ \{}
    \FunctionTok{assert}\NormalTok{(scale\_and\_bias(}\FloatTok{5.0}\NormalTok{) }\OperatorTok{==} \FloatTok{5.0}\NormalTok{)}
    \FunctionTok{assert}\NormalTok{(scale\_and\_bias(}\FloatTok{5.0}\NormalTok{, gain }\OperatorTok{=} \FloatTok{2.0}\NormalTok{, bias }\OperatorTok{=} \FloatTok{3.0}\NormalTok{) }\OperatorTok{==} \FloatTok{13.0}\NormalTok{)}
    \FunctionTok{assert}\NormalTok{(difference(right }\OperatorTok{=} \DecValTok{3}\NormalTok{, left }\OperatorTok{=} \DecValTok{10}\NormalTok{) }\OperatorTok{==} \DecValTok{7}\NormalTok{)}
    \FunctionTok{assert}\NormalTok{(}\DecValTok{21}\OperatorTok{.}\NormalTok{doubled() }\OperatorTok{==} \DecValTok{42}\NormalTok{)}
\NormalTok{    println(}\StringTok{"named args and ufcs: ok"}\NormalTok{)}
    \ControlFlowTok{return} \DecValTok{0}
\NormalTok{\}}
\end{Highlighting}
\end{Shaded}

Named arguments are reordered cleanly at compile time to match the
callee's signature, and missing trailing/gapped arguments with defaults
are filled automatically. Builtins and \texttt{extern\ def} declarations
do not expose parameter names, so pass their arguments positionally.

\begin{center}\rule{0.5\linewidth}{0.5pt}\end{center}

\section{3.3 Recursion and Tail-Call
Optimization}\label{recursion-and-tail-call-optimization}

IRIS implements \textbf{IR-level tail-call optimization
(\texttt{TailCallPass})} inside the standard compiler pipeline,
guaranteeing constant stack space on both the native LLVM backend and
the SSA interpreter for self-recursive tail calls:

\begin{itemize}
\tightlist
\item
  \textbf{Self-recursion in tail expression form} lowers directly into a
  branch back to a fresh loop header inside SSA IR---before
  loop-invariant code motion (\texttt{LicmPass}), so both the native
  backend and the interpreter run deep self-recursion in constant stack
  space.
\item
  \textbf{Limitation}: \texttt{TailCallPass} rewrites
  \emph{self-recursion} whose result returns directly (such as
  \texttt{if\ cond\ \{\ base\ \}\ else\ \{\ self\_call(...)\ \}}).
  Mutually recursive cycles between distinct functions are left as
  normal calls.
\end{itemize}

\begin{Shaded}
\begin{Highlighting}[]
\KeywordTok{def}\NormalTok{ sum\_tail(n: }\DataTypeTok{i64}\NormalTok{, acc: }\DataTypeTok{i64}\NormalTok{) }\OperatorTok{{-}\textgreater{}} \DataTypeTok{i64}\NormalTok{ \{}
    \ControlFlowTok{if}\NormalTok{ n }\OperatorTok{\textless{}=} \DecValTok{0}\NormalTok{ \{}
\NormalTok{        acc}
\NormalTok{    \} }\ControlFlowTok{else}\NormalTok{ \{}
\NormalTok{        sum\_tail(n }\OperatorTok{{-}} \DecValTok{1}\NormalTok{, acc }\OperatorTok{+}\NormalTok{ n)}
\NormalTok{    \}}
\NormalTok{\}}

\KeywordTok{def}\NormalTok{ main() }\OperatorTok{{-}\textgreater{}} \DataTypeTok{i64}\NormalTok{ \{}
    \KeywordTok{val}\NormalTok{ total }\OperatorTok{=}\NormalTok{ sum\_tail(}\DecValTok{50000}\NormalTok{, }\DecValTok{0}\NormalTok{)}
    \FunctionTok{assert}\NormalTok{(total }\OperatorTok{==} \DecValTok{1250025000}\NormalTok{)}
\NormalTok{    println(}\StringTok{"tail call optimization: ok"}\NormalTok{)}
    \ControlFlowTok{return} \DecValTok{0}
\NormalTok{\}}
\end{Highlighting}
\end{Shaded}

Running 50,000 recursive calls in \texttt{sum\_tail} runs in constant
stack frames on both native binaries and the interpreter without
tripping the default call-depth guard (\texttt{-\/-max-depth\ 250}).

\begin{center}\rule{0.5\linewidth}{0.5pt}\end{center}

\section{3.4 First-Class Functions and
Closures}\label{first-class-functions-and-closures}

Functions can be passed by name or constructed inline as
\textbf{closures} with \texttt{\textbar{}params\textbar{}\ expr} or
\texttt{\textbar{}params\textbar{}\ \{\ ...\ \}}. Closures capture outer
immutable values cleanly:

\begin{Shaded}
\begin{Highlighting}[]
\KeywordTok{def}\NormalTok{ apply\_twice(value: }\DataTypeTok{i64}\NormalTok{, f: (}\DataTypeTok{i64}\NormalTok{) }\OperatorTok{{-}\textgreater{}} \DataTypeTok{i64}\NormalTok{) }\OperatorTok{{-}\textgreater{}} \DataTypeTok{i64}\NormalTok{ \{}
    \ControlFlowTok{return}\NormalTok{ f(f(value))}
\NormalTok{\}}

\KeywordTok{def}\NormalTok{ main() }\OperatorTok{{-}\textgreater{}} \DataTypeTok{i64}\NormalTok{ \{}
    \KeywordTok{val}\NormalTok{ step }\OperatorTok{=} \DecValTok{5}
    \KeywordTok{val}\NormalTok{ advance }\OperatorTok{=} \OperatorTok{|}\NormalTok{x: }\DataTypeTok{i64}\OperatorTok{|}\NormalTok{ x }\OperatorTok{+}\NormalTok{ step}
    \KeywordTok{val}\NormalTok{ square }\OperatorTok{=} \OperatorTok{|}\NormalTok{x: }\DataTypeTok{i64}\OperatorTok{|}\NormalTok{ x }\OperatorTok{*}\NormalTok{ x}

    \FunctionTok{assert}\NormalTok{(apply\_twice(}\DecValTok{10}\NormalTok{, advance) }\OperatorTok{==} \DecValTok{20}\NormalTok{)}
    \FunctionTok{assert}\NormalTok{(apply\_twice(}\DecValTok{3}\NormalTok{, square) }\OperatorTok{==} \DecValTok{81}\NormalTok{)}
\NormalTok{    println(}\StringTok{"closures: ok"}\NormalTok{)}
    \ControlFlowTok{return} \DecValTok{0}
\NormalTok{\}}
\end{Highlighting}
\end{Shaded}

For closures whose bodies allocate, perform I/O, or accept
effect-polymorphic callbacks under \texttt{-\/-strict-effects}, use
pipe-arrow syntax with an optional effect clause:
\texttt{\textbar{}x:\ i64\textbar{}\ -\textgreater{}\ i64\ effect\ alloc}.

\begin{center}\rule{0.5\linewidth}{0.5pt}\end{center}

\section{3.5 Modules, Visibility, and Transitive
Imports}\label{modules-visibility-and-transitive-imports}

IRIS organizes code across files and embedded standard libraries via
\texttt{bring}:

\begin{Shaded}
\begin{Highlighting}[]
\KeywordTok{bring}\NormalTok{ std}\OperatorTok{.}\NormalTok{math              }\CommentTok{// Embedded standard{-}library module}
\KeywordTok{bring}\NormalTok{ std}\OperatorTok{.}\NormalTok{json              }\CommentTok{// Another embedded module}
\KeywordTok{bring} \StringTok{"physics.iris"}        \CommentTok{// Local sibling module relative to this file}
\end{Highlighting}
\end{Shaded}

Inside a library module (for instance \texttt{math\_helpers.iris}),
\textbf{definitions are private by default}. Only declarations marked
\texttt{pub} are exported to importers:

\begin{itemize}
\tightlist
\item
  \texttt{pub\ def\ ...} exports a function.
\item
  \texttt{pub\ record\ ...} and \texttt{pub\ choice\ ...} export structs
  and algebraic enums.
\item
  \texttt{pub\ const\ ...} and \texttt{pub\ type\ ...} export constants
  and type aliases.
\item
  \texttt{pub\ bring\ "submod.iris"} (or
  \texttt{pub\ bring\ std.ui\_theme}) \textbf{re-exports} a brought
  module's public functions, records, choices, and type aliases to
  downstream files---just like \texttt{std.ui} re-exports
  \texttt{std.ui\_theme}, \texttt{std.ui\_protocol},
  \texttt{std.ui\_motion}, \texttt{std.ui\_icons},
  \texttt{std.ui\_layout}, and \texttt{std.ui\_widgets} so callers only
  need \texttt{bring\ std.ui}.
\end{itemize}

Attempting to call a private helper in another file fails at compile
time with \texttt{E0100}, naming the module and instructing you to add
\texttt{pub}.

\begin{quote}
\textbf{Cross-Module Trait Limit}: Co-locate
\texttt{impl\ Trait\ for\ Record} blocks in the same module as the
defining \texttt{record} (and mark helper methods carefully), since
traits themselves have module-global names while record visibility
follows \texttt{pub}.
\end{quote}

\begin{center}\rule{0.5\linewidth}{0.5pt}\end{center}

\section{3.6 Documentation Comments and HTML
Docs}\label{documentation-comments-and-html-docs}

Prefix public declarations with \texttt{///} doc comments:

\begin{Shaded}
\begin{Highlighting}[]
\CommentTok{/// Compute kinetic energy \textasciigrave{}0.5 * m * v * v\textasciigrave{} for a non{-}negative mass \textasciigrave{}m\textasciigrave{}.}
\KeywordTok{pub} \KeywordTok{def}\NormalTok{ kinetic\_energy(mass: }\DataTypeTok{f64}\NormalTok{, speed: }\DataTypeTok{f64}\NormalTok{) }\OperatorTok{{-}\textgreater{}} \DataTypeTok{f64}\NormalTok{ \{}
    \ControlFlowTok{return} \FloatTok{0.5} \OperatorTok{*}\NormalTok{ mass }\OperatorTok{*}\NormalTok{ speed }\OperatorTok{*}\NormalTok{ speed}
\NormalTok{\}}
\end{Highlighting}
\end{Shaded}

These doc comments surface automatically in two places:

\begin{enumerate}
\def\labelenumi{\arabic{enumi}.}
\tightlist
\item
  \textbf{In your editor}: \texttt{iris\ lsp} displays \texttt{///}
  summaries on hover and inside autocomplete popups.
\item
  \textbf{Generated API documentation}:
  \texttt{iris\ docs\ \textless{}path\textgreater{}} parses \texttt{///}
  comments across your project and renders navigable HTML documentation.
\end{enumerate}

\subsection{Common Mistakes}\label{common-mistakes-2}

\begin{itemize}
\tightlist
\item
  \textbf{Calling \texttt{pub} items without \texttt{pub} keyword}:
  forgetting \texttt{pub} on a helper record or function inside
  \texttt{"helpers.iris"} causes \texttt{E0100} in \texttt{main.iris}.
\item
  \textbf{Naming a helper after a builtin by accident}: avoid shadowing
  core names like \texttt{len}, \texttt{min}, \texttt{max},
  \texttt{split}, \texttt{trim}, \texttt{filter}, or \texttt{map} unless
  intentionally overriding in local scope; \texttt{std.*} modules
  already provide rich namespaced helpers.
\item
  \textbf{Putting non-tail work after a recursive call}:
  \texttt{return\ 1\ +\ sum(n\ -\ 1)} is \emph{not} a tail call because
  addition happens after the call returns. Pass an accumulator parameter
  instead.
\end{itemize}

\subsection{Try It Yourself}\label{try-it-yourself-2}

\begin{enumerate}
\def\labelenumi{\arabic{enumi}.}
\tightlist
\item
  Split a two-file mini-project (\texttt{geometry.iris} +
  \texttt{main.iris}) where \texttt{geometry.iris} exports
  \texttt{pub\ record\ Rect}, \texttt{pub\ def\ rect\_area}, and a
  private helper \texttt{def\ clamp\_non\_neg}. Verify that calling
  \texttt{clamp\_non\_neg} from \texttt{main.iris} fails to compile
  until marked \texttt{pub}.
\item
  Rewrite Fibonacci recursively with tail accumulators
  (\texttt{fib\_tail(n,\ a,\ b)}) and compute
  \texttt{fib\_tail(50,\ 0,\ 1)}.
\item
  Document three \texttt{pub\ def} functions with \texttt{///} comments
  and run \texttt{iris\ docs} to inspect the generated HTML output.
\end{enumerate}

\begin{center}\rule{0.5\linewidth}{0.5pt}\end{center}

\chapter{Chapter 4: Control Flow}\label{chapter-4-control-flow}

\section{4.1 Expression-Oriented
Conditionals}\label{expression-oriented-conditionals}

In IRIS, \texttt{if\ /\ else} is an \textbf{expression} that evaluates
to a value when both branches return the same type:

\begin{Shaded}
\begin{Highlighting}[]
\KeywordTok{def}\NormalTok{ classify\_temp(c: }\DataTypeTok{f64}\NormalTok{) }\OperatorTok{{-}\textgreater{}} \DataTypeTok{str}\NormalTok{ \{}
    \KeywordTok{val} \FunctionTok{band} \OperatorTok{=} \ControlFlowTok{if}\NormalTok{ c }\OperatorTok{\textless{}} \FloatTok{0.0}\NormalTok{ \{}
        \StringTok{"freezing"}
\NormalTok{    \} }\ControlFlowTok{else} \ControlFlowTok{if}\NormalTok{ c }\OperatorTok{\textless{}} \FloatTok{25.0}\NormalTok{ \{}
        \StringTok{"temperate"}
\NormalTok{    \} }\ControlFlowTok{else}\NormalTok{ \{}
        \StringTok{"hot"}
\NormalTok{    \}}
    \ControlFlowTok{return} \FunctionTok{band}
\NormalTok{\}}

\KeywordTok{def}\NormalTok{ main() }\OperatorTok{{-}\textgreater{}} \DataTypeTok{i64}\NormalTok{ \{}
    \FunctionTok{assert}\NormalTok{(classify\_temp(}\OperatorTok{{-}}\FloatTok{3.0}\NormalTok{) }\OperatorTok{==} \StringTok{"freezing"}\NormalTok{)}
    \FunctionTok{assert}\NormalTok{(classify\_temp(}\FloatTok{18.5}\NormalTok{) }\OperatorTok{==} \StringTok{"temperate"}\NormalTok{)}
    \FunctionTok{assert}\NormalTok{(classify\_temp(}\FloatTok{35.0}\NormalTok{) }\OperatorTok{==} \StringTok{"hot"}\NormalTok{)}
\NormalTok{    println(}\StringTok{"conditionals: ok"}\NormalTok{)}
    \ControlFlowTok{return} \DecValTok{0}
\NormalTok{\}}
\end{Highlighting}
\end{Shaded}

When using \texttt{if} solely for side effects inside a statement
(\texttt{if\ ready\ \{\ notify();\ \}}), an \texttt{else} branch is
optional. When using the result of \texttt{if} in a \texttt{val},
argument, or \texttt{return}, an \texttt{else} branch with a compatible
type is \textbf{mandatory} (omitting it raises \texttt{E0119}, or
mismatched branch types raise \texttt{E0208}).

\begin{center}\rule{0.5\linewidth}{0.5pt}\end{center}

\section{4.2 Loops: Ranges, Collections, Conditionals, and Infinite
Loops}\label{loops-ranges-collections-conditionals-and-infinite-loops}

IRIS offers four clean iteration forms:

\begin{enumerate}
\def\labelenumi{\arabic{enumi}.}
\tightlist
\item
  \textbf{Exclusive integer range}: \texttt{for\ i\ in\ start..end}
  iterates \texttt{start,\ start\ +\ 1,\ ...,\ end\ -\ 1}.
\item
  \textbf{Stepped range}: \texttt{for\ i\ in\ start..end\ by\ step}
  advances by \texttt{step} each iteration (\texttt{by} is a contextual
  keyword after a range).
\item
  \textbf{Collection iteration}: \texttt{for\ item\ in\ list\_val}
  iterates every element of a \texttt{list\textless{}T\textgreater{}},
  fixed array \texttt{{[}T;\ N{]}}, or map/collection view.
\item
  \textbf{Conditional \texttt{while} and unconditional \texttt{loop}}:
  repeat while a condition holds, or until \texttt{break}.
\end{enumerate}

\begin{Shaded}
\begin{Highlighting}[]
\KeywordTok{def}\NormalTok{ main() }\OperatorTok{{-}\textgreater{}} \DataTypeTok{i64}\NormalTok{ \{}
    \CommentTok{// 1. Exclusive range 1..11 (1 through 10)}
    \KeywordTok{var}\NormalTok{ sum }\OperatorTok{=} \DecValTok{0}
    \ControlFlowTok{for}\NormalTok{ n }\ControlFlowTok{in} \DecValTok{1}\OperatorTok{..}\DecValTok{11}\NormalTok{ \{}
\NormalTok{        sum }\OperatorTok{=}\NormalTok{ sum }\OperatorTok{+}\NormalTok{ n}
\NormalTok{    \}}

    \CommentTok{// 2. Stepped range with \textasciigrave{}by\textasciigrave{}}
    \KeywordTok{var}\NormalTok{ evens }\OperatorTok{=} \DecValTok{0}
    \ControlFlowTok{for}\NormalTok{ k }\ControlFlowTok{in} \DecValTok{0}\OperatorTok{..}\DecValTok{10}\NormalTok{ by }\DecValTok{2}\NormalTok{ \{}
\NormalTok{        evens }\OperatorTok{=}\NormalTok{ evens }\OperatorTok{+}\NormalTok{ k}
\NormalTok{    \}}

    \CommentTok{// 3. Collection iteration}
    \KeywordTok{val}\NormalTok{ weights }\OperatorTok{=} \DataTypeTok{list}\NormalTok{(}\DecValTok{3}\NormalTok{, }\DecValTok{5}\NormalTok{, }\DecValTok{8}\NormalTok{)}
    \KeywordTok{var}\NormalTok{ total\_w }\OperatorTok{=} \DecValTok{0}
    \ControlFlowTok{for}\NormalTok{ w }\ControlFlowTok{in}\NormalTok{ weights \{}
\NormalTok{        total\_w }\OperatorTok{=}\NormalTok{ total\_w }\OperatorTok{+}\NormalTok{ w}
\NormalTok{    \}}

    \CommentTok{// 4. Conditional while}
    \KeywordTok{var}\NormalTok{ countdown }\OperatorTok{=} \DecValTok{4}
    \ControlFlowTok{while}\NormalTok{ countdown }\OperatorTok{\textgreater{}} \DecValTok{0}\NormalTok{ \{}
\NormalTok{        countdown }\OperatorTok{=}\NormalTok{ countdown }\OperatorTok{{-}} \DecValTok{1}
\NormalTok{    \}}

    \FunctionTok{assert}\NormalTok{(sum }\OperatorTok{==} \DecValTok{55}\NormalTok{)}
    \FunctionTok{assert}\NormalTok{(evens }\OperatorTok{==} \DecValTok{20}\NormalTok{)}
    \FunctionTok{assert}\NormalTok{(total\_w }\OperatorTok{==} \DecValTok{16}\NormalTok{)}
    \FunctionTok{assert}\NormalTok{(countdown }\OperatorTok{==} \DecValTok{0}\NormalTok{)}
\NormalTok{    println(}\StringTok{"loops: ok"}\NormalTok{)}
    \ControlFlowTok{return} \DecValTok{0}
\NormalTok{\}}
\end{Highlighting}
\end{Shaded}

Notice that \texttt{while\ countdown\ \textgreater{}\ 0} explicitly
mutates \texttt{countdown} inside the loop body. The compiler's lint
pass warns when a \texttt{while} loop condition references variables
that never change inside the loop body (a common bug when forgetting
\texttt{i\ =\ i\ +\ 1}).

\begin{center}\rule{0.5\linewidth}{0.5pt}\end{center}

\section{4.3 Breaking, Continuing, and Loop
Labels}\label{breaking-continuing-and-loop-labels}

Inside \texttt{while}, \texttt{for}, or \texttt{loop}, use
\texttt{continue} to skip to the next iteration and \texttt{break} to
exit the innermost loop:

\begin{Shaded}
\begin{Highlighting}[]
\KeywordTok{def}\NormalTok{ find\_first\_product\_above(limit: }\DataTypeTok{i64}\NormalTok{) }\OperatorTok{{-}\textgreater{}} \DataTypeTok{i64}\NormalTok{ \{}
    \KeywordTok{var} \DataTypeTok{result} \OperatorTok{=} \OperatorTok{{-}}\DecValTok{1}
    \KeywordTok{var}\NormalTok{ row }\OperatorTok{=} \DecValTok{1}
    \ControlFlowTok{while}\NormalTok{ row }\OperatorTok{\textless{}=} \DecValTok{10}\NormalTok{ \{}
        \KeywordTok{var}\NormalTok{ col }\OperatorTok{=} \DecValTok{1}
        \ControlFlowTok{while}\NormalTok{ col }\OperatorTok{\textless{}=} \DecValTok{10}\NormalTok{ \{}
            \KeywordTok{val}\NormalTok{ prod }\OperatorTok{=}\NormalTok{ row }\OperatorTok{*}\NormalTok{ col}
            \ControlFlowTok{if}\NormalTok{ prod }\OperatorTok{\textgreater{}}\NormalTok{ limit \{}
                \DataTypeTok{result} \OperatorTok{=}\NormalTok{ prod}
                \ControlFlowTok{break}
\NormalTok{            \}}
\NormalTok{            col }\OperatorTok{=}\NormalTok{ col }\OperatorTok{+} \DecValTok{1}
\NormalTok{        \}}
        \ControlFlowTok{if} \DataTypeTok{result} \OperatorTok{\textgreater{}=} \DecValTok{0}\NormalTok{ \{}
            \ControlFlowTok{break}
\NormalTok{        \}}
\NormalTok{        row }\OperatorTok{=}\NormalTok{ row }\OperatorTok{+} \DecValTok{1}
\NormalTok{    \}}
    \ControlFlowTok{return} \DataTypeTok{result}
\NormalTok{\}}

\KeywordTok{def}\NormalTok{ main() }\OperatorTok{{-}\textgreater{}} \DataTypeTok{i64}\NormalTok{ \{}
    \KeywordTok{var}\NormalTok{ odd\_sum }\OperatorTok{=} \DecValTok{0}
    \ControlFlowTok{for}\NormalTok{ n }\ControlFlowTok{in} \DecValTok{1}\OperatorTok{..}\DecValTok{10}\NormalTok{ \{}
        \ControlFlowTok{if}\NormalTok{ n }\OperatorTok{\%} \DecValTok{2} \OperatorTok{==} \DecValTok{0}\NormalTok{ \{}
            \ControlFlowTok{continue}
\NormalTok{        \}}
\NormalTok{        odd\_sum }\OperatorTok{=}\NormalTok{ odd\_sum }\OperatorTok{+}\NormalTok{ n}
\NormalTok{    \}}
    \FunctionTok{assert}\NormalTok{(odd\_sum }\OperatorTok{==} \DecValTok{25}\NormalTok{)}
    \FunctionTok{assert}\NormalTok{(find\_first\_product\_above(}\DecValTok{45}\NormalTok{) }\OperatorTok{==} \DecValTok{50}\NormalTok{)}
\NormalTok{    println(}\StringTok{"break and continue: ok"}\NormalTok{)}
    \ControlFlowTok{return} \DecValTok{0}
\NormalTok{\}}
\end{Highlighting}
\end{Shaded}

\texttt{loop\ \{\ ...\ \}} also supports breaking directly when a
sentinel state is reached, and early \texttt{return} immediately unwinds
out of nested loops and exits the enclosing function cleanly.

\begin{center}\rule{0.5\linewidth}{0.5pt}\end{center}

\section{4.4 Logical and Bitwise
Operators}\label{logical-and-bitwise-operators}

IRIS supports both symbolic boolean operators and readable keyword
forms:

{\def\LTcaptype{none} % do not increment counter
\begin{longtable}[]{@{}
  >{\raggedright\arraybackslash}p{(\linewidth - 6\tabcolsep) * \real{0.2500}}
  >{\raggedright\arraybackslash}p{(\linewidth - 6\tabcolsep) * \real{0.2500}}
  >{\raggedright\arraybackslash}p{(\linewidth - 6\tabcolsep) * \real{0.2500}}
  >{\raggedright\arraybackslash}p{(\linewidth - 6\tabcolsep) * \real{0.2500}}@{}}
\toprule\noalign{}
\begin{minipage}[b]{\linewidth}\raggedright
Operation
\end{minipage} & \begin{minipage}[b]{\linewidth}\raggedright
Symbolic
\end{minipage} & \begin{minipage}[b]{\linewidth}\raggedright
Keyword Equivalent
\end{minipage} & \begin{minipage}[b]{\linewidth}\raggedright
Behavior
\end{minipage} \\
\midrule\noalign{}
\endhead
\bottomrule\noalign{}
\endlastfoot
Logical AND & \texttt{a\ \&\&\ b} & \texttt{a\ and\ b} & Short-circuits
if \texttt{a} is \texttt{false} \\
Logical OR &
\texttt{a\ \textbackslash{}\textbar{}\textbackslash{}\textbar{}\ b} &
\texttt{a\ or\ b} & Short-circuits if \texttt{a} is \texttt{true} \\
Logical NOT & \texttt{!a} & \texttt{not\ a} & Inverts \texttt{bool} \\
\end{longtable}
}

Bitwise operations on integers use explicit functions rather than
\texttt{\&}/\texttt{\textbar{}} symbols (which are reserved for
references and closure parameter lists):

\begin{Shaded}
\begin{Highlighting}[]
\KeywordTok{def}\NormalTok{ main() }\OperatorTok{{-}\textgreater{}} \DataTypeTok{i64}\NormalTok{ \{}
    \KeywordTok{val}\NormalTok{ flags }\OperatorTok{=} \FunctionTok{bor}\NormalTok{(}\FunctionTok{shl}\NormalTok{(}\DecValTok{1}\NormalTok{, }\DecValTok{3}\NormalTok{), }\FunctionTok{shl}\NormalTok{(}\DecValTok{1}\NormalTok{, }\DecValTok{1}\NormalTok{))   }\CommentTok{// bits 3 and 1 {-}\textgreater{} 0b1010 = 10}
    \KeywordTok{val}\NormalTok{ has\_bit\_3 }\OperatorTok{=} \FunctionTok{band}\NormalTok{(flags, }\DecValTok{8}\NormalTok{) }\OperatorTok{!=} \DecValTok{0}
    \KeywordTok{val}\NormalTok{ toggled }\OperatorTok{=} \FunctionTok{bxor}\NormalTok{(flags, }\DecValTok{2}\NormalTok{)            }\CommentTok{// 0b1000 = 8}
    \KeywordTok{val}\NormalTok{ shifted }\OperatorTok{=} \FunctionTok{shr}\NormalTok{(flags, }\DecValTok{1}\NormalTok{)             }\CommentTok{// 0b0101 = 5}

    \FunctionTok{assert}\NormalTok{(flags }\OperatorTok{==} \DecValTok{10}\NormalTok{)}
    \FunctionTok{assert}\NormalTok{(has\_bit\_3)}
    \FunctionTok{assert}\NormalTok{(toggled }\OperatorTok{==} \DecValTok{8}\NormalTok{)}
    \FunctionTok{assert}\NormalTok{(shifted }\OperatorTok{==} \DecValTok{5}\NormalTok{)}
\NormalTok{    println(}\StringTok{"logical and bitwise: ok"}\NormalTok{)}
    \ControlFlowTok{return} \DecValTok{0}
\NormalTok{\}}
\end{Highlighting}
\end{Shaded}

All bitwise builtins (\texttt{band}, \texttt{bor}, \texttt{bxor},
\texttt{shl}, \texttt{shr}, \texttt{bitnot}) are globally available
without imports.

\begin{center}\rule{0.5\linewidth}{0.5pt}\end{center}

\section{4.5 Tuple Destructuring and Expression
Tails}\label{tuple-destructuring-and-expression-tails}

Functions that return multiple values typically return a tuple
\texttt{(A,\ B)}. Destructure tuples directly in a \texttt{val} binding
or index their components with \texttt{.0}, \texttt{.1}:

\begin{Shaded}
\begin{Highlighting}[]
\KeywordTok{def}\NormalTok{ divide\_with\_rem(num: }\DataTypeTok{i64}\NormalTok{, den: }\DataTypeTok{i64}\NormalTok{) }\OperatorTok{{-}\textgreater{}}\NormalTok{ (}\DataTypeTok{i64}\NormalTok{, }\DataTypeTok{i64}\NormalTok{) \{}
    \ControlFlowTok{return}\NormalTok{ (num }\OperatorTok{/}\NormalTok{ den, num }\OperatorTok{\%}\NormalTok{ den)}
\NormalTok{\}}

\KeywordTok{def}\NormalTok{ main() }\OperatorTok{{-}\textgreater{}} \DataTypeTok{i64}\NormalTok{ \{}
    \KeywordTok{val}\NormalTok{ (q, r) }\OperatorTok{=}\NormalTok{ divide\_with\_rem(}\DecValTok{29}\NormalTok{, }\DecValTok{6}\NormalTok{)}
    \KeywordTok{val}\NormalTok{ pair }\OperatorTok{=}\NormalTok{ divide\_with\_rem(}\DecValTok{17}\NormalTok{, }\DecValTok{5}\NormalTok{)}
    \FunctionTok{assert}\NormalTok{(q }\OperatorTok{==} \DecValTok{4} \OperatorTok{\&\&}\NormalTok{ r }\OperatorTok{==} \DecValTok{5}\NormalTok{)}
    \FunctionTok{assert}\NormalTok{(pair}\OperatorTok{.}\DecValTok{0} \OperatorTok{==} \DecValTok{3} \OperatorTok{\&\&}\NormalTok{ pair}\OperatorTok{.}\DecValTok{1} \OperatorTok{==} \DecValTok{2}\NormalTok{)}

    \KeywordTok{val}\NormalTok{ adjusted }\OperatorTok{=}\NormalTok{ \{}
        \KeywordTok{val}\NormalTok{ base }\OperatorTok{=}\NormalTok{ q }\OperatorTok{*} \DecValTok{10}
\NormalTok{        base }\OperatorTok{+}\NormalTok{ r}
\NormalTok{    \}}
    \FunctionTok{assert}\NormalTok{(adjusted }\OperatorTok{==} \DecValTok{45}\NormalTok{)}
\NormalTok{    println(}\StringTok{"destructuring and expression tails: ok"}\NormalTok{)}
    \ControlFlowTok{return} \DecValTok{0}
\NormalTok{\}}
\end{Highlighting}
\end{Shaded}

\subsection{Common Mistakes}\label{common-mistakes-3}

\begin{itemize}
\tightlist
\item
  \textbf{Using \texttt{..=} in \texttt{for} loops}: \texttt{..=} is an
  inclusive range pattern for \texttt{when} arms (\texttt{1..=9}). For
  \texttt{for} loops, write exclusive \texttt{0..n} (so \texttt{0..10}
  visits \texttt{0..9}, or \texttt{0..11} visits \texttt{0..10}).
\item
  \textbf{Omitting \texttt{else} on a value-producing \texttt{if}}:
  \texttt{val\ x\ =\ if\ cond\ \{\ 1\ \}} fails with \texttt{E0119};
  supply \texttt{else\ \{\ 0\ \}}.
\item
  \textbf{Chaining comparisons
  (\texttt{0\ \textless{}\ x\ \textless{}\ 10})}: comparisons do not
  chain. Write \texttt{0\ \textless{}\ x\ \&\&\ x\ \textless{}\ 10}.
\end{itemize}

\subsection{Try It Yourself}\label{try-it-yourself-3}

\begin{enumerate}
\def\labelenumi{\arabic{enumi}.}
\tightlist
\item
  Write \texttt{fizzbuzz\_score(n:\ i64)\ -\textgreater{}\ str} using
  expression-oriented \texttt{if\ /\ else\ if\ /\ else}.
\item
  Sum all multiples of 3 or 5 below 100 using a single \texttt{for} loop
  and \texttt{continue}.
\item
  Write a function returning \texttt{(min\_val,\ max\_val)} for a
  non-empty \texttt{list\textless{}i64\textgreater{}} and unpack it with
  \texttt{val\ (lo,\ hi)\ =\ ...}.
\end{enumerate}

\begin{center}\rule{0.5\linewidth}{0.5pt}\end{center}

\part{Part II: Data, Abstraction, and
Safety}\label{part-ii-data-abstraction-and-safety}

\chapter{Chapter 5: Records, Algebraic Data Types, and
Collections}\label{chapter-5-records-algebraic-data-types-and-collections}

\section{5.1 Records and Methods}\label{records-and-methods}

Records group named fields into strongly typed data structures. Fields
may carry default values, and methods attach cleanly via \texttt{impl}
blocks (keep \texttt{impl} blocks co-located with the record and write
methods inside \texttt{impl} with \texttt{def}, not \texttt{pub\ def}):

\begin{Shaded}
\begin{Highlighting}[]
\KeywordTok{pub} \KeywordTok{record}\NormalTok{ Telemetry \{}
\NormalTok{    altitude\_m: }\DataTypeTok{f64} \OperatorTok{=} \FloatTok{0.0}\NormalTok{,}
\NormalTok{    velocity\_mps: }\DataTypeTok{f64} \OperatorTok{=} \FloatTok{0.0}\NormalTok{,}
\NormalTok{    active: }\DataTypeTok{bool} \OperatorTok{=} \ConstantTok{true}\NormalTok{,}
\NormalTok{\}}

\KeywordTok{impl}\NormalTok{ Telemetry \{}
    \KeywordTok{def}\NormalTok{ kinetic\_index(self) }\OperatorTok{{-}\textgreater{}} \DataTypeTok{f64}\NormalTok{ \{}
        \ControlFlowTok{return}\NormalTok{ self}\OperatorTok{.}\NormalTok{altitude\_m }\OperatorTok{+} \FloatTok{0.5} \OperatorTok{*}\NormalTok{ self}\OperatorTok{.}\NormalTok{velocity\_mps }\OperatorTok{*}\NormalTok{ self}\OperatorTok{.}\NormalTok{velocity\_mps}
\NormalTok{    \}}

    \KeywordTok{def}\NormalTok{ boost(self, dv: }\DataTypeTok{f64}\NormalTok{) }\OperatorTok{{-}\textgreater{}}\NormalTok{ Telemetry \{}
        \ControlFlowTok{return}\NormalTok{ Telemetry \{}
\NormalTok{            altitude\_m: self}\OperatorTok{.}\NormalTok{altitude\_m,}
\NormalTok{            velocity\_mps: self}\OperatorTok{.}\NormalTok{velocity\_mps }\OperatorTok{+}\NormalTok{ dv,}
\NormalTok{            active: self}\OperatorTok{.}\NormalTok{active,}
\NormalTok{        \}}
\NormalTok{    \}}
\NormalTok{\}}

\KeywordTok{def}\NormalTok{ main() }\OperatorTok{{-}\textgreater{}} \DataTypeTok{i64}\NormalTok{ \{}
    \KeywordTok{val}\NormalTok{ craft }\OperatorTok{=}\NormalTok{ Telemetry \{ altitude\_m: }\FloatTok{120.0}\NormalTok{, velocity\_mps: }\FloatTok{10.0}\NormalTok{ \}}
    \KeywordTok{val}\NormalTok{ boosted }\OperatorTok{=}\NormalTok{ craft}\OperatorTok{.}\NormalTok{boost(}\FloatTok{4.0}\NormalTok{)}
    \FunctionTok{assert}\NormalTok{(craft}\OperatorTok{.}\NormalTok{kinetic\_index() }\OperatorTok{==} \FloatTok{170.0}\NormalTok{)}
    \FunctionTok{assert}\NormalTok{(boosted}\OperatorTok{.}\NormalTok{velocity\_mps }\OperatorTok{==} \FloatTok{14.0}\NormalTok{)}
    \FunctionTok{assert}\NormalTok{(boosted}\OperatorTok{.}\NormalTok{active)}
\NormalTok{    println(}\StringTok{"records and methods: ok"}\NormalTok{)}
    \ControlFlowTok{return} \DecValTok{0}
\NormalTok{\}}
\end{Highlighting}
\end{Shaded}

When constructing a record, any field with a default value can be
omitted cleanly, which keeps state-record updates concise as data
schemas grow.

\begin{center}\rule{0.5\linewidth}{0.5pt}\end{center}

\section{5.2 Algebraic Data Types with Choice
Enums}\label{algebraic-data-types-with-choice-enums}

Use \texttt{choice} to model values that take one of several mutually
exclusive variants---optionally carrying typed payloads:

\begin{Shaded}
\begin{Highlighting}[]
\KeywordTok{choice}\NormalTok{ Alert \{}
\NormalTok{    Nominal,}
\NormalTok{    Throttle(}\DataTypeTok{f64}\NormalTok{),}
\NormalTok{    Fault(}\DataTypeTok{i64}\NormalTok{, }\DataTypeTok{str}\NormalTok{),}
\NormalTok{\}}

\KeywordTok{def}\NormalTok{ render\_alert(a: Alert) }\OperatorTok{{-}\textgreater{}} \DataTypeTok{str}\NormalTok{ \{}
    \ControlFlowTok{return} \ControlFlowTok{when}\NormalTok{ a \{}
\NormalTok{        Alert}\OperatorTok{.}\NormalTok{Nominal }\OperatorTok{=\textgreater{}} \StringTok{"OK"}\NormalTok{,}
\NormalTok{        Alert}\OperatorTok{.}\NormalTok{Throttle(ratio) }\OperatorTok{=\textgreater{}} \StringTok{f"THROTTLE=}\ControlFlowTok{\{}\NormalTok{ratio}\ControlFlowTok{\}}\StringTok{"}\NormalTok{,}
\NormalTok{        Alert}\OperatorTok{.}\NormalTok{Fault(code, msg) }\OperatorTok{=\textgreater{}} \StringTok{f"ERR(}\ControlFlowTok{\{}\NormalTok{code}\ControlFlowTok{\}}\StringTok{): }\ControlFlowTok{\{}\NormalTok{msg}\ControlFlowTok{\}}\StringTok{"}\NormalTok{,}
\NormalTok{    \}}
\NormalTok{\}}

\KeywordTok{def}\NormalTok{ main() }\OperatorTok{{-}\textgreater{}} \DataTypeTok{i64}\NormalTok{ \{}
    \KeywordTok{val}\NormalTok{ ok\_state }\OperatorTok{=}\NormalTok{ Alert}\OperatorTok{.}\NormalTok{Nominal}
    \KeywordTok{val}\NormalTok{ warn\_state }\OperatorTok{=}\NormalTok{ Alert}\OperatorTok{.}\NormalTok{Throttle(}\FloatTok{0.75}\NormalTok{)}
    \KeywordTok{val}\NormalTok{ err\_state }\OperatorTok{=}\NormalTok{ Alert}\OperatorTok{.}\NormalTok{Fault(}\DecValTok{42}\NormalTok{, }\StringTok{"sensor timeout"}\NormalTok{)}

    \FunctionTok{assert}\NormalTok{(ok\_state }\OperatorTok{==}\NormalTok{ Alert}\OperatorTok{.}\NormalTok{Nominal)}
    \FunctionTok{assert}\NormalTok{(render\_alert(warn\_state) }\OperatorTok{==} \StringTok{"THROTTLE=0.75"}\NormalTok{)}
    \FunctionTok{assert}\NormalTok{(render\_alert(err\_state) }\OperatorTok{==} \StringTok{"ERR(42): sensor timeout"}\NormalTok{)}
\NormalTok{    println(}\StringTok{"choice enums: ok"}\NormalTok{)}
    \ControlFlowTok{return} \DecValTok{0}
\NormalTok{\}}
\end{Highlighting}
\end{Shaded}

Note that \texttt{choice} declarations represent concrete closed variant
families (they do not take generic \texttt{{[}T{]}} type parameters; use
\texttt{record\ Name{[}T{]}} or
\texttt{option\textless{}T\textgreater{}} /
\texttt{result\textless{}T,\ E\textgreater{}} when you need parametric
containers). Structural \texttt{==} and \texttt{!=} work directly on
\texttt{choice} values.

\begin{center}\rule{0.5\linewidth}{0.5pt}\end{center}

\section{5.3 Tuples and Fixed Arrays}\label{tuples-and-fixed-arrays}

IRIS provides two stack-friendly fixed-structure aggregates:

\begin{itemize}
\tightlist
\item
  \textbf{Tuples \texttt{(T1,\ T2,\ ...)}}: lightweight heterogeneous
  grouping accessed by \texttt{.0}, \texttt{.1}, or destructured via
  \texttt{val\ (a,\ b)\ =\ pair}.
\item
  \textbf{Fixed arrays \texttt{{[}T;\ N{]}}}: fixed-length homogeneous
  sequences with compile-time size \texttt{N}. Mutate elements in place
  via indexing when bound with \texttt{var}, and read length with
  \texttt{len(arr)}. Out-of-bounds access halts safely with runtime
  diagnostic \texttt{E0403}.
\end{itemize}

\begin{Shaded}
\begin{Highlighting}[]
\KeywordTok{def}\NormalTok{ main() }\OperatorTok{{-}\textgreater{}} \DataTypeTok{i64}\NormalTok{ \{}
    \KeywordTok{var}\NormalTok{ samples: [}\DataTypeTok{i64}\NormalTok{; }\DecValTok{4}\NormalTok{] }\OperatorTok{=}\NormalTok{ [}\DecValTok{10}\NormalTok{, }\DecValTok{20}\NormalTok{, }\DecValTok{30}\NormalTok{, }\DecValTok{40}\NormalTok{]}
\NormalTok{    samples[}\DecValTok{2}\NormalTok{] }\OperatorTok{=} \DecValTok{99}
    \FunctionTok{assert}\NormalTok{(}\FunctionTok{len}\NormalTok{(samples) }\OperatorTok{==} \DecValTok{4}\NormalTok{)}
    \FunctionTok{assert}\NormalTok{(samples[}\DecValTok{2}\NormalTok{] }\OperatorTok{==} \DecValTok{99}\NormalTok{)}

    \KeywordTok{val}\NormalTok{ reading }\OperatorTok{=}\NormalTok{ (}\StringTok{"alpha"}\NormalTok{, }\DecValTok{256}\NormalTok{, }\ConstantTok{true}\NormalTok{)}
    \FunctionTok{assert}\NormalTok{(reading}\OperatorTok{.}\DecValTok{0} \OperatorTok{==} \StringTok{"alpha"} \OperatorTok{\&\&}\NormalTok{ reading}\OperatorTok{.}\DecValTok{1} \OperatorTok{==} \DecValTok{256} \OperatorTok{\&\&}\NormalTok{ reading}\OperatorTok{.}\DecValTok{2}\NormalTok{)}
\NormalTok{    println(}\StringTok{"tuples and arrays: ok"}\NormalTok{)}
    \ControlFlowTok{return} \DecValTok{0}
\NormalTok{\}}
\end{Highlighting}
\end{Shaded}

\begin{center}\rule{0.5\linewidth}{0.5pt}\end{center}

\section{5.4 Dynamic Lists and Maps}\label{dynamic-lists-and-maps}

When sizes vary dynamically at runtime, reach for
\texttt{list\textless{}T\textgreater{}} and
\texttt{map\textless{}K,\ V\textgreater{}}, which are reference-counted
heap collections backed by global builtins:

\begin{Shaded}
\begin{Highlighting}[]
\KeywordTok{def}\NormalTok{ main() }\OperatorTok{{-}\textgreater{}} \DataTypeTok{i64}\NormalTok{ \{}
    \KeywordTok{val}\NormalTok{ nums: }\DataTypeTok{list}\OperatorTok{\textless{}}\DataTypeTok{i64}\OperatorTok{\textgreater{}} \OperatorTok{=} \DataTypeTok{list}\NormalTok{(}\DecValTok{10}\NormalTok{, }\DecValTok{20}\NormalTok{)}
\NormalTok{    list\_push(nums, }\DecValTok{30}\NormalTok{)}
    \FunctionTok{list\_set}\NormalTok{(nums, }\DecValTok{1}\NormalTok{, }\DecValTok{25}\NormalTok{)}
    \FunctionTok{assert}\NormalTok{(}\FunctionTok{list\_len}\NormalTok{(nums) }\OperatorTok{==} \DecValTok{3}\NormalTok{)}
    \FunctionTok{assert}\NormalTok{(}\FunctionTok{list\_get}\NormalTok{(nums, }\DecValTok{1}\NormalTok{) }\OperatorTok{==} \DecValTok{25}\NormalTok{)}

    \KeywordTok{val}\NormalTok{ popped }\OperatorTok{=} \FunctionTok{list\_pop}\NormalTok{(nums)}
    \FunctionTok{assert}\NormalTok{(popped }\OperatorTok{==} \DecValTok{30}\NormalTok{)}
    \FunctionTok{assert}\NormalTok{(}\FunctionTok{list\_len}\NormalTok{(nums) }\OperatorTok{==} \DecValTok{2}\NormalTok{)}

    \KeywordTok{val}\NormalTok{ ports: }\DataTypeTok{map}\OperatorTok{\textless{}}\DataTypeTok{str}\NormalTok{, }\DataTypeTok{i64}\OperatorTok{\textgreater{}} \OperatorTok{=} \DataTypeTok{map}\NormalTok{()}
    \FunctionTok{map\_set}\NormalTok{(ports, }\StringTok{"http"}\NormalTok{, }\DecValTok{8080}\NormalTok{)}
    \FunctionTok{map\_set}\NormalTok{(ports, }\StringTok{"metrics"}\NormalTok{, }\DecValTok{9090}\NormalTok{)}
    \FunctionTok{assert}\NormalTok{(}\FunctionTok{map\_contains}\NormalTok{(ports, }\StringTok{"http"}\NormalTok{))}
    \FunctionTok{assert}\NormalTok{(unwrap(}\FunctionTok{map\_get}\NormalTok{(ports, }\StringTok{"http"}\NormalTok{)) }\OperatorTok{==} \DecValTok{8080}\NormalTok{)}

    \KeywordTok{val}\NormalTok{ keys }\OperatorTok{=} \FunctionTok{map\_keys}\NormalTok{(ports)}
    \FunctionTok{assert}\NormalTok{(}\FunctionTok{list\_len}\NormalTok{(keys) }\OperatorTok{==} \DecValTok{2}\NormalTok{)}
\NormalTok{    println(}\StringTok{"lists and maps: ok"}\NormalTok{)}
    \ControlFlowTok{return} \DecValTok{0}
\NormalTok{\}}
\end{Highlighting}
\end{Shaded}

Notice \texttt{map\_get(m,\ key)} returns
\texttt{option\textless{}V\textgreater{}} (\texttt{some(value)} or
\texttt{none()}), preventing missing-key crashes. Notice also that
constructing an empty \texttt{list()} or \texttt{map()} without
initializers requires a binding type annotation such as
\texttt{:\ list\textless{}i64\textgreater{}} or
\texttt{:\ map\textless{}str,\ i64\textgreater{}} (\texttt{E0207}
otherwise).

\begin{center}\rule{0.5\linewidth}{0.5pt}\end{center}

\section{5.5 Sets, Deques, Heaps, Queues, BitSets, and
Buffers}\label{sets-deques-heaps-queues-bitsets-and-buffers}

Beyond lists and maps, the embedded standard library supplies
specialized containers tuned for distinct access patterns:

{\def\LTcaptype{none} % do not increment counter
\begin{longtable}[]{@{}
  >{\raggedright\arraybackslash}p{(\linewidth - 4\tabcolsep) * \real{0.3333}}
  >{\raggedright\arraybackslash}p{(\linewidth - 4\tabcolsep) * \real{0.3333}}
  >{\raggedright\arraybackslash}p{(\linewidth - 4\tabcolsep) * \real{0.3333}}@{}}
\toprule\noalign{}
\begin{minipage}[b]{\linewidth}\raggedright
Module
\end{minipage} & \begin{minipage}[b]{\linewidth}\raggedright
Types / Constructors
\end{minipage} & \begin{minipage}[b]{\linewidth}\raggedright
Core Operations
\end{minipage} \\
\midrule\noalign{}
\endhead
\bottomrule\noalign{}
\endlastfoot
\texttt{std.set} & \texttt{Set\textless{}T\textgreater{}},
\texttt{set\_new()}, \texttt{set\_from\_list(xs)} & \texttt{set\_add},
\texttt{set\_remove}, \texttt{set\_contains}, \texttt{set\_union},
\texttt{set\_intersection}, \texttt{set\_difference},
\texttt{set\_len} \\
\texttt{std.deque} & \texttt{deque\_create()} (backed by
\texttt{list\textless{}i64\textgreater{}}) & \texttt{push\_front},
\texttt{push\_back}, \texttt{pop\_front}, \texttt{pop\_back},
\texttt{peek\_front}, \texttt{peek\_back}, \texttt{deque\_size} \\
\texttt{std.heap} & \texttt{MinHeap\textless{}T\textgreater{}},
\texttt{heap\_new()}, \texttt{heap\_from\_list(xs)} &
\texttt{heap\_push}, \texttt{heap\_pop}, \texttt{heap\_pop\_val},
\texttt{heap\_peek}, \texttt{heap\_len}, \texttt{heap\_merge} \\
\texttt{std.queue} & \texttt{Queue}, \texttt{queue\_new()} &
\texttt{enqueue}, \texttt{dequeue\_val}, \texttt{dequeue\_queue},
\texttt{queue\_peek}, \texttt{queue\_len} \\
\texttt{std.bitset} & \texttt{bitset\_create()} & \texttt{set\_bit},
\texttt{get\_bit}, \texttt{clear\_bit}, \texttt{toggle\_bit},
\texttt{popcount}, \texttt{bitset\_union},
\texttt{bitset\_intersection} \\
\texttt{std.buffer} & \texttt{ByteBuffer},
\texttt{byte\_buffer\_from\_str(s)}, \texttt{byte\_buffer\_from\_hex(h)}
& \texttt{byte\_buffer\_read\_u16\_be},
\texttt{byte\_buffer\_read\_u32\_be}, \texttt{byte\_buffer\_slice},
\texttt{byte\_buffer\_to\_hex} \\
\end{longtable}
}

Here is a compact example exercising sets and bitsets side by side:

\begin{Shaded}
\begin{Highlighting}[]
\KeywordTok{bring}\NormalTok{ std}\OperatorTok{.}\NormalTok{set}
\KeywordTok{bring}\NormalTok{ std}\OperatorTok{.}\DataTypeTok{bitset}

\KeywordTok{def}\NormalTok{ main() }\OperatorTok{{-}\textgreater{}} \DataTypeTok{i64}\NormalTok{ \{}
    \KeywordTok{var}\NormalTok{ tags: Set}\OperatorTok{\textless{}}\DataTypeTok{str}\OperatorTok{\textgreater{}} \OperatorTok{=}\NormalTok{ set\_new()}
\NormalTok{    tags }\OperatorTok{=}\NormalTok{ set\_add(tags, }\StringTok{"telemetry"}\NormalTok{)}
\NormalTok{    tags }\OperatorTok{=}\NormalTok{ set\_add(tags, }\StringTok{"vision"}\NormalTok{)}
\NormalTok{    tags }\OperatorTok{=}\NormalTok{ set\_add(tags, }\StringTok{"telemetry"}\NormalTok{)}
    \FunctionTok{assert}\NormalTok{(set\_len(tags) }\OperatorTok{==} \DecValTok{2}\NormalTok{)}
    \FunctionTok{assert}\NormalTok{(set\_contains(tags, }\StringTok{"vision"}\NormalTok{))}

    \KeywordTok{var}\NormalTok{ flags }\OperatorTok{=}\NormalTok{ bitset\_create()}
\NormalTok{    flags }\OperatorTok{=}\NormalTok{ set\_bit(flags, }\DecValTok{2}\NormalTok{)}
\NormalTok{    flags }\OperatorTok{=}\NormalTok{ set\_bit(flags, }\DecValTok{7}\NormalTok{)}
    \FunctionTok{assert}\NormalTok{(get\_bit(flags, }\DecValTok{2}\NormalTok{))}
    \FunctionTok{assert}\NormalTok{(}\OperatorTok{!}\NormalTok{get\_bit(flags, }\DecValTok{3}\NormalTok{))}
    \FunctionTok{assert}\NormalTok{(popcount(flags) }\OperatorTok{==} \DecValTok{2}\NormalTok{)}

\NormalTok{    println(}\StringTok{"specialized collections: ok"}\NormalTok{)}
    \ControlFlowTok{return} \DecValTok{0}
\NormalTok{\}}
\end{Highlighting}
\end{Shaded}

\begin{center}\rule{0.5\linewidth}{0.5pt}\end{center}

\section{5.6 Strings, Formatting, Regular Expressions, and
DateTime}\label{strings-formatting-regular-expressions-and-datetime}

IRIS includes UTF-8 string primitives globally, rich Unicode/text
helpers in \texttt{std.string}, \texttt{std.unicode}, and
\texttt{std.fmt}, regex builtins, and calendar/clock utilities in
\texttt{std.time}:

\begin{Shaded}
\begin{Highlighting}[]
\KeywordTok{bring}\NormalTok{ std}\OperatorTok{.}\NormalTok{string}
\KeywordTok{bring}\NormalTok{ std}\OperatorTok{.}\NormalTok{fmt}
\KeywordTok{bring}\NormalTok{ std}\OperatorTok{.}\NormalTok{time}

\KeywordTok{def}\NormalTok{ main() }\OperatorTok{{-}\textgreater{}} \DataTypeTok{i64}\NormalTok{ \{}
    \CommentTok{// Core string builtins + std.string}
    \KeywordTok{val}\NormalTok{ raw }\OperatorTok{=} \StringTok{"  iris,vision,control  "}
    \KeywordTok{val}\NormalTok{ cleaned }\OperatorTok{=} \FunctionTok{trim}\NormalTok{(raw)}
    \KeywordTok{val}\NormalTok{ parts }\OperatorTok{=} \FunctionTok{split}\NormalTok{(cleaned, }\StringTok{","}\NormalTok{)}
    \FunctionTok{assert}\NormalTok{(}\FunctionTok{list\_len}\NormalTok{(parts) }\OperatorTok{==} \DecValTok{3}\NormalTok{)}
    \FunctionTok{assert}\NormalTok{(str\_join(parts, }\StringTok{" | "}\NormalTok{) }\OperatorTok{==} \StringTok{"iris | vision | control"}\NormalTok{)}
    \FunctionTok{assert}\NormalTok{(pad\_left(}\StringTok{"42"}\NormalTok{, }\DecValTok{5}\NormalTok{, }\StringTok{"0"}\NormalTok{) }\OperatorTok{==} \StringTok{"00042"}\NormalTok{)}

    \CommentTok{// Built{-}in regular expressions}
    \FunctionTok{assert}\NormalTok{(}\FunctionTok{regex\_match}\NormalTok{(}\StringTok{"\textbackslash{}\textbackslash{}d+"}\NormalTok{, }\StringTok{"port=8080"}\NormalTok{))}
    \KeywordTok{val}\NormalTok{ nums }\OperatorTok{=} \FunctionTok{regex\_find\_all}\NormalTok{(}\StringTok{"\textbackslash{}\textbackslash{}d+"}\NormalTok{, }\StringTok{"a12 b34 c56"}\NormalTok{)}
    \FunctionTok{assert}\NormalTok{(}\FunctionTok{list\_len}\NormalTok{(nums) }\OperatorTok{==} \DecValTok{3}\NormalTok{)}
    \FunctionTok{assert}\NormalTok{(regex\_replace\_all(}\StringTok{"\textbackslash{}\textbackslash{}d+"}\NormalTok{, }\StringTok{"v1 {-}\textgreater{} v2"}\NormalTok{, }\StringTok{"X"}\NormalTok{) }\OperatorTok{==} \StringTok{"vX {-}\textgreater{} vX"}\NormalTok{)}

    \CommentTok{// Formatting tables / sprintf and ISO{-}8601 UTC date formatting}
    \KeywordTok{val}\NormalTok{ summary }\OperatorTok{=}\NormalTok{ sprintf(}\StringTok{"node=\%s count=\%d"}\NormalTok{, }\DataTypeTok{list}\NormalTok{(}\StringTok{"edge{-}1"}\NormalTok{, }\StringTok{"7"}\NormalTok{))}
    \FunctionTok{assert}\NormalTok{(summary }\OperatorTok{==} \StringTok{"node=edge{-}1 count=7"}\NormalTok{)}
    \KeywordTok{val}\NormalTok{ iso }\OperatorTok{=}\NormalTok{ format\_iso8601(}\DecValTok{0}\NormalTok{)}
    \FunctionTok{assert}\NormalTok{(}\FunctionTok{starts\_with}\NormalTok{(iso, }\StringTok{"1970{-}01{-}01"}\NormalTok{))}

\NormalTok{    println(}\StringTok{"strings, regex, and time: ok"}\NormalTok{)}
    \ControlFlowTok{return} \DecValTok{0}
\NormalTok{\}}
\end{Highlighting}
\end{Shaded}

For Unicode code-point iteration and UTF-8 scalar handling beyond ASCII
byte helpers, import \texttt{std.unicode}.

\subsection{Common Mistakes}\label{common-mistakes-4}

\begin{itemize}
\tightlist
\item
  \textbf{Treating \texttt{set\_add} or \texttt{set\_bit} as mutating in
  place}: functional-style stdlib collection helpers (\texttt{set\_add},
  \texttt{heap\_push}, \texttt{set\_bit}, \texttt{enqueue}) return the
  updated value. Reassign \texttt{s\ =\ set\_add(s,\ x)} or thread the
  returned value forward.
\item
  \textbf{Writing \texttt{regex\_replace} instead of
  \texttt{regex\_replace\_all}}: the built-in substitution function is
  \texttt{regex\_replace\_all(pattern,\ text,\ replacement)}.
\item
  \textbf{Mutating fields directly on immutable record instances}:
  records follow value-update style---construct an updated copy cleanly
  (as in \texttt{boost} above), or keep mutable aggregates inside
  explicit handles when needed.
\end{itemize}

\subsection{Try It Yourself}\label{try-it-yourself-4}

\begin{enumerate}
\def\labelenumi{\arabic{enumi}.}
\tightlist
\item
  Model a simple order book or task queue with
  \texttt{record\ Task\ \{\ id:\ i64,\ priority:\ i64,\ title:\ str\ \}}
  and a
  \texttt{choice\ Status\ \{\ Queued,\ Running,\ Done,\ Failed(str)\ \}}.
\item
  Write a program using \texttt{MinHeap\textless{}i64\textgreater{}}
  from \texttt{std.heap} to push five out-of-order integers and pop them
  into ascending sorted order verified with \texttt{assert}.
\item
  Parse \texttt{"temp=23.5;hum=61.0;node=alpha"} with
  \texttt{regex\_find\_all} or \texttt{split} and
  \texttt{std.string.strip\_prefix}, verifying every field with
  \texttt{assert}.
\end{enumerate}

\begin{center}\rule{0.5\linewidth}{0.5pt}\end{center}

\chapter{Chapter 6: Pattern Matching with
When}\label{chapter-6-pattern-matching-with-when}

\section{6.1 Matching Choice Variants}\label{matching-choice-variants}

The \texttt{when} expression dispatches on a value's structure, binding
variant payloads directly in each arm:

\begin{Shaded}
\begin{Highlighting}[]
\KeywordTok{choice}\NormalTok{ Packet \{}
\NormalTok{    Ping(}\DataTypeTok{i64}\NormalTok{),}
\NormalTok{    Data(}\DataTypeTok{str}\NormalTok{, }\DataTypeTok{i64}\NormalTok{),}
\NormalTok{    Close,}
\NormalTok{\}}

\KeywordTok{def}\NormalTok{ handle\_packet(pkt: Packet) }\OperatorTok{{-}\textgreater{}} \DataTypeTok{i64}\NormalTok{ \{}
    \ControlFlowTok{return} \ControlFlowTok{when}\NormalTok{ pkt \{}
\NormalTok{        Packet}\OperatorTok{.}\NormalTok{Ping(seq) }\OperatorTok{=\textgreater{}}\NormalTok{ seq,}
\NormalTok{        Packet}\OperatorTok{.}\NormalTok{Data(topic, bytes) }\OperatorTok{=\textgreater{}} \FunctionTok{len}\NormalTok{(topic) }\OperatorTok{+}\NormalTok{ bytes,}
\NormalTok{        Packet}\OperatorTok{.}\NormalTok{Close }\OperatorTok{=\textgreater{}} \DecValTok{0}\NormalTok{,}
\NormalTok{    \}}
\NormalTok{\}}

\KeywordTok{def}\NormalTok{ main() }\OperatorTok{{-}\textgreater{}} \DataTypeTok{i64}\NormalTok{ \{}
    \FunctionTok{assert}\NormalTok{(handle\_packet(Packet}\OperatorTok{.}\NormalTok{Ping(}\DecValTok{12}\NormalTok{)) }\OperatorTok{==} \DecValTok{12}\NormalTok{)}
    \FunctionTok{assert}\NormalTok{(handle\_packet(Packet}\OperatorTok{.}\NormalTok{Data(}\StringTok{"imu"}\NormalTok{, }\DecValTok{64}\NormalTok{)) }\OperatorTok{==} \DecValTok{67}\NormalTok{)}
    \FunctionTok{assert}\NormalTok{(handle\_packet(Packet}\OperatorTok{.}\NormalTok{Close) }\OperatorTok{==} \DecValTok{0}\NormalTok{)}
\NormalTok{    println(}\StringTok{"when choice matching: ok"}\NormalTok{)}
    \ControlFlowTok{return} \DecValTok{0}
\NormalTok{\}}
\end{Highlighting}
\end{Shaded}

\begin{center}\rule{0.5\linewidth}{0.5pt}\end{center}

\section{6.2 Matching Options and
Results}\label{matching-options-and-results}

Instead of risking panics with \texttt{unwrap(...)}, handle
\texttt{option\textless{}T\textgreater{}} and
\texttt{result\textless{}T,\ E\textgreater{}} explicitly through
\texttt{when}:

\begin{Shaded}
\begin{Highlighting}[]
\KeywordTok{def}\NormalTok{ parse\_port(raw: }\DataTypeTok{str}\NormalTok{) }\OperatorTok{{-}\textgreater{}} \DataTypeTok{result}\OperatorTok{\textless{}}\DataTypeTok{i64}\NormalTok{, }\DataTypeTok{str}\OperatorTok{\textgreater{}}\NormalTok{ \{}
    \KeywordTok{val}\NormalTok{ parsed: }\DataTypeTok{option}\OperatorTok{\textless{}}\DataTypeTok{i64}\OperatorTok{\textgreater{}} \OperatorTok{=} \FunctionTok{parse\_i64}\NormalTok{(raw)}
    \ControlFlowTok{return} \ControlFlowTok{when}\NormalTok{ parsed \{}
        \FunctionTok{some}\NormalTok{(p) }\OperatorTok{=\textgreater{}}\NormalTok{ \{}
            \ControlFlowTok{if}\NormalTok{ p }\OperatorTok{\textgreater{}=} \DecValTok{1} \OperatorTok{\&\&}\NormalTok{ p }\OperatorTok{\textless{}=} \DecValTok{65535}\NormalTok{ \{}
                \FunctionTok{ok}\NormalTok{(p)}
\NormalTok{            \} }\ControlFlowTok{else}\NormalTok{ \{}
                \FunctionTok{err}\NormalTok{(}\StringTok{"port out of range"}\NormalTok{)}
\NormalTok{            \}}
\NormalTok{        \},}
        \ConstantTok{none} \OperatorTok{=\textgreater{}} \FunctionTok{err}\NormalTok{(}\StringTok{"not an integer"}\NormalTok{),}
\NormalTok{    \}}
\NormalTok{\}}

\KeywordTok{def}\NormalTok{ describe\_port(r: }\DataTypeTok{result}\OperatorTok{\textless{}}\DataTypeTok{i64}\NormalTok{, }\DataTypeTok{str}\OperatorTok{\textgreater{}}\NormalTok{) }\OperatorTok{{-}\textgreater{}} \DataTypeTok{str}\NormalTok{ \{}
    \ControlFlowTok{return} \ControlFlowTok{when}\NormalTok{ r \{}
        \FunctionTok{ok}\NormalTok{(port) }\OperatorTok{=\textgreater{}} \StringTok{f"listening on }\ControlFlowTok{\{}\NormalTok{port}\ControlFlowTok{\}}\StringTok{"}\NormalTok{,}
        \FunctionTok{err}\NormalTok{(reason) }\OperatorTok{=\textgreater{}} \StringTok{f"invalid config: }\ControlFlowTok{\{}\NormalTok{reason}\ControlFlowTok{\}}\StringTok{"}\NormalTok{,}
\NormalTok{    \}}
\NormalTok{\}}

\KeywordTok{def}\NormalTok{ main() }\OperatorTok{{-}\textgreater{}} \DataTypeTok{i64}\NormalTok{ \{}
    \FunctionTok{assert}\NormalTok{(describe\_port(parse\_port(}\StringTok{"8080"}\NormalTok{)) }\OperatorTok{==} \StringTok{"listening on 8080"}\NormalTok{)}
    \FunctionTok{assert}\NormalTok{(describe\_port(parse\_port(}\StringTok{"99999"}\NormalTok{)) }\OperatorTok{==} \StringTok{"invalid config: port out of range"}\NormalTok{)}
    \FunctionTok{assert}\NormalTok{(describe\_port(parse\_port(}\StringTok{"abc"}\NormalTok{)) }\OperatorTok{==} \StringTok{"invalid config: not an integer"}\NormalTok{)}
\NormalTok{    println(}\StringTok{"option and result matching: ok"}\NormalTok{)}
    \ControlFlowTok{return} \DecValTok{0}
\NormalTok{\}}
\end{Highlighting}
\end{Shaded}

Note an important IRIS standard-library rule illustrated here:
\texttt{parse\_i64(s)} and \texttt{parse\_f64(s)} return
\textbf{\texttt{option\textless{}i64\textgreater{}} /
\texttt{option\textless{}f64\textgreater{}}} (not \texttt{result}), so
match \texttt{some(n)} and \texttt{none}---or coalesce directly with
\texttt{??}.

\begin{center}\rule{0.5\linewidth}{0.5pt}\end{center}

\section{6.3 Literals, Inclusive Ranges, Tuples, and
Wildcards}\label{literals-inclusive-ranges-tuples-and-wildcards}

\texttt{when} also matches integers, booleans, strings, inclusive
integer ranges (\texttt{lo..=hi}), flat tuples, and the catch-all
wildcard \texttt{\_}:

\begin{Shaded}
\begin{Highlighting}[]
\KeywordTok{def}\NormalTok{ http\_category(status: }\DataTypeTok{i64}\NormalTok{) }\OperatorTok{{-}\textgreater{}} \DataTypeTok{str}\NormalTok{ \{}
    \ControlFlowTok{return} \ControlFlowTok{when}\NormalTok{ status \{}
        \DecValTok{200} \OperatorTok{=\textgreater{}} \StringTok{"OK"}\NormalTok{,}
        \DecValTok{201}\OperatorTok{..=}\DecValTok{299} \OperatorTok{=\textgreater{}} \StringTok{"Success"}\NormalTok{,}
        \DecValTok{400}\OperatorTok{..=}\DecValTok{499} \OperatorTok{=\textgreater{}} \StringTok{"Client Error"}\NormalTok{,}
        \DecValTok{500}\OperatorTok{..=}\DecValTok{599} \OperatorTok{=\textgreater{}} \StringTok{"Server Error"}\NormalTok{,}
\NormalTok{        \_ }\OperatorTok{=\textgreater{}} \StringTok{"Other"}\NormalTok{,}
\NormalTok{    \}}
\NormalTok{\}}

\KeywordTok{def}\NormalTok{ quadrant(x\_pos: }\DataTypeTok{bool}\NormalTok{, y\_pos: }\DataTypeTok{bool}\NormalTok{) }\OperatorTok{{-}\textgreater{}} \DataTypeTok{i64}\NormalTok{ \{}
    \ControlFlowTok{return} \ControlFlowTok{when}\NormalTok{ (x\_pos, y\_pos) \{}
\NormalTok{        (}\ConstantTok{true}\NormalTok{, }\ConstantTok{true}\NormalTok{) }\OperatorTok{=\textgreater{}} \DecValTok{1}\NormalTok{,}
\NormalTok{        (}\ConstantTok{false}\NormalTok{, }\ConstantTok{true}\NormalTok{) }\OperatorTok{=\textgreater{}} \DecValTok{2}\NormalTok{,}
\NormalTok{        (}\ConstantTok{false}\NormalTok{, }\ConstantTok{false}\NormalTok{) }\OperatorTok{=\textgreater{}} \DecValTok{3}\NormalTok{,}
\NormalTok{        (}\ConstantTok{true}\NormalTok{, }\ConstantTok{false}\NormalTok{) }\OperatorTok{=\textgreater{}} \DecValTok{4}\NormalTok{,}
\NormalTok{    \}}
\NormalTok{\}}

\KeywordTok{def}\NormalTok{ main() }\OperatorTok{{-}\textgreater{}} \DataTypeTok{i64}\NormalTok{ \{}
    \FunctionTok{assert}\NormalTok{(http\_category(}\DecValTok{200}\NormalTok{) }\OperatorTok{==} \StringTok{"OK"}\NormalTok{)}
    \FunctionTok{assert}\NormalTok{(http\_category(}\DecValTok{204}\NormalTok{) }\OperatorTok{==} \StringTok{"Success"}\NormalTok{)}
    \FunctionTok{assert}\NormalTok{(http\_category(}\DecValTok{404}\NormalTok{) }\OperatorTok{==} \StringTok{"Client Error"}\NormalTok{)}
    \FunctionTok{assert}\NormalTok{(quadrant(}\ConstantTok{true}\NormalTok{, }\ConstantTok{false}\NormalTok{) }\OperatorTok{==} \DecValTok{4}\NormalTok{)}
\NormalTok{    println(}\StringTok{"ranges and tuple matching: ok"}\NormalTok{)}
    \ControlFlowTok{return} \DecValTok{0}
\NormalTok{\}}
\end{Highlighting}
\end{Shaded}

\begin{center}\rule{0.5\linewidth}{0.5pt}\end{center}

\section{6.4 Match Guards and Value-Producing
Arms}\label{match-guards-and-value-producing-arms}

Add conditional precision with \texttt{if\ guard} clauses on any arm, or
compute multi-statement blocks inside braced arms whose last expression
produces the arm's result:

\begin{Shaded}
\begin{Highlighting}[]
\KeywordTok{choice}\NormalTok{ Measurement \{}
\NormalTok{    Valid(}\DataTypeTok{f64}\NormalTok{),}
\NormalTok{    Missing,}
\NormalTok{\}}

\KeywordTok{def}\NormalTok{ score\_reading(m: Measurement, limit: }\DataTypeTok{f64}\NormalTok{) }\OperatorTok{{-}\textgreater{}} \DataTypeTok{str}\NormalTok{ \{}
    \ControlFlowTok{return} \ControlFlowTok{when}\NormalTok{ m \{}
\NormalTok{        Measurement}\OperatorTok{.}\NormalTok{Valid(v) }\ControlFlowTok{if}\NormalTok{ v }\OperatorTok{\textless{}} \FloatTok{0.0} \OperatorTok{=\textgreater{}} \StringTok{"negative"}\NormalTok{,}
\NormalTok{        Measurement}\OperatorTok{.}\NormalTok{Valid(v) }\ControlFlowTok{if}\NormalTok{ v }\OperatorTok{\textgreater{}}\NormalTok{ limit }\OperatorTok{=\textgreater{}}\NormalTok{ \{}
            \KeywordTok{val}\NormalTok{ excess }\OperatorTok{=}\NormalTok{ v }\OperatorTok{{-}}\NormalTok{ limit}
            \StringTok{f"over by }\ControlFlowTok{\{}\NormalTok{excess}\ControlFlowTok{\}}\StringTok{"}
\NormalTok{        \},}
\NormalTok{        Measurement}\OperatorTok{.}\NormalTok{Valid(\_) }\OperatorTok{=\textgreater{}} \StringTok{"nominal"}\NormalTok{,}
\NormalTok{        Measurement}\OperatorTok{.}\NormalTok{Missing }\OperatorTok{=\textgreater{}} \StringTok{"missing"}\NormalTok{,}
\NormalTok{    \}}
\NormalTok{\}}

\KeywordTok{def}\NormalTok{ main() }\OperatorTok{{-}\textgreater{}} \DataTypeTok{i64}\NormalTok{ \{}
    \FunctionTok{assert}\NormalTok{(score\_reading(Measurement}\OperatorTok{.}\NormalTok{Valid(}\OperatorTok{{-}}\FloatTok{1.5}\NormalTok{), }\FloatTok{10.0}\NormalTok{) }\OperatorTok{==} \StringTok{"negative"}\NormalTok{)}
    \FunctionTok{assert}\NormalTok{(score\_reading(Measurement}\OperatorTok{.}\NormalTok{Valid(}\FloatTok{12.5}\NormalTok{), }\FloatTok{10.0}\NormalTok{) }\OperatorTok{==} \StringTok{"over by 2.5"}\NormalTok{)}
    \FunctionTok{assert}\NormalTok{(score\_reading(Measurement}\OperatorTok{.}\NormalTok{Valid(}\FloatTok{4.0}\NormalTok{), }\FloatTok{10.0}\NormalTok{) }\OperatorTok{==} \StringTok{"nominal"}\NormalTok{)}
    \FunctionTok{assert}\NormalTok{(score\_reading(Measurement}\OperatorTok{.}\NormalTok{Missing, }\FloatTok{10.0}\NormalTok{) }\OperatorTok{==} \StringTok{"missing"}\NormalTok{)}
\NormalTok{    println(}\StringTok{"match guards: ok"}\NormalTok{)}
    \ControlFlowTok{return} \DecValTok{0}
\NormalTok{\}}
\end{Highlighting}
\end{Shaded}

Notice the clean fallback \texttt{Measurement.Valid(\_)} after the two
guarded \texttt{Measurement.Valid(v)\ if\ ...} arms: guards do not
promise exhaustiveness by themselves, so always cover the unguarded case
or \texttt{\_}.

\begin{center}\rule{0.5\linewidth}{0.5pt}\end{center}

\section{6.5 Exhaustiveness Checking and Practical Pattern
Rules}\label{exhaustiveness-checking-and-practical-pattern-rules}

The compiler's \texttt{ExhaustivePass} validates \texttt{when}
expressions across \texttt{choice} variants, \texttt{bool},
\texttt{option\textless{}T\textgreater{}}, and
\texttt{result\textless{}T,\ E\textgreater{}} (reporting \texttt{E0210}
if a case is unhandled). Keep two practical pattern boundaries in mind:

\begin{enumerate}
\def\labelenumi{\arabic{enumi}.}
\tightlist
\item
  \textbf{No nested destructuring inside constructor patterns}:
  \texttt{when\ opt\ \{\ some((a,\ b))\ =\textgreater{}\ ...\ \}} is
  rejected; bind \texttt{some(pair)} cleanly and unpack
  \texttt{val\ (a,\ b)\ =\ pair} (or return a small \texttt{record} like
  \texttt{UiSize} instead of a nested tuple).
\item
  \textbf{All arms of a value-producing \texttt{when} must agree on
  type}: mismatched arm types trigger \texttt{E0208}.
\end{enumerate}

\subsection{Common Mistakes}\label{common-mistakes-5}

\begin{itemize}
\tightlist
\item
  \textbf{Nesting tuple patterns inside \texttt{some(...)} or
  \texttt{ok(...)}}: bind \texttt{some(v)} or \texttt{ok(rec)} as a
  single identifier and destructure inside the arm block.
\item
  \textbf{Using \texttt{match} instead of \texttt{when}}: \texttt{match}
  is reserved; IRIS uses \texttt{when} for pattern matching.
\item
  \textbf{Forgetting a fallback arm after guarded arms}: guarded
  patterns may all evaluate to \texttt{false} at runtime, so include an
  unguarded catch-all arm.
\end{itemize}

\subsection{Try It Yourself}\label{try-it-yourself-5}

\begin{enumerate}
\def\labelenumi{\arabic{enumi}.}
\tightlist
\item
  Write a
  \texttt{choice\ Expr\ \{\ Lit(i64),\ Add(i64,\ i64),\ Neg(i64)\ \}}
  evaluator using \texttt{when}.
\item
  Write \texttt{classify\_char\_code(code:\ i64)\ -\textgreater{}\ str}
  using inclusive range arms \texttt{48..=57}, \texttt{65..=90},
  \texttt{97..=122}, and \texttt{\_}.
\item
  Remove one \texttt{choice} variant from a \texttt{when} block
  temporarily and inspect the compiler's \texttt{E0210} error message
  via \texttt{iris\ check}.
\end{enumerate}

\begin{center}\rule{0.5\linewidth}{0.5pt}\end{center}

\chapter{Chapter 7: Generics, Const Generics, and
Traits}\label{chapter-7-generics-const-generics-and-traits}

\section{7.1 Generic Functions and Generic
Records}\label{generic-functions-and-generic-records}

IRIS monomorphizes generic functions and generic records
(\texttt{MonoPass}) into type-specialized SSA code, giving static type
safety with zero dynamic boxing overhead:

\begin{Shaded}
\begin{Highlighting}[]
\KeywordTok{record}\NormalTok{ PairBox[T, U] \{}
\NormalTok{    first: T,}
\NormalTok{    second: U,}
\NormalTok{\}}

\KeywordTok{def}\NormalTok{ identity[T](value: T) }\OperatorTok{{-}\textgreater{}}\NormalTok{ T \{}
    \ControlFlowTok{return}\NormalTok{ value}
\NormalTok{\}}

\KeywordTok{def}\NormalTok{ swap\_pair[A, B](p: PairBox}\OperatorTok{\textless{}}\NormalTok{A, B}\OperatorTok{\textgreater{}}\NormalTok{) }\OperatorTok{{-}\textgreater{}}\NormalTok{ PairBox}\OperatorTok{\textless{}}\NormalTok{B, A}\OperatorTok{\textgreater{}}\NormalTok{ \{}
    \ControlFlowTok{return}\NormalTok{ PairBox \{ first: p}\OperatorTok{.}\NormalTok{second, second: p}\OperatorTok{.}\NormalTok{first \}}
\NormalTok{\}}

\KeywordTok{def}\NormalTok{ main() }\OperatorTok{{-}\textgreater{}} \DataTypeTok{i64}\NormalTok{ \{}
    \KeywordTok{val}\NormalTok{ n }\OperatorTok{=}\NormalTok{ identity(}\DecValTok{42}\NormalTok{)}
    \KeywordTok{val}\NormalTok{ s }\OperatorTok{=}\NormalTok{ identity(}\StringTok{"IRIS"}\NormalTok{)}
    \KeywordTok{val}\NormalTok{ p: PairBox}\OperatorTok{\textless{}}\DataTypeTok{str}\NormalTok{, }\DataTypeTok{i64}\OperatorTok{\textgreater{}} \OperatorTok{=}\NormalTok{ PairBox \{ first: }\StringTok{"port"}\NormalTok{, second: }\DecValTok{8080}\NormalTok{ \}}
    \KeywordTok{val}\NormalTok{ swapped }\OperatorTok{=}\NormalTok{ swap\_pair(p)}

    \FunctionTok{assert}\NormalTok{(n }\OperatorTok{==} \DecValTok{42} \OperatorTok{\&\&}\NormalTok{ s }\OperatorTok{==} \StringTok{"IRIS"}\NormalTok{)}
    \FunctionTok{assert}\NormalTok{(swapped}\OperatorTok{.}\NormalTok{first }\OperatorTok{==} \DecValTok{8080} \OperatorTok{\&\&}\NormalTok{ swapped}\OperatorTok{.}\NormalTok{second }\OperatorTok{==} \StringTok{"port"}\NormalTok{)}
\NormalTok{    println(}\StringTok{"generics: ok"}\NormalTok{)}
    \ControlFlowTok{return} \DecValTok{0}
\NormalTok{\}}
\end{Highlighting}
\end{Shaded}

Observe the angle vs.~bracket convention across IRIS definitions and
call sites: declarations write \texttt{def\ f{[}T{]}} and
\texttt{record\ R{[}T{]}}, whereas type annotations write
\texttt{R\textless{}i64\textgreater{}} (just like
\texttt{list\textless{}T\textgreater{}} and
\texttt{option\textless{}T\textgreater{}}).

\begin{center}\rule{0.5\linewidth}{0.5pt}\end{center}

\section{7.2 Const Generics and Higher-Kinded
Records}\label{const-generics-and-higher-kinded-records}

Records and functions can also parameterize over \textbf{compile-time
integer sizes (\texttt{const\ N:\ usize})} and \textbf{higher-kinded
type constructors (\texttt{F{[}\_{]}})}:

\begin{Shaded}
\begin{Highlighting}[]
\KeywordTok{record}\NormalTok{ FixedVec[T, }\KeywordTok{const}\NormalTok{ N: }\DataTypeTok{usize}\NormalTok{] \{}
\NormalTok{    data: [T; N],}
\NormalTok{\}}

\KeywordTok{record}\NormalTok{ Wrapped[F[\_], A] \{}
\NormalTok{    inner: F}\OperatorTok{\textless{}}\NormalTok{A}\OperatorTok{\textgreater{}}\NormalTok{,}
\NormalTok{\}}

\KeywordTok{record}\NormalTok{ Box[T] \{}
\NormalTok{    item: T,}
\NormalTok{\}}

\KeywordTok{def}\NormalTok{ main() }\OperatorTok{{-}\textgreater{}} \DataTypeTok{i64}\NormalTok{ \{}
    \KeywordTok{val}\NormalTok{ vec3: FixedVec}\OperatorTok{\textless{}}\DataTypeTok{f64}\NormalTok{, }\DecValTok{3}\OperatorTok{\textgreater{}} \OperatorTok{=}\NormalTok{ FixedVec \{ data: [}\FloatTok{1.0}\NormalTok{, }\FloatTok{2.5}\NormalTok{, }\FloatTok{4.0}\NormalTok{] \}}
    \KeywordTok{val}\NormalTok{ wrapped: Wrapped}\OperatorTok{\textless{}}\NormalTok{Box, }\DataTypeTok{str}\OperatorTok{\textgreater{}} \OperatorTok{=}\NormalTok{ Wrapped \{}
\NormalTok{        inner: Box \{ item: }\StringTok{"tensor{-}ready"}\NormalTok{ \},}
\NormalTok{    \}}

    \FunctionTok{assert}\NormalTok{(vec3}\OperatorTok{.}\NormalTok{data[}\DecValTok{1}\NormalTok{] }\OperatorTok{==} \FloatTok{2.5}\NormalTok{)}
    \FunctionTok{assert}\NormalTok{(wrapped}\OperatorTok{.}\NormalTok{inner}\OperatorTok{.}\NormalTok{item }\OperatorTok{==} \StringTok{"tensor{-}ready"}\NormalTok{)}
\NormalTok{    println(}\StringTok{"const generics and higher{-}kinded records: ok"}\NormalTok{)}
    \ControlFlowTok{return} \DecValTok{0}
\NormalTok{\}}
\end{Highlighting}
\end{Shaded}

Const-generic array bounds propagate directly into field metadata
(\texttt{IrModule::struct\_defs}) across unannotated struct literals and
cross-function calls.

\begin{center}\rule{0.5\linewidth}{0.5pt}\end{center}

\section{7.3 Defining and Implementing
Traits}\label{defining-and-implementing-traits}

A \texttt{trait} declares shared behavior that any \texttt{record} can
implement via \texttt{impl\ Trait\ for\ Type}:

\begin{Shaded}
\begin{Highlighting}[]
\KeywordTok{trait}\NormalTok{ Metric \{}
    \KeywordTok{def}\NormalTok{ score(self) }\OperatorTok{{-}\textgreater{}} \DataTypeTok{f64}
\NormalTok{\}}

\KeywordTok{record}\NormalTok{ SensorNode \{}
\NormalTok{    id: }\DataTypeTok{i64}\NormalTok{,}
\NormalTok{    confidence: }\DataTypeTok{f64}\NormalTok{,}
\NormalTok{\}}

\KeywordTok{impl}\NormalTok{ Metric }\ControlFlowTok{for}\NormalTok{ SensorNode \{}
    \KeywordTok{def}\NormalTok{ score(self) }\OperatorTok{{-}\textgreater{}} \DataTypeTok{f64}\NormalTok{ \{}
        \ControlFlowTok{return}\NormalTok{ self}\OperatorTok{.}\NormalTok{confidence }\OperatorTok{*} \FloatTok{100.0}
\NormalTok{    \}}
\NormalTok{\}}

\KeywordTok{def}\NormalTok{ main() }\OperatorTok{{-}\textgreater{}} \DataTypeTok{i64}\NormalTok{ \{}
    \KeywordTok{val}\NormalTok{ node }\OperatorTok{=}\NormalTok{ SensorNode \{ id: }\DecValTok{1}\NormalTok{, confidence: }\FloatTok{0.94}\NormalTok{ \}}
    \FunctionTok{assert}\NormalTok{(node}\OperatorTok{.}\NormalTok{score() }\OperatorTok{==} \FloatTok{94.0}\NormalTok{)}
\NormalTok{    println(}\StringTok{"traits: ok"}\NormalTok{)}
    \ControlFlowTok{return} \DecValTok{0}
\NormalTok{\}}
\end{Highlighting}
\end{Shaded}

Inside \texttt{impl\ ...\ \{\ ...\ \}}, declare methods with
\texttt{def} (do not write \texttt{pub\ def}) and write bare
\texttt{self} (or \texttt{self:\ RecordName}) so method bodies resolve
record fields cleanly on the concrete implementing type rather than
abstract \texttt{Self}. Traits may also declare effect bounds on
signatures---for instance
\texttt{def\ forward(self,\ input:\ DiffTensor)\ -\textgreater{}\ DiffTensor\ effect\ alloc}.

\begin{center}\rule{0.5\linewidth}{0.5pt}\end{center}

\section{7.4 Trait Bounds and Dynamic
Dispatch}\label{trait-bounds-and-dynamic-dispatch}

IRIS supports \textbf{two complementary polymorphism mechanisms}:

\begin{enumerate}
\def\labelenumi{\arabic{enumi}.}
\tightlist
\item
  \textbf{Static dispatch via bounded generics
  (\texttt{{[}T\ where\ T:\ Trait{]}} / blanket \texttt{impl})}:
  monomorphized at compile time for maximum speed.
\item
  \textbf{Dynamic dispatch via trait objects (\texttt{dyn\ Trait})}:
  stores heterogeneous record values behind a uniform interface when
  types only unify at runtime (such as a
  \texttt{list\textless{}dyn\ Show\textgreater{}} or simulation entity
  list).
\end{enumerate}

\begin{Shaded}
\begin{Highlighting}[]
\KeywordTok{trait}\NormalTok{ Show \{}
    \KeywordTok{def}\NormalTok{ show(self) }\OperatorTok{{-}\textgreater{}} \DataTypeTok{str}
\NormalTok{\}}

\KeywordTok{trait}\NormalTok{ Describe \{}
    \KeywordTok{def}\NormalTok{ describe(self) }\OperatorTok{{-}\textgreater{}} \DataTypeTok{str}
\NormalTok{\}}

\KeywordTok{record}\NormalTok{ Thruster \{}
\NormalTok{    name: }\DataTypeTok{str}\NormalTok{,}
\NormalTok{\}}

\KeywordTok{impl}\NormalTok{ Show }\ControlFlowTok{for}\NormalTok{ Thruster \{}
    \KeywordTok{def}\NormalTok{ show(self) }\OperatorTok{{-}\textgreater{}} \DataTypeTok{str}\NormalTok{ \{}
        \ControlFlowTok{return}\NormalTok{ self}\OperatorTok{.}\NormalTok{name}
\NormalTok{    \}}
\NormalTok{\}}

\KeywordTok{impl}\NormalTok{[T }\KeywordTok{where}\NormalTok{ T: Show] Describe }\ControlFlowTok{for}\NormalTok{ T \{}
    \KeywordTok{def}\NormalTok{ describe(self) }\OperatorTok{{-}\textgreater{}} \DataTypeTok{str}\NormalTok{ \{}
        \ControlFlowTok{return} \StringTok{"unit: "} \OperatorTok{+}\NormalTok{ self}\OperatorTok{.}\NormalTok{show()}
\NormalTok{    \}}
\NormalTok{\}}

\KeywordTok{def}\NormalTok{ main() }\OperatorTok{{-}\textgreater{}} \DataTypeTok{i64}\NormalTok{ \{}
    \KeywordTok{val} \ConstantTok{unit} \OperatorTok{=}\NormalTok{ Thruster \{ name: }\StringTok{"RCS{-}Alpha"}\NormalTok{ \}}
    \FunctionTok{assert}\NormalTok{(}\ConstantTok{unit}\OperatorTok{.}\NormalTok{describe() }\OperatorTok{==} \StringTok{"unit: RCS{-}Alpha"}\NormalTok{)}

    \KeywordTok{val}\NormalTok{ obj: dyn Show }\OperatorTok{=} \ConstantTok{unit}
    \FunctionTok{assert}\NormalTok{(obj}\OperatorTok{.}\NormalTok{show() }\OperatorTok{==} \StringTok{"RCS{-}Alpha"}\NormalTok{)}
\NormalTok{    println(}\StringTok{"trait bounds and dyn dispatch: ok"}\NormalTok{)}
    \ControlFlowTok{return} \DecValTok{0}
\NormalTok{\}}
\end{Highlighting}
\end{Shaded}

\begin{center}\rule{0.5\linewidth}{0.5pt}\end{center}

\section{7.5 Macros and Compile-Time
Functions}\label{macros-and-compile-time-functions}

For lightweight syntactic abstraction and compile-time scalar
evaluation, IRIS provides \textbf{hygienic macros (\texttt{defmacro})}
and \textbf{compile-time functions (\texttt{const\ def})}:

\begin{Shaded}
\begin{Highlighting}[]
\NormalTok{defmacro double\_expr(x) }\OperatorTok{=\textgreater{}}\NormalTok{ x }\OperatorTok{*} \DecValTok{2}

\KeywordTok{const} \KeywordTok{def}\NormalTok{ square\_i64(n: }\DataTypeTok{i64}\NormalTok{) }\OperatorTok{{-}\textgreater{}} \DataTypeTok{i64}\NormalTok{ \{}
    \ControlFlowTok{return}\NormalTok{ n }\OperatorTok{*}\NormalTok{ n}
\NormalTok{\}}

\KeywordTok{def}\NormalTok{ main() }\OperatorTok{{-}\textgreater{}} \DataTypeTok{i64}\NormalTok{ \{}
    \KeywordTok{val}\NormalTok{ a }\OperatorTok{=}\NormalTok{ double\_expr}\OperatorTok{!}\NormalTok{(}\DecValTok{21}\NormalTok{)}
    \KeywordTok{val}\NormalTok{ b }\OperatorTok{=}\NormalTok{ square\_i64(}\DecValTok{9}\NormalTok{)}
    \FunctionTok{assert}\NormalTok{(a }\OperatorTok{==} \DecValTok{42}\NormalTok{)}
    \FunctionTok{assert}\NormalTok{(b }\OperatorTok{==} \DecValTok{81}\NormalTok{)}
\NormalTok{    println(}\StringTok{"macros and const def: ok"}\NormalTok{)}
    \ControlFlowTok{return} \DecValTok{0}
\NormalTok{\}}
\end{Highlighting}
\end{Shaded}

Macro invocations use \texttt{!} (\texttt{double\_expr!(21)}), while
\texttt{const\ def} functions evaluate deterministically during
compilation when supplied constant inputs.

\subsection{Common Mistakes}\label{common-mistakes-6}

\begin{itemize}
\tightlist
\item
  \textbf{Writing \texttt{pub\ def} inside an \texttt{impl} block}:
  \texttt{impl} methods use plain \texttt{def}.
\item
  \textbf{Omitting return type parameter annotations on empty generic
  constructors}: \texttt{val\ b\ =\ Box\ \{\ ...\ \}} works when fields
  determine \texttt{T}, whereas \texttt{set\_new()} needs
  \texttt{val\ s:\ Set\textless{}i64\textgreater{}\ =\ set\_new()}.
\item
  \textbf{Separating \texttt{impl\ Trait\ for\ R} into another module
  from \texttt{R}}: keep trait implementations in the defining module
  alongside \texttt{record\ R}.
\end{itemize}

\subsection{Try It Yourself}\label{try-it-yourself-6}

\begin{enumerate}
\def\labelenumi{\arabic{enumi}.}
\tightlist
\item
  Write a generic \texttt{record\ RingSlot{[}T,\ const\ N:\ usize{]}}
  holding \texttt{{[}T;\ N{]}} and instantiate it for \texttt{i64} and
  \texttt{str}.
\item
  Define a
  \texttt{trait\ CelestialMetric\ \{\ def\ kinetic\_energy(self:\ Self)\ -\textgreater{}\ f64\ \}},
  implement it for a \texttt{record\ Orbiter}, and query an instance
  through \texttt{dyn\ CelestialMetric}.
\item
  Write \texttt{defmacro\ clamp01(x)} or
  \texttt{const\ def\ cube(n:\ i64)\ -\textgreater{}\ i64} and assert
  their results natively.
\end{enumerate}

\begin{center}\rule{0.5\linewidth}{0.5pt}\end{center}

\chapter{Chapter 8: Error Handling, Ownership, and Memory
Management}\label{chapter-8-error-handling-ownership-and-memory-management}

\section{8.1 Options, Results, and Error
Propagation}\label{options-results-and-error-propagation}

IRIS represents absence and recoverable failure explicitly in the type
system:

\begin{itemize}
\tightlist
\item
  \textbf{\texttt{option\textless{}T\textgreater{}}} (\texttt{some(v)} /
  \texttt{none()}): when a value may legitimately be absent
  (\texttt{map\_get}, \texttt{parse\_i64}, \texttt{find}).
\item
  \textbf{\texttt{result\textless{}T,\ E\textgreater{}}} (\texttt{ok(v)}
  / \texttt{err(e)}): when an operation can fail with diagnostic
  context.
\end{itemize}

Use three ergonomic tools to compose them cleanly without deeply nested
boilerplate:

\begin{enumerate}
\def\labelenumi{\arabic{enumi}.}
\tightlist
\item
  \textbf{\texttt{?} early propagation}: inside a function returning
  \texttt{result\textless{}U,\ E\textgreater{}} or
  \texttt{option\textless{}U\textgreater{}}, suffixing an expression
  with \texttt{?} unwraps the success value or immediately returns the
  \texttt{err}/\texttt{none}.
\item
  \textbf{\texttt{??} default coalescing}: \texttt{expr\ ??\ fallback}
  yields the inner value or evaluates \texttt{fallback}.
\item
  \textbf{\texttt{try\ \{\ ...\ \}\ catch\ reason\ \{\ ...\ \}}}:
  intercepts \texttt{err(reason)} propagated via \texttt{?} locally
  inside a block!
\end{enumerate}

\begin{Shaded}
\begin{Highlighting}[]
\KeywordTok{def}\NormalTok{ safe\_div(a: }\DataTypeTok{i64}\NormalTok{, b: }\DataTypeTok{i64}\NormalTok{) }\OperatorTok{{-}\textgreater{}} \DataTypeTok{result}\OperatorTok{\textless{}}\DataTypeTok{i64}\NormalTok{, }\DataTypeTok{str}\OperatorTok{\textgreater{}}\NormalTok{ \{}
    \ControlFlowTok{if}\NormalTok{ b }\OperatorTok{==} \DecValTok{0}\NormalTok{ \{}
        \ControlFlowTok{return} \FunctionTok{err}\NormalTok{(}\StringTok{"division by zero"}\NormalTok{)}
\NormalTok{    \}}
    \ControlFlowTok{return} \FunctionTok{ok}\NormalTok{(a }\OperatorTok{/}\NormalTok{ b)}
\NormalTok{\}}

\KeywordTok{def}\NormalTok{ stepped(a: }\DataTypeTok{i64}\NormalTok{, b: }\DataTypeTok{i64}\NormalTok{) }\OperatorTok{{-}\textgreater{}} \DataTypeTok{result}\OperatorTok{\textless{}}\DataTypeTok{i64}\NormalTok{, }\DataTypeTok{str}\OperatorTok{\textgreater{}}\NormalTok{ \{}
    \KeywordTok{val}\NormalTok{ q }\OperatorTok{=}\NormalTok{ safe\_div(a, b)}\OperatorTok{?}
    \ControlFlowTok{return} \FunctionTok{ok}\NormalTok{(q }\OperatorTok{*} \DecValTok{10}\NormalTok{)}
\NormalTok{\}}

\KeywordTok{def}\NormalTok{ main() }\OperatorTok{{-}\textgreater{}} \DataTypeTok{i64}\NormalTok{ \{}
    \FunctionTok{assert}\NormalTok{(unwrap(stepped(}\DecValTok{24}\NormalTok{, }\DecValTok{4}\NormalTok{)) }\OperatorTok{==} \DecValTok{60}\NormalTok{)}

    \KeywordTok{val}\NormalTok{ recovered }\OperatorTok{=}\NormalTok{ try \{}
\NormalTok{        stepped(}\DecValTok{24}\NormalTok{, }\DecValTok{0}\NormalTok{)}\OperatorTok{?}
\NormalTok{    \} catch reason \{}
        \FunctionTok{assert}\NormalTok{(reason }\OperatorTok{==} \StringTok{"division by zero"}\NormalTok{)}
        \OperatorTok{{-}}\DecValTok{1}
\NormalTok{    \}}
    \FunctionTok{assert}\NormalTok{(recovered }\OperatorTok{==} \OperatorTok{{-}}\DecValTok{1}\NormalTok{)}

    \KeywordTok{val}\NormalTok{ port }\OperatorTok{=} \FunctionTok{parse\_i64}\NormalTok{(}\StringTok{"bad"}\NormalTok{) }\OperatorTok{??} \DecValTok{8080}
    \FunctionTok{assert}\NormalTok{(port }\OperatorTok{==} \DecValTok{8080}\NormalTok{)}
\NormalTok{    println(}\StringTok{"error handling: ok"}\NormalTok{)}
    \ControlFlowTok{return} \DecValTok{0}
\NormalTok{\}}
\end{Highlighting}
\end{Shaded}

\begin{center}\rule{0.5\linewidth}{0.5pt}\end{center}

\section{8.2 Assertions, Panics, and Runtime
Diagnostics}\label{assertions-panics-and-runtime-diagnostics}

Use \texttt{assert(condition)} for invariants and test checks, or
\texttt{panic("reason")} on unrecoverable states. When a runtime check
triggers (integer overflow, division by zero, out-of-bounds access, or
failed \texttt{assert}/\texttt{unwrap}), both backends exit with code
\textbf{101}, preserve stdout already flushed before the failure, and
format a precise source-mapped diagnostic:

{\def\LTcaptype{none} % do not increment counter
\begin{longtable}[]{@{}
  >{\raggedright\arraybackslash}p{(\linewidth - 2\tabcolsep) * \real{0.5000}}
  >{\raggedright\arraybackslash}p{(\linewidth - 2\tabcolsep) * \real{0.5000}}@{}}
\toprule\noalign{}
\begin{minipage}[b]{\linewidth}\raggedright
Code
\end{minipage} & \begin{minipage}[b]{\linewidth}\raggedright
Meaning
\end{minipage} \\
\midrule\noalign{}
\endhead
\bottomrule\noalign{}
\endlastfoot
\texttt{E0402} & Integer division or remainder by zero \\
\texttt{E0403} & Index out of bounds (\texttt{list}, array,
\texttt{str}, or tensor rank) \\
\texttt{E0406} & \texttt{panic(...)}, failed \texttt{assert(...)}, or
\texttt{unwrap} on \texttt{none}/\texttt{err} \\
\texttt{E0408} & Stack overflow / recursion depth limit exceeded \\
\texttt{E0409} & Checked integer arithmetic overflow/underflow \\
\end{longtable}
}

Run \texttt{iris\ explain\ \textless{}code\textgreater{}} (for example
\texttt{iris\ explain\ E0409}) to inspect detailed causes and fixes at
any time.

\begin{center}\rule{0.5\linewidth}{0.5pt}\end{center}

\section{8.3 Borrowing, Mutable Write-Back, and Explicit
Moves}\label{borrowing-mutable-write-back-and-explicit-moves}

IRIS supports block- and CFG-checked \textbf{references (\texttt{\&T},
\texttt{\&mut\ T})}, \textbf{dereference write-back
(\texttt{*p\ =\ val}, \texttt{*p\ +=\ k})}, and \textbf{explicit
\texttt{move}} transfers:

\begin{Shaded}
\begin{Highlighting}[]
\KeywordTok{def}\NormalTok{ read\_plus\_one(r: }\OperatorTok{\&}\DataTypeTok{i64}\NormalTok{) }\OperatorTok{{-}\textgreater{}} \DataTypeTok{i64}\NormalTok{ \{}
    \ControlFlowTok{return} \OperatorTok{*}\NormalTok{r }\OperatorTok{+} \DecValTok{1}
\NormalTok{\}}

\KeywordTok{def}\NormalTok{ bump\_in\_place(p: }\OperatorTok{\&}\NormalTok{mut }\DataTypeTok{i64}\NormalTok{, delta: }\DataTypeTok{i64}\NormalTok{) }\OperatorTok{{-}\textgreater{}} \DataTypeTok{i64}\NormalTok{ \{}
    \OperatorTok{*}\NormalTok{p }\OperatorTok{+=}\NormalTok{ delta}
    \ControlFlowTok{return} \OperatorTok{*}\NormalTok{p}
\NormalTok{\}}

\KeywordTok{def}\NormalTok{ main() }\OperatorTok{{-}\textgreater{}} \DataTypeTok{i64}\NormalTok{ \{}
    \KeywordTok{val}\NormalTok{ base }\OperatorTok{=} \DecValTok{41}
    \KeywordTok{val}\NormalTok{ ref\_base: }\OperatorTok{\&}\DataTypeTok{i64} \OperatorTok{=} \OperatorTok{\&}\NormalTok{base}
    \FunctionTok{assert}\NormalTok{(read\_plus\_one(ref\_base) }\OperatorTok{==} \DecValTok{42}\NormalTok{)}

    \KeywordTok{var}\NormalTok{ counter: }\DataTypeTok{i64} \OperatorTok{=} \DecValTok{10}
    \KeywordTok{val}\NormalTok{ p: }\OperatorTok{\&}\NormalTok{mut }\DataTypeTok{i64} \OperatorTok{=} \OperatorTok{\&}\NormalTok{mut counter}
    \FunctionTok{assert}\NormalTok{(bump\_in\_place(p, }\DecValTok{5}\NormalTok{) }\OperatorTok{==} \DecValTok{15}\NormalTok{)}
    \FunctionTok{assert}\NormalTok{(counter }\OperatorTok{==} \DecValTok{15}\NormalTok{)}

    \CommentTok{// Explicit move invalidates the original binding at compile time (E0382 on reuse)}
    \KeywordTok{val}\NormalTok{ items: }\DataTypeTok{list}\OperatorTok{\textless{}}\DataTypeTok{i64}\OperatorTok{\textgreater{}} \OperatorTok{=} \DataTypeTok{list}\NormalTok{(}\DecValTok{3}\NormalTok{, }\DecValTok{5}\NormalTok{, }\DecValTok{8}\NormalTok{)}
    \KeywordTok{val}\NormalTok{ owned }\OperatorTok{=}\NormalTok{ move items}
    \FunctionTok{assert}\NormalTok{(}\FunctionTok{list\_get}\NormalTok{(owned, }\DecValTok{2}\NormalTok{) }\OperatorTok{==} \DecValTok{8}\NormalTok{)}

\NormalTok{    println(}\StringTok{"borrowing and moves: ok"}\NormalTok{)}
    \ControlFlowTok{return} \DecValTok{0}
\NormalTok{\}}
\end{Highlighting}
\end{Shaded}

Dereference write-back supports conditional \texttt{\&mut} branch
selection
(\texttt{val\ p\ =\ if\ cond\ \{\ \&mut\ a\ \}\ else\ \{\ \&mut\ b\ \};\ *p\ =\ 77}),
loop mutations, compound operators (\texttt{+=}, \texttt{-=},
\texttt{*=}), and projected struct/tuple targets
(\texttt{(*p).field\ =\ v}, \texttt{(*p).0\ =\ v}).

\begin{center}\rule{0.5\linewidth}{0.5pt}\end{center}

\section{8.4 CFG Non-Lexical Lifetimes and Explicit Drop
Scopes}\label{cfg-non-lexical-lifetimes-and-explicit-drop-scopes}

Unlike naïve lexical borrow scopes, IRIS runs
\textbf{\texttt{CfgBorrowCheckerPass}} over control-flow-graph liveness
(NLL) alongside \textbf{\texttt{DropElaborationPass}}:

\begin{itemize}
\tightlist
\item
  \textbf{Fine-grained path tracking}: move state tracks field/tuple
  projections (\texttt{move\ pair.a} invalidates \texttt{pair.a},
  leaving \texttt{pair.b} valid while rejecting moving all of
  \texttt{pair}).
\item
  \textbf{Scope-driven drops}: \texttt{EnterScope} / \texttt{ExitScope}
  / \texttt{Drop} instructions ensure destructors run cleanly once per
  owned value---and skip values already moved out so a moved binding is
  never double-freed.
\item
  \textbf{Opt-in strict record moves}: setting
  \texttt{IRIS\_STRICT\_RECORD\_OWNERSHIP=1} turns by-value record
  parameter passing, record-field initialization from another record
  binding, and binding reassignment into implicit ownership moves
  (\texttt{E0382} on subsequent use of the moved source). Explicit
  \texttt{move} is always checked regardless of this environment
  variable.
\end{itemize}

\begin{center}\rule{0.5\linewidth}{0.5pt}\end{center}

\section{8.5 Two-Tier Memory Reclamation and Frame
Pools}\label{two-tier-memory-reclamation-and-frame-pools}

Long-running services and interactive 60 FPS desktop GUIs need
predictable, leak-free memory behavior. IRIS combines \textbf{two
reclamation tiers} (see \href{memory-management.md}{Memory Management}):

\begin{enumerate}
\def\labelenumi{\arabic{enumi}.}
\item
  \textbf{Tier 1 --- Compiler-Inserted Reference Counting + Cycle
  Collector}: Strings, lists, maps, options, results, channels, mutexes,
  atomics, closures, tensors, and native records
  (\texttt{IRIS\_RC\_NATIVE\_OBJECT}) are reference-counted.
  \texttt{GcAnnotatePass} inserts \texttt{retain}/\texttt{release} calls
  at \textbf{last use}, freeing short-lived values promptly, while a
  Bacon--Rajan cycle collector reclaims reference cycles. \emph{(Note:
  because \texttt{dyn\ Trait} wrappers hold raw pointers to wrapped
  records, records are reclaimed via the RC table's native-object
  destructor and cycle collector rather than raw last-use pointer
  free.)}
\item
  \textbf{Tier 2 --- Explicit Arena Frame Pools
  (\texttt{ui\_frame\_pool})}: Interactive UI frames construct fresh
  node trees, styles, hit-test rectangles, and draw command lists every
  frame. Wrapping each frame inside \texttt{ui\_frame\_pool\_begin()}
  \ldots{} \texttt{ui\_frame\_pool\_destroy(pool)} allocates transient
  data on a fast per-frame arena and reclaims every byte in bulk at
  frame end. \textbf{\texttt{FramePoolScopePass}} additionally verifies
  that any open \texttt{ui\_frame\_pool\_begin()} has
  \texttt{iris\_ui\_pool\_destroy} inserted before \texttt{Return} or
  \texttt{Panic} terminators so pools cannot leak on early exits.
\end{enumerate}

\begin{quote}
\textbf{Critical Rule for Persistent State Inside a Frame Pool}: Record
field writes and struct literals allocate on whichever allocator is
currently active! Always \textbf{pause the active frame pool},
\textbf{detach surviving values (\texttt{ui\_state\_detach},
\texttt{list}/\texttt{str} clones)} onto the main heap, mutate
persistent app state while paused, and only \textbf{resume and destroy}
afterward:

\begin{Shaded}
\begin{Highlighting}[]
\CommentTok{// illustrative UI frame loop pattern (requires \textasciigrave{}bring std.ui\textasciigrave{})}
\KeywordTok{val}\NormalTok{ pool }\OperatorTok{=}\NormalTok{ ui\_frame\_pool\_begin()}
\KeywordTok{val}\NormalTok{ tree }\OperatorTok{=}\NormalTok{ build\_ui(app, metrics)}
\KeywordTok{val}\NormalTok{ buf }\OperatorTok{=}\NormalTok{ ui\_paint(tree, ui\_st, buf)}

\CommentTok{// Pause pool BEFORE detaching or mutating persistent \textasciigrave{}app\textasciigrave{} state!}
\KeywordTok{val}\NormalTok{ paused }\OperatorTok{=}\NormalTok{ ui\_frame\_pool\_pause()}
\KeywordTok{val}\NormalTok{ clean\_ui\_st }\OperatorTok{=}\NormalTok{ ui\_state\_detach(next\_ui\_st)}
\NormalTok{app }\OperatorTok{=}\NormalTok{ apply\_events(app, clean\_ui\_st, events)}
\KeywordTok{val}\NormalTok{ \_resumed }\OperatorTok{=}\NormalTok{ ui\_frame\_pool\_resume(paused)}
\KeywordTok{val}\NormalTok{ \_freed }\OperatorTok{=}\NormalTok{ ui\_frame\_pool\_destroy(pool)}
\end{Highlighting}
\end{Shaded}
\end{quote}

\begin{center}\rule{0.5\linewidth}{0.5pt}\end{center}

\section{8.6 Inspecting Allocations with GC
Statistics}\label{inspecting-allocations-with-gc-statistics}

You can audit memory health and hunt leaks directly from IRIS code via
\texttt{std.async}:

\begin{Shaded}
\begin{Highlighting}[]
\KeywordTok{bring}\NormalTok{ std}\OperatorTok{.}\KeywordTok{async}

\KeywordTok{def}\NormalTok{ main() }\OperatorTok{{-}\textgreater{}} \DataTypeTok{i64}\NormalTok{ \{}
    \KeywordTok{val}\NormalTok{ before }\OperatorTok{=}\NormalTok{ gc\_get\_stats()}
    \KeywordTok{var}\NormalTok{ sink }\OperatorTok{=} \DecValTok{0}
    \ControlFlowTok{for}\NormalTok{ i }\ControlFlowTok{in} \DecValTok{0}\OperatorTok{..}\DecValTok{100}\NormalTok{ \{}
        \KeywordTok{val}\NormalTok{ tmp }\OperatorTok{=} \StringTok{f"packet{-}}\ControlFlowTok{\{}\NormalTok{i}\ControlFlowTok{\}}\StringTok{"}
\NormalTok{        sink }\OperatorTok{=}\NormalTok{ sink }\OperatorTok{+} \FunctionTok{len}\NormalTok{(tmp)}
\NormalTok{    \}}
\NormalTok{    gc\_collect\_sweep()}
    \KeywordTok{val}\NormalTok{ after }\OperatorTok{=}\NormalTok{ gc\_get\_stats()}

    \KeywordTok{val}\NormalTok{ allocs }\OperatorTok{=} \FunctionTok{map\_get}\NormalTok{(after, }\StringTok{"allocated"}\NormalTok{) }\OperatorTok{??} \DecValTok{0}
    \KeywordTok{val}\NormalTok{ freed }\OperatorTok{=} \FunctionTok{map\_get}\NormalTok{(after, }\StringTok{"freed"}\NormalTok{) }\OperatorTok{??} \DecValTok{0}
    \KeywordTok{val}\NormalTok{ live\_bytes }\OperatorTok{=} \FunctionTok{map\_get}\NormalTok{(after, }\StringTok{"raw\_bytes\_live"}\NormalTok{) }\OperatorTok{??} \DecValTok{0}
    \KeywordTok{val}\NormalTok{ peak\_pool }\OperatorTok{=} \FunctionTok{map\_get}\NormalTok{(after, }\StringTok{"pool\_bytes\_peak"}\NormalTok{) }\OperatorTok{??} \DecValTok{0}

    \FunctionTok{assert}\NormalTok{(sink }\OperatorTok{\textgreater{}} \DecValTok{0}\NormalTok{)}
    \FunctionTok{assert}\NormalTok{(allocs }\OperatorTok{\textgreater{}=} \DecValTok{0} \OperatorTok{\&\&}\NormalTok{ freed }\OperatorTok{\textgreater{}=} \DecValTok{0} \OperatorTok{\&\&}\NormalTok{ live\_bytes }\OperatorTok{\textgreater{}=} \DecValTok{0} \OperatorTok{\&\&}\NormalTok{ peak\_pool }\OperatorTok{\textgreater{}=} \DecValTok{0}\NormalTok{)}
\NormalTok{    println(}\StringTok{f"gc check ok: alloc=}\ControlFlowTok{\{}\NormalTok{allocs}\ControlFlowTok{\}}\StringTok{ freed=}\ControlFlowTok{\{}\NormalTok{freed}\ControlFlowTok{\}}\StringTok{ live\_bytes=}\ControlFlowTok{\{}\NormalTok{live\_bytes}\ControlFlowTok{\}}\StringTok{"}\NormalTok{)}
    \ControlFlowTok{return} \DecValTok{0}
\NormalTok{\}}
\end{Highlighting}
\end{Shaded}

\texttt{gc\_get\_stats()} reports \texttt{allocated}, \texttt{freed},
\texttt{cycles}, \texttt{weak\_invalidated}, raw heap metrics
(\texttt{raw\_allocs}, \texttt{raw\_frees}, \texttt{raw\_bytes\_live},
\texttt{raw\_bytes\_peak}), frame-pool counters
(\texttt{pool\_bytes\_live}, \texttt{pool\_bytes\_peak}), and per-site
counts (\texttt{site:\textless{}location\textgreater{}}). In a
steady-state server or desktop app, \texttt{raw\_bytes\_live} should
plateau cleanly rather than climbing monotonically.

\subsection{Common Mistakes}\label{common-mistakes-7}

\begin{itemize}
\tightlist
\item
  \textbf{Unwrapping \texttt{option}/\texttt{result} inside production
  logic with \texttt{unwrap(...)}}: prefer \texttt{?}, \texttt{??}, or
  exhaustive \texttt{when}.
\item
  \textbf{Mutating persistent application records while
  \texttt{ui\_frame\_pool\_begin()} is active}: if any child string,
  list, or nested record allocated in the pool ends up stored in
  long-lived state, destroying the pool triggers a use-after-free crash.
  Always \texttt{ui\_frame\_pool\_pause()} +
  \texttt{ui\_state\_detach(...)} first.
\item
  \textbf{Reusing a binding after \texttt{move\ x} or partial field move
  of the whole record}: moving \texttt{p.a} leaves \texttt{p.b} usable,
  but \texttt{move\ p} afterward is rejected by
  \texttt{CfgBorrowCheckerPass}.
\end{itemize}

\subsection{Try It Yourself}\label{try-it-yourself-7}

\begin{enumerate}
\def\labelenumi{\arabic{enumi}.}
\tightlist
\item
  Write a chained config parser using \texttt{parse\_i64} + \texttt{?} +
  \texttt{try\ /\ catch} and assert both valid and malformed cases.
\item
  Write a function with two \texttt{var} counters, select \texttt{\&mut}
  between them conditionally via \texttt{if}, mutate via
  \texttt{*ptr\ +=\ 10}, and assert both final values.
\item
  Run a 1,000-iteration allocation loop alongside
  \texttt{gc\_get\_stats()} before and after
  \texttt{gc\_collect\_sweep()} to observe \texttt{allocated}
  vs.~\texttt{freed}.
\end{enumerate}

\begin{center}\rule{0.5\linewidth}{0.5pt}\end{center}

\part{Part III: Effects, Concurrency, and
Systems}\label{part-iii-effects-concurrency-and-systems}

\chapter{Chapter 9: Static Effects, Handlers, and
Speculation}\label{chapter-9-static-effects-handlers-and-speculation}

\section{9.1 Effect Annotations and Strict Effect
Checking}\label{effect-annotations-and-strict-effect-checking}

In safety-critical software, reading a type signature like
\texttt{-\textgreater{}\ i64} does not tell you whether a function
touches the filesystem, allocates on the heap, invokes FFI, or sleeps on
the network. IRIS makes side effects explicit directly in function and
closure types via \texttt{effect} clauses:

\begin{Shaded}
\begin{Highlighting}[]
\CommentTok{/// Pass \textasciigrave{}{-}{-}strict{-}effects\textasciigrave{} to enforce effect bounds statically.}
\KeywordTok{def}\NormalTok{ compute\_pure(x: }\DataTypeTok{i64}\NormalTok{) }\OperatorTok{{-}\textgreater{}} \DataTypeTok{i64}\NormalTok{ \{}
    \ControlFlowTok{return}\NormalTok{ x }\OperatorTok{*} \DecValTok{2} \OperatorTok{+} \DecValTok{1}
\NormalTok{\}}

\KeywordTok{def}\NormalTok{ log\_result(value: }\DataTypeTok{i64}\NormalTok{) }\OperatorTok{{-}\textgreater{}} \DataTypeTok{i64}\NormalTok{ effect io, alloc \{}
\NormalTok{    println(}\StringTok{f"result = }\ControlFlowTok{\{}\NormalTok{value}\ControlFlowTok{\}}\StringTok{"}\NormalTok{)}
    \ControlFlowTok{return}\NormalTok{ value}
\NormalTok{\}}

\KeywordTok{def}\NormalTok{ main() }\OperatorTok{{-}\textgreater{}} \DataTypeTok{i64}\NormalTok{ effect io, alloc \{}
    \KeywordTok{val}\NormalTok{ ans }\OperatorTok{=}\NormalTok{ compute\_pure(}\DecValTok{20}\NormalTok{)}
    \ControlFlowTok{return}\NormalTok{ log\_result(ans) }\OperatorTok{{-}} \DecValTok{41}
\NormalTok{\}}
\end{Highlighting}
\end{Shaded}

Common tracked effect labels in IRIS and its standard library include
\texttt{io}, \texttt{alloc}, \texttt{fs}, \texttt{net}, \texttt{ffi},
\texttt{thread}, \texttt{time}, \texttt{env}, \texttt{sys},
\texttt{random}, \texttt{throw}, \texttt{reflect}, and
\texttt{transaction}.

By default, effect annotations document intent while migration stays
frictionless. Enable \textbf{\texttt{-\/-strict-effects}}
(\texttt{iris\ check\ -\/-strict-effects} or
\texttt{iris\ build\ -\/-strict-effects}) to enforce that every
function's \texttt{effect} clause covers every effect triggered by its
body and callees (\texttt{E0304} on violations). Both
\texttt{iris\ embedded} and \texttt{iris\ evolve} turn on strict effect
checking automatically.

\begin{center}\rule{0.5\linewidth}{0.5pt}\end{center}

\section{9.2 Polymorphic Effect Rows}\label{polymorphic-effect-rows}

Higher-order helpers rarely know in advance which effects a caller's
callback might trigger. Use an uppercase effect variable such as
\texttt{effect\ E} to stay agnostic cleanly:

\begin{Shaded}
\begin{Highlighting}[]
\CommentTok{/// Run with: iris {-}{-}strict{-}effects run effects.iris}
\KeywordTok{def}\NormalTok{ apply(value: }\DataTypeTok{i64}\NormalTok{, operation: }\OperatorTok{|}\DataTypeTok{i64}\OperatorTok{|} \OperatorTok{{-}\textgreater{}} \DataTypeTok{i64}\NormalTok{ effect E) }\OperatorTok{{-}\textgreater{}} \DataTypeTok{i64}\NormalTok{ effect E \{}
    \ControlFlowTok{return}\NormalTok{ operation(value)}
\NormalTok{\}}

\KeywordTok{def}\NormalTok{ increment(value: }\DataTypeTok{i64}\NormalTok{) }\OperatorTok{{-}\textgreater{}} \DataTypeTok{i64}\NormalTok{ \{}
    \ControlFlowTok{return}\NormalTok{ value }\OperatorTok{+} \DecValTok{1}
\NormalTok{\}}

\KeywordTok{def}\NormalTok{ main() }\OperatorTok{{-}\textgreater{}} \DataTypeTok{i64}\NormalTok{ effect io, alloc, throw \{}
    \KeywordTok{val}\NormalTok{ plus\_one }\OperatorTok{=} \OperatorTok{|}\NormalTok{value: }\DataTypeTok{i64}\OperatorTok{|}\NormalTok{ value }\OperatorTok{+} \DecValTok{1}
    \FunctionTok{assert}\NormalTok{(apply(}\DecValTok{41}\NormalTok{, plus\_one) }\OperatorTok{==} \DecValTok{42}\NormalTok{)}
    \FunctionTok{assert}\NormalTok{(apply(}\DecValTok{41}\NormalTok{, increment) }\OperatorTok{==} \DecValTok{42}\NormalTok{)}
\NormalTok{    println(}\StringTok{"polymorphic effects: ok"}\NormalTok{)}
    \ControlFlowTok{return} \DecValTok{0}
\NormalTok{\}}
\end{Highlighting}
\end{Shaded}

Because \texttt{apply} quantifies over \texttt{effect\ E}, passing a
pure function (\texttt{increment}) requires zero caller effects, whereas
passing an effectful callback propagates that callback's exact effects
to the caller under \texttt{-\/-strict-effects}.

\begin{center}\rule{0.5\linewidth}{0.5pt}\end{center}

\section{9.3 Algebraic Effect Handlers}\label{algebraic-effect-handlers}

Beyond checking effects, IRIS supports scoped interception through
\textbf{\texttt{handle\ \{\ ...\ \}\ with\ \{\ ...\ \}}} expressions. An
external declaration (\texttt{extern\ def}) or effectful operation can
be intercepted cleanly inside a test harness or sandbox block without
rewriting production callers:

\begin{Shaded}
\begin{Highlighting}[]
\KeywordTok{extern} \KeywordTok{def}\NormalTok{ echo(value: }\DataTypeTok{str}\NormalTok{) }\OperatorTok{{-}\textgreater{}} \DataTypeTok{str}

\KeywordTok{def}\NormalTok{ main() }\OperatorTok{{-}\textgreater{}} \DataTypeTok{i64}\NormalTok{ \{}
    \KeywordTok{val}\NormalTok{ intercepted }\OperatorTok{=}\NormalTok{ handle \{}
\NormalTok{        echo(}\StringTok{"hello"}\NormalTok{)}
\NormalTok{    \} with \{}
\NormalTok{        echo(value) }\OperatorTok{=\textgreater{}} \StringTok{"handled"}\NormalTok{,}
\NormalTok{    \}}

    \FunctionTok{assert}\NormalTok{(intercepted }\OperatorTok{==} \StringTok{"handled"}\NormalTok{)}
\NormalTok{    println(}\StringTok{"effect handlers: ok"}\NormalTok{)}
    \ControlFlowTok{return} \DecValTok{0}
\NormalTok{\}}
\end{Highlighting}
\end{Shaded}

Notice \texttt{handle\ ...\ with} evaluates directly as an expression,
so you can return its result, bind it to a \texttt{val}, or assign it to
state.

\begin{center}\rule{0.5\linewidth}{0.5pt}\end{center}

\section{9.4 Transactional Speculation and
Rollback}\label{transactional-speculation-and-rollback}

When testing an untrusted optimization, heuristic policy, or multi-step
mutation, \texttt{std.speculation} gives you \textbf{speculative
evaluation and reversible thread-local transactions}:

\begin{Shaded}
\begin{Highlighting}[]
\KeywordTok{bring}\NormalTok{ std}\OperatorTok{.}\NormalTok{speculation}

\KeywordTok{def}\NormalTok{ main() }\OperatorTok{{-}\textgreater{}} \DataTypeTok{i64}\NormalTok{ \{}
    \KeywordTok{val}\NormalTok{ candidate }\OperatorTok{=} \OperatorTok{||} \DecValTok{500}
    \KeywordTok{val}\NormalTok{ fallback }\OperatorTok{=} \OperatorTok{||} \DecValTok{10}
    \KeywordTok{val}\NormalTok{ in\_bounds }\OperatorTok{=} \OperatorTok{|}\NormalTok{value: }\DataTypeTok{i64}\OperatorTok{|}\NormalTok{ value }\OperatorTok{\textgreater{}=} \DecValTok{0} \OperatorTok{\&\&}\NormalTok{ value }\OperatorTok{\textless{}} \DecValTok{100}

    \KeywordTok{val}\NormalTok{ chosen }\OperatorTok{=}\NormalTok{ speculate\_with\_fallback(candidate, fallback, in\_bounds)}
    \FunctionTok{assert}\NormalTok{(chosen }\OperatorTok{==} \DecValTok{10}\NormalTok{)}

    \KeywordTok{val}\NormalTok{ verified: }\DataTypeTok{option}\OperatorTok{\textless{}}\DataTypeTok{i64}\OperatorTok{\textgreater{}} \OperatorTok{=}\NormalTok{ speculate(}\OperatorTok{||} \DecValTok{42}\NormalTok{, in\_bounds)}
    \FunctionTok{assert}\NormalTok{(unwrap(verified) }\OperatorTok{==} \DecValTok{42}\NormalTok{)}

    \CommentTok{// Transactional speculation snapshots lists, maps, atomics, and staged file writes:}
    \KeywordTok{val}\NormalTok{ xs: }\DataTypeTok{list}\OperatorTok{\textless{}}\DataTypeTok{i64}\OperatorTok{\textgreater{}} \OperatorTok{=} \DataTypeTok{list}\NormalTok{(}\DecValTok{1}\NormalTok{, }\DecValTok{2}\NormalTok{)}
    \KeywordTok{val}\NormalTok{ outcome }\OperatorTok{=}\NormalTok{ speculate\_transactional(}
        \OperatorTok{|}\NormalTok{n: }\DataTypeTok{i64}\OperatorTok{|}\NormalTok{ n }\OperatorTok{\textless{}} \DecValTok{0}\NormalTok{,}
        \OperatorTok{||}\NormalTok{ \{}
\NormalTok{            list\_push(xs, }\DecValTok{999}\NormalTok{)}
            \DecValTok{999}
\NormalTok{        \},}
\NormalTok{    )}
    \FunctionTok{assert}\NormalTok{(}\OperatorTok{!}\FunctionTok{is\_some}\NormalTok{(outcome))}
    \FunctionTok{assert}\NormalTok{(}\FunctionTok{list\_len}\NormalTok{(xs) }\OperatorTok{==} \DecValTok{2}\NormalTok{)}

\NormalTok{    println(}\StringTok{"speculation: ok"}\NormalTok{)}
    \ControlFlowTok{return} \DecValTok{0}
\NormalTok{\}}
\end{Highlighting}
\end{Shaded}

Under \texttt{speculate\_transactional}, \texttt{transaction\_begin()},
\texttt{transaction\_commit()}, and \texttt{transaction\_rollback()}
maintain nestable thread-local checkpoints over in-memory lists, maps,
atomics, and staged whole-file writes (\texttt{file\_write\_all} staged
through a synced two-phase \texttt{.new}/\texttt{.bak} WAL until
commit). External network services, open file streams, and physical
hardware actuators are excluded from the transactional callback effect
row and are intentionally not claimed reversible.

\begin{center}\rule{0.5\linewidth}{0.5pt}\end{center}

\section{9.5 Security Policies and
Sandboxing}\label{security-policies-and-sandboxing}

Defense in depth requires runtime mediation alongside static checks:

\begin{itemize}
\tightlist
\item
  \textbf{\texttt{-\/-sandbox} / \texttt{IRIS\_SECURITY\_POLICY}}:
  propagates an IRIS security policy into compiled LLVM binaries and the
  C runtime (\texttt{iris\_sandbox\_set\_policy\_spec}), mediating
  C-runtime filesystem access, SQLite opens, ML model loaders, sockets
  (TCP/UDP/HTTP/TLS/async), processes, environment mutation, Python FFI,
  and dynamic shared library loads (enforcing path traversal checks,
  allowlists/blocklists, byte limits, and open-handle quotas).
  Unmediated \texttt{extern\ def} declarations fail closed unless
  explicitly allowed via \texttt{allow\_ffi} / \texttt{ffi\_allowlist}.
\item
  \textbf{OS Worker Sandbox}: because arbitrary pre-linked third-party
  native machine code outside IRIS C-runtime mediation cannot be policed
  solely in-process, \texttt{iris\ evolve\ -\/-validate\ sandbox} (and
  host sandboxing APIs) isolates candidate binaries inside OS-native
  kernel sandboxes: Windows AppContainer/LPAC Job objects
  (CPU/memory/process/pipe caps), Linux Landlock (\texttt{5.13+}) +
  \texttt{seccomp-bpf} + \texttt{rlimit}s, and macOS Seatbelt
  \texttt{sandbox-exec} + \texttt{rlimit}s. Unsupported platforms fail
  closed.
\end{itemize}

\subsection{Common Mistakes}\label{common-mistakes-8}

\begin{itemize}
\tightlist
\item
  \textbf{Declaring \texttt{extern\ def} without an effect clause when
  \texttt{-\/-strict-effects} is enabled}: mark pure foreign
  declarations with \texttt{effect\ pure}, or list \texttt{effect\ ffi}
  / \texttt{effect\ io}.
\item
  \textbf{Expecting transactional rollback to undo an HTTP request or
  GPIO write}: only lists, maps, atomics, and staged whole-file writes
  participate in \texttt{std.speculation} transactions.
\end{itemize}

\subsection{Try It Yourself}\label{try-it-yourself-8}

\begin{enumerate}
\def\labelenumi{\arabic{enumi}.}
\tightlist
\item
  Create a small program mixing pure math with \texttt{println} and
  check it with \texttt{iris\ check\ -\/-strict-effects}. Observe
  \texttt{E0304} when an \texttt{effect\ io,\ alloc} clause is omitted,
  then fix the signature.
\item
  Wrap a mocked
  \texttt{extern\ def\ fetch\_sensor(id:\ i64)\ -\textgreater{}\ i64}
  inside
  \texttt{handle\ \{\ ...\ \}\ with\ \{\ fetch\_sensor(id)\ =\textgreater{}\ id\ *\ 10,\ \}}
  and run both interpreter and native builds.
\item
  Run a test file that attempts filesystem access using
  \texttt{iris\ -\/-sandbox\ -\/-emit\ eval\ ...} to see runtime policy
  denial (\texttt{E0407}).
\end{enumerate}

\begin{center}\rule{0.5\linewidth}{0.5pt}\end{center}

\chapter{Chapter 10: Concurrency, Async/Await, and
Parallelism}\label{chapter-10-concurrency-asyncawait-and-parallelism}

\section{10.1 Typed Channels and Spawned
Tasks}\label{typed-channels-and-spawned-tasks}

IRIS schedules \texttt{spawn} tasks onto a bounded native worker
executor (capped at 32 workers) and coordinates typed message passing
across tasks via \textbf{\texttt{chan\textless{}T\textgreater{}}}:

\begin{Shaded}
\begin{Highlighting}[]
\KeywordTok{def}\NormalTok{ main() }\OperatorTok{{-}\textgreater{}} \DataTypeTok{i64}\NormalTok{ \{}
    \KeywordTok{val}\NormalTok{ results: }\DataTypeTok{chan}\OperatorTok{\textless{}}\DataTypeTok{i64}\OperatorTok{\textgreater{}} \OperatorTok{=} \FunctionTok{channel}\NormalTok{()}

    \KeywordTok{spawn}\NormalTok{ \{}
        \FunctionTok{send}\NormalTok{(results, }\DecValTok{21} \OperatorTok{*} \DecValTok{2}\NormalTok{)}
\NormalTok{    \}}

    \KeywordTok{val}\NormalTok{ answer }\OperatorTok{=} \FunctionTok{recv}\NormalTok{(results)}
    \FunctionTok{assert}\NormalTok{(answer }\OperatorTok{==} \DecValTok{42}\NormalTok{)}
\NormalTok{    println(}\StringTok{"channels and spawn: ok"}\NormalTok{)}
    \ControlFlowTok{return} \DecValTok{0}
\NormalTok{\}}
\end{Highlighting}
\end{Shaded}

Waiters blocked on an empty channel help execute queued work rather than
spinning idle OS threads. For bounded queues, timeouts, and non-blocking
polling, import \texttt{std.async} for
\texttt{channel\_bounded(capacity)}, \texttt{try\_recv(ch)},
\texttt{recv\_timeout\_ms(ch,\ ms)}, \texttt{select\_first(channels)},
and \texttt{channel\_len(ch)}.

\begin{center}\rule{0.5\linewidth}{0.5pt}\end{center}

\section{10.2 Structured Concurrency with Task
Groups}\label{structured-concurrency-with-task-groups}

Detached \texttt{spawn} blocks are easy to lose track of. When a fan-out
computation must complete together---or support cooperative cancellation
on deadlines---group tasks explicitly with \texttt{new\_task\_group()}
from \texttt{std.async}:

\begin{Shaded}
\begin{Highlighting}[]
\KeywordTok{bring}\NormalTok{ std}\OperatorTok{.}\KeywordTok{async}

\KeywordTok{def}\NormalTok{ main() }\OperatorTok{{-}\textgreater{}} \DataTypeTok{i64}\NormalTok{ \{}
    \KeywordTok{val}\NormalTok{ completed }\OperatorTok{=} \DataTypeTok{atomic}\NormalTok{(}\DecValTok{0}\NormalTok{)}
    \KeywordTok{val}\NormalTok{ workers }\OperatorTok{=}\NormalTok{ new\_task\_group()}

    \ControlFlowTok{for}\NormalTok{ \_i }\ControlFlowTok{in} \DecValTok{0}\OperatorTok{..}\DecValTok{32}\NormalTok{ \{}
        \KeywordTok{spawn}\NormalTok{(workers) \{}
            \ControlFlowTok{if} \OperatorTok{!}\NormalTok{cancellation\_requested() \{}
                \FunctionTok{atomic\_add}\NormalTok{(completed, }\DecValTok{1}\NormalTok{)}
\NormalTok{            \}}
\NormalTok{        \}}
\NormalTok{    \}}

\NormalTok{    task\_group\_join(workers)}
    \FunctionTok{assert}\NormalTok{(}\FunctionTok{atomic\_load}\NormalTok{(completed) }\OperatorTok{==} \DecValTok{32}\NormalTok{)}

    \KeywordTok{val}\NormalTok{ cancelled }\OperatorTok{=}\NormalTok{ new\_task\_group()}
\NormalTok{    task\_group\_cancel(cancelled)}
    \KeywordTok{spawn}\NormalTok{(cancelled) \{}
        \FunctionTok{atomic\_add}\NormalTok{(completed, }\DecValTok{1000}\NormalTok{)}
\NormalTok{    \}}
\NormalTok{    task\_group\_join(cancelled)}

    \FunctionTok{assert}\NormalTok{(task\_group\_cancelled(cancelled))}
    \FunctionTok{assert}\NormalTok{(}\FunctionTok{atomic\_load}\NormalTok{(completed) }\OperatorTok{==} \DecValTok{32}\NormalTok{)}
\NormalTok{    println(}\StringTok{"task groups: ok"}\NormalTok{)}
    \ControlFlowTok{return} \DecValTok{0}
\NormalTok{\}}
\end{Highlighting}
\end{Shaded}

\texttt{task\_group\_join(group)} closes the group to new work and
blocks while helping drain queued tasks. Cancellation is
\textbf{cooperative}: queued tasks skip execution on entry when
cancelled, while long-running CPU loops should poll
\texttt{cancellation\_requested()} or
\texttt{task\_group\_cancelled(group)} at bounded intervals.

\begin{center}\rule{0.5\linewidth}{0.5pt}\end{center}

\section{10.3 Async/Await on the Cooperative
Executor}\label{asyncawait-on-the-cooperative-executor}

Mark asynchronous functions with \texttt{async\ def}. Calling an
\texttt{async\ def} returns a single-result awaitable
(\texttt{chan\textless{}T\textgreater{}} under the current RC ABI) that
completes cleanly with \texttt{await}:

\begin{Shaded}
\begin{Highlighting}[]
\KeywordTok{async} \KeywordTok{def}\NormalTok{ square\_async(value: }\DataTypeTok{i64}\NormalTok{) }\OperatorTok{{-}\textgreater{}} \DataTypeTok{i64}\NormalTok{ \{}
    \ControlFlowTok{return}\NormalTok{ value }\OperatorTok{*}\NormalTok{ value}
\NormalTok{\}}

\KeywordTok{def}\NormalTok{ main() }\OperatorTok{{-}\textgreater{}} \DataTypeTok{i64}\NormalTok{ \{}
    \KeywordTok{val}\NormalTok{ a }\OperatorTok{=}\NormalTok{ square\_async(}\DecValTok{3}\NormalTok{)}
    \KeywordTok{val}\NormalTok{ b }\OperatorTok{=}\NormalTok{ square\_async(}\DecValTok{4}\NormalTok{)}
    \KeywordTok{val}\NormalTok{ total }\OperatorTok{=} \KeywordTok{await}\NormalTok{ a }\OperatorTok{+} \KeywordTok{await}\NormalTok{ b}
    \FunctionTok{assert}\NormalTok{(total }\OperatorTok{==} \DecValTok{25}\NormalTok{)}
\NormalTok{    println(}\StringTok{"async await: ok"}\NormalTok{)}
    \ControlFlowTok{return} \DecValTok{0}
\NormalTok{\}}
\end{Highlighting}
\end{Shaded}

Under the compiler's \texttt{AsyncLowerPass} (the final pass in the
standard pipeline, running \emph{after} \texttt{GcAnnotatePass}),
\texttt{async\ def} bodies lower cleanly into pollable state machines on
both native LLVM builds and the interpreter.

\begin{quote}
\textbf{Async Runtime Scope}: The cooperative executor supports
non-blocking wakeups for timers, channels, \texttt{tcp\_accept\_async},
and \texttt{tcp\_read\_async}. Synchronous socket connects, writes, and
\texttt{std.http} requests remain blocking calls---offload or pool them
explicitly in high-concurrency servers.
\end{quote}

\begin{center}\rule{0.5\linewidth}{0.5pt}\end{center}

\section{10.4 Data Parallelism with Parallel For Loops, SIMD, and
GPU}\label{data-parallelism-with-parallel-for-loops-simd-and-gpu}

When a computation maps independently across an integer index range, use
\texttt{par\ for} to unlock thread, vector, or device parallelism from a
single construct:

\begin{Shaded}
\begin{Highlighting}[]
\KeywordTok{def}\NormalTok{ main() }\OperatorTok{{-}\textgreater{}} \DataTypeTok{i64}\NormalTok{ \{}
    \KeywordTok{val}\NormalTok{ total }\OperatorTok{=} \DataTypeTok{atomic}\NormalTok{(}\DecValTok{0}\NormalTok{)}
    \KeywordTok{par} \ControlFlowTok{for}\NormalTok{ index }\ControlFlowTok{in} \DecValTok{1}\OperatorTok{..}\DecValTok{101}\NormalTok{ \{}
        \FunctionTok{atomic\_add}\NormalTok{(total, index)}
\NormalTok{    \}}
    \FunctionTok{assert}\NormalTok{(}\FunctionTok{atomic\_load}\NormalTok{(total) }\OperatorTok{==} \DecValTok{5050}\NormalTok{)}
\NormalTok{    println(}\StringTok{"par for: ok"}\NormalTok{)}
    \ControlFlowTok{return} \DecValTok{0}
\NormalTok{\}}
\end{Highlighting}
\end{Shaded}

Notice how \texttt{par\ for} accepts flag-selected acceleration without
changing source logic:

\begin{itemize}
\tightlist
\item
  \textbf{Default native (\texttt{iris\ build} / \texttt{iris\ run})}:
  executes iterations concurrently across worker threads with closure
  captures marshaled safely.
\item
  \textbf{\texttt{-\/-simd} / \texttt{-\/-emit\ simd}}: lowers
  element-wise \texttt{par\ for} loops into explicit vector kernels
  (AVX2 on x86-64, NEON on AArch64) and enables LLVM's loop vectorizer
  across loops.
\item
  \textbf{\texttt{-\/-gpu\ cuda} / \texttt{-\/-emit\ cuda} /
  \texttt{-\/-emit\ cuda-ptx}}: emits NVIDIA PTX device kernels for
  numeric and tensor \texttt{par\ for} loops. \emph{(Verification
  boundary: CUDA PTX emission and CPU fallback simulation pass automated
  tests, while physical NVIDIA GPU kernel execution has not yet been
  tested on GPU hardware.)}
\end{itemize}

\begin{center}\rule{0.5\linewidth}{0.5pt}\end{center}

\section{10.5 Lock-Free Atomics and
Mutexes}\label{lock-free-atomics-and-mutexes}

For lightweight counters shared across \texttt{spawn} or
\texttt{par\ for}, prefer lock-free atomics:

\begin{Shaded}
\begin{Highlighting}[]
\KeywordTok{def}\NormalTok{ main() }\OperatorTok{{-}\textgreater{}} \DataTypeTok{i64}\NormalTok{ \{}
    \KeywordTok{val}\NormalTok{ acc }\OperatorTok{=} \DataTypeTok{atomic}\NormalTok{(}\DecValTok{10}\NormalTok{)}
    \FunctionTok{atomic\_add}\NormalTok{(acc, }\DecValTok{32}\NormalTok{)}
    \FunctionTok{assert}\NormalTok{(}\FunctionTok{atomic\_load}\NormalTok{(acc) }\OperatorTok{==} \DecValTok{42}\NormalTok{)}
    \FunctionTok{atomic\_store}\NormalTok{(acc, }\DecValTok{7}\NormalTok{)}
    \FunctionTok{assert}\NormalTok{(}\FunctionTok{atomic\_load}\NormalTok{(acc) }\OperatorTok{==} \DecValTok{7}\NormalTok{)}
\NormalTok{    println(}\StringTok{"atomics: ok"}\NormalTok{)}
    \ControlFlowTok{return} \DecValTok{0}
\NormalTok{\}}
\end{Highlighting}
\end{Shaded}

When coordinating shared multi-field state or serialized resource access
across threads, use \texttt{mutex\textless{}T\textgreater{}}
(\texttt{mutex}, \texttt{mutex\_lock}, \texttt{mutex\_unlock}) or
communicate exclusively over typed
\texttt{chan\textless{}T\textgreater{}} queues.

\subsection{Common Mistakes}\label{common-mistakes-9}

\begin{itemize}
\tightlist
\item
  \textbf{Calling \texttt{recv(ch)} on the main thread when nobody calls
  \texttt{send(ch,\ ...)}}: deadlocks execution until interrupted.
  Ensure a producer runs inside \texttt{spawn} or check non-blockingly
  with \texttt{std.async.try\_recv}.
\item
  \textbf{Sharing unsynchronized mutable lists across parallel workers}:
  protect shared state behind a \texttt{chan\textless{}T\textgreater{}},
  \texttt{atomic}, or \texttt{mutex\textless{}T\textgreater{}}.
\item
  \textbf{Expecting pre-emptive thread termination on
  \texttt{task\_group\_cancel}}: long-running compute tasks must check
  \texttt{cancellation\_requested()}.
\end{itemize}

\subsection{Try It Yourself}\label{try-it-yourself-9}

\begin{enumerate}
\def\labelenumi{\arabic{enumi}.}
\tightlist
\item
  Build a three-worker fan-in pipeline using a shared
  \texttt{chan\textless{}i64\textgreater{}} and a \texttt{task\_group}
  that waits for all workers to emit their partial sums.
\item
  Benchmark a numeric array loop under default \texttt{iris\ build}
  vs.~\texttt{iris\ build\ -\/-simd}.
\item
  Write an \texttt{async\ def\ fetch\_pair(x:\ i64,\ y:\ i64)} helper
  that launches two \texttt{async\ def} tasks concurrently before
  \texttt{await}-ing both results.
\end{enumerate}

\begin{center}\rule{0.5\linewidth}{0.5pt}\end{center}

\chapter{Chapter 11: Files, Data Formats, SQLite, and
Networking}\label{chapter-11-files-data-formats-sqlite-and-networking}

\section{11.1 Filesystem, Paths, and Operating System
Queries}\label{filesystem-paths-and-operating-system-queries}

IRIS separates checked file I/O (\texttt{std.fs}), portable path
manipulation (\texttt{std.path}), and environment/process metadata
(\texttt{std.os}):

\begin{Shaded}
\begin{Highlighting}[]
\KeywordTok{bring}\NormalTok{ std}\OperatorTok{.}\NormalTok{path}
\KeywordTok{bring}\NormalTok{ std}\OperatorTok{.}\NormalTok{os}

\KeywordTok{def}\NormalTok{ main() }\OperatorTok{{-}\textgreater{}} \DataTypeTok{i64}\NormalTok{ \{}
    \KeywordTok{val}\NormalTok{ joined }\OperatorTok{=}\NormalTok{ normalize\_path(join\_path(}\StringTok{"data/logs"}\NormalTok{, }\StringTok{"../metrics.json"}\NormalTok{))}
    \FunctionTok{assert}\NormalTok{(basename(joined) }\OperatorTok{==} \StringTok{"metrics.json"}\NormalTok{)}
    \FunctionTok{assert}\NormalTok{(extension(joined) }\OperatorTok{==} \StringTok{"json"}\NormalTok{)}
    \FunctionTok{assert}\NormalTok{(stem(joined) }\OperatorTok{==} \StringTok{"metrics"}\NormalTok{)}

    \KeywordTok{val}\NormalTok{ cpus }\OperatorTok{=}\NormalTok{ cpu\_count()}
    \KeywordTok{val} \FunctionTok{pid} \OperatorTok{=}\NormalTok{ get\_pid()}
    \FunctionTok{assert}\NormalTok{(cpus }\OperatorTok{\textgreater{}=} \DecValTok{1} \OperatorTok{\&\&} \FunctionTok{pid} \OperatorTok{\textgreater{}} \DecValTok{0}\NormalTok{)}
\NormalTok{    println(}\StringTok{f"os check: cpus=}\ControlFlowTok{\{}\NormalTok{cpus}\ControlFlowTok{\}}\StringTok{ pid=}\ControlFlowTok{\{}\FunctionTok{pid}\ControlFlowTok{\}}\StringTok{ file=}\ControlFlowTok{\{}\NormalTok{joined}\ControlFlowTok{\}}\StringTok{"}\NormalTok{)}
    \ControlFlowTok{return} \DecValTok{0}
\NormalTok{\}}
\end{Highlighting}
\end{Shaded}

In \texttt{std.fs}, high-level helpers (\texttt{read\_text},
\texttt{write\_text}, \texttt{append\_text}, \texttt{read\_lines},
\texttt{copy\_file}) and buffered streams (\texttt{file\_open\_stream},
\texttt{file\_read\_stream}, \texttt{file\_write\_stream},
\texttt{file\_close\_stream}) return
\texttt{result\textless{}...,\ str\textgreater{}} so filesystem
permission errors or missing paths never fail silently.

\begin{center}\rule{0.5\linewidth}{0.5pt}\end{center}

\section{11.2 JSON, CSV, Tables, and
DataFrames}\label{json-csv-tables-and-dataframes}

Structured data formats live in \texttt{std.json}, \texttt{std.csv},
\texttt{std.table}, and \texttt{std.dataframe}:

\begin{Shaded}
\begin{Highlighting}[]
\KeywordTok{bring}\NormalTok{ std}\OperatorTok{.}\NormalTok{json}
\KeywordTok{bring}\NormalTok{ std}\OperatorTok{.}\NormalTok{csv}
\KeywordTok{bring}\NormalTok{ std}\OperatorTok{.}\NormalTok{dataframe}

\KeywordTok{def}\NormalTok{ main() }\OperatorTok{{-}\textgreater{}} \DataTypeTok{i64}\NormalTok{ \{}
    \CommentTok{// RFC{-}4180 CSV rows}
    \KeywordTok{val}\NormalTok{ row }\OperatorTok{=}\NormalTok{ csv\_parse\_row(}\StringTok{"telemetry,42,ok"}\NormalTok{)}
    \FunctionTok{assert}\NormalTok{(}\FunctionTok{list\_len}\NormalTok{(row) }\OperatorTok{==} \DecValTok{3} \OperatorTok{\&\&} \FunctionTok{list\_get}\NormalTok{(row, }\DecValTok{1}\NormalTok{) }\OperatorTok{==} \StringTok{"42"}\NormalTok{)}

    \CommentTok{// Structured JSON encoding \& path query}
    \KeywordTok{val}\NormalTok{ doc }\OperatorTok{=}\NormalTok{ json\_obj(}
        \DataTypeTok{list}\NormalTok{(}\StringTok{"service"}\NormalTok{, }\StringTok{"replicas"}\NormalTok{, }\StringTok{"healthy"}\NormalTok{),}
        \DataTypeTok{list}\NormalTok{(json\_str(}\StringTok{"iris{-}node"}\NormalTok{), json\_num\_i64(}\DecValTok{3}\NormalTok{), json\_bool(}\ConstantTok{true}\NormalTok{)),}
\NormalTok{    )}
    \FunctionTok{assert}\NormalTok{(json\_query(doc, }\StringTok{"service"}\NormalTok{) }\OperatorTok{==} \StringTok{"iris{-}node"}\NormalTok{)}
    \FunctionTok{assert}\NormalTok{(json\_query(doc, }\StringTok{"replicas"}\NormalTok{) }\OperatorTok{==} \StringTok{"3"}\NormalTok{)}

    \CommentTok{// Typed numeric DataFrame}
    \KeywordTok{var}\NormalTok{ df }\OperatorTok{=}\NormalTok{ df\_new(}\DataTypeTok{list}\NormalTok{(}\StringTok{"step"}\NormalTok{, }\StringTok{"loss"}\NormalTok{))}
\NormalTok{    df }\OperatorTok{=}\NormalTok{ df\_add\_row(df, }\DataTypeTok{list}\NormalTok{(}\FloatTok{1.0}\NormalTok{, }\FloatTok{0.85}\NormalTok{))}
\NormalTok{    df }\OperatorTok{=}\NormalTok{ df\_add\_row(df, }\DataTypeTok{list}\NormalTok{(}\FloatTok{2.0}\NormalTok{, }\FloatTok{0.22}\NormalTok{))}
    \KeywordTok{val}\NormalTok{ filtered }\OperatorTok{=}\NormalTok{ df\_filter(df, }\StringTok{"loss"}\NormalTok{, }\FloatTok{0.0}\NormalTok{, }\FloatTok{0.5}\NormalTok{)}
    \FunctionTok{assert}\NormalTok{(df\_len(filtered) }\OperatorTok{==} \DecValTok{1}\NormalTok{)}

\NormalTok{    println(}\StringTok{"json, csv, and dataframe: ok"}\NormalTok{)}
    \ControlFlowTok{return} \DecValTok{0}
\NormalTok{\}}
\end{Highlighting}
\end{Shaded}

Use \texttt{json\_escape(text)} before manually assembling JSON
fragments from arbitrary user strings, and remember that
\texttt{std.table} stores string columns (\texttt{Table},
\texttt{table\_new}, \texttt{table\_insert}, \texttt{table\_where},
\texttt{table\_to\_csv}) whereas \texttt{std.dataframe} stores
\texttt{f64} numeric matrices optimized for ML/scientific filtering and
analytics.

\begin{center}\rule{0.5\linewidth}{0.5pt}\end{center}

\section{11.3 Embedded SQLite and Parameterized
Queries}\label{embedded-sqlite-and-parameterized-queries}

IRIS bundles SQLite directly (\texttt{std.sql}), requiring zero external
database driver installation:

\begin{Shaded}
\begin{Highlighting}[]
\KeywordTok{bring}\NormalTok{ std}\OperatorTok{.}\NormalTok{sql}

\KeywordTok{def}\NormalTok{ main() }\OperatorTok{{-}\textgreater{}} \DataTypeTok{i64}\NormalTok{ \{}
    \KeywordTok{val}\NormalTok{ db }\OperatorTok{=}\NormalTok{ sql\_open(}\StringTok{":memory:"}\NormalTok{)}
    \FunctionTok{assert}\NormalTok{(db }\OperatorTok{!=} \DecValTok{0}\NormalTok{)}
    \FunctionTok{assert}\NormalTok{(sql\_exec(db, }\StringTok{"CREATE TABLE tasks (title TEXT, done INTEGER)"}\NormalTok{) }\OperatorTok{\textgreater{}=} \DecValTok{0}\NormalTok{)}
    \FunctionTok{assert}\NormalTok{(sql\_exec\_params(db, }\StringTok{"INSERT INTO tasks VALUES (?, ?)"}\NormalTok{, }\DataTypeTok{list}\NormalTok{(}\StringTok{"build UI"}\NormalTok{, }\StringTok{"1"}\NormalTok{)) }\OperatorTok{\textgreater{}=} \DecValTok{0}\NormalTok{)}
    \FunctionTok{assert}\NormalTok{(sql\_exec\_params(db, }\StringTok{"INSERT INTO tasks VALUES (?, ?)"}\NormalTok{, }\DataTypeTok{list}\NormalTok{(}\StringTok{"write docs"}\NormalTok{, }\StringTok{"0"}\NormalTok{)) }\OperatorTok{\textgreater{}=} \DecValTok{0}\NormalTok{)}

    \KeywordTok{val}\NormalTok{ pending }\OperatorTok{=}\NormalTok{ sql\_query\_params(db, }\StringTok{"SELECT title FROM tasks WHERE done = ?"}\NormalTok{, }\DataTypeTok{list}\NormalTok{(}\StringTok{"0"}\NormalTok{))}
    \FunctionTok{assert}\NormalTok{(}\FunctionTok{list\_len}\NormalTok{(pending) }\OperatorTok{==} \DecValTok{1}\NormalTok{)}
    \FunctionTok{assert}\NormalTok{(}\FunctionTok{list\_get}\NormalTok{(}\FunctionTok{list\_get}\NormalTok{(pending, }\DecValTok{0}\NormalTok{), }\DecValTok{0}\NormalTok{) }\OperatorTok{==} \StringTok{"write docs"}\NormalTok{)}

\NormalTok{    sql\_close(db)}
\NormalTok{    println(}\StringTok{"sqlite: ok"}\NormalTok{)}
    \ControlFlowTok{return} \DecValTok{0}
\NormalTok{\}}
\end{Highlighting}
\end{Shaded}

Always keep values out of raw SQL strings using
\texttt{sql\_exec\_params} and \texttt{sql\_query\_params}. Parameter
binding avoids SQL injection cleanly and prevents quote-escaping bugs.

\begin{center}\rule{0.5\linewidth}{0.5pt}\end{center}

\section{11.4 Key-Value Storage, Cryptography, and
Visualization}\label{key-value-storage-cryptography-and-visualization}

For persistent flat configs, cryptographic checksums, HMAC
authentication, and crisp charting exports, IRIS includes:

\begin{itemize}
\tightlist
\item
  \textbf{\texttt{std.kv}}: simple file-backed key-value store
  (\texttt{kv\_set}, \texttt{kv\_get}, \texttt{kv\_delete},
  \texttt{kv\_keys}).
\item
  \textbf{\texttt{std.crypto}}: \texttt{hash\_sha256},
  \texttt{hmac\_simple}, \texttt{hmac\_verify},
  \texttt{timing\_safe\_equal}, \texttt{to\_hex}/\texttt{from\_hex},
  \texttt{to\_base64}/\texttt{from\_base64}, \texttt{generate\_uuid},
  \texttt{websocket\_accept\_key}, and \texttt{aes128\_encrypt\_cbc} /
  \texttt{aes128\_decrypt\_cbc}.
\item
  \textbf{\texttt{std.svg} and \texttt{std.termplot}}: programmatic SVG
  vector export (\texttt{svg\_new}, \texttt{svg\_rect},
  \texttt{svg\_circle}, \texttt{svg\_line}, \texttt{svg\_text},
  \texttt{svg\_save}) and instant ASCII terminal bar plots
  (\texttt{term\_plot\_bar}).
\end{itemize}

\begin{Shaded}
\begin{Highlighting}[]
\KeywordTok{bring}\NormalTok{ std}\OperatorTok{.}\NormalTok{crypto}

\KeywordTok{def}\NormalTok{ main() }\OperatorTok{{-}\textgreater{}} \DataTypeTok{i64}\NormalTok{ \{}
    \KeywordTok{val}\NormalTok{ digest }\OperatorTok{=}\NormalTok{ hash\_sha256(}\StringTok{"iris{-}1.0"}\NormalTok{)}
    \FunctionTok{assert}\NormalTok{(}\FunctionTok{len}\NormalTok{(digest) }\OperatorTok{==} \DecValTok{64}\NormalTok{)}

    \KeywordTok{val}\NormalTok{ mac }\OperatorTok{=}\NormalTok{ hmac\_simple(}\StringTok{"secret{-}key"}\NormalTok{, }\StringTok{"payload"}\NormalTok{)}
    \FunctionTok{assert}\NormalTok{(hmac\_verify(}\StringTok{"secret{-}key"}\NormalTok{, }\StringTok{"payload"}\NormalTok{, mac))}
    \FunctionTok{assert}\NormalTok{(timing\_safe\_equal(mac, mac))}
\NormalTok{    println(}\StringTok{"crypto: ok"}\NormalTok{)}
    \ControlFlowTok{return} \DecValTok{0}
\NormalTok{\}}
\end{Highlighting}
\end{Shaded}

\begin{center}\rule{0.5\linewidth}{0.5pt}\end{center}

\section{11.5 TCP/UDP, HTTP Client/Server, TLS, and Packet
Inspection}\label{tcpudp-http-clientserver-tls-and-packet-inspection}

IRIS ships a complete networking and protocol analysis suite:

\begin{itemize}
\item
  \textbf{\texttt{std.net}}: bounded, timeout-checked TCP/UDP sockets
  (\texttt{tcp\_connect\_checked}, \texttt{tcp\_listen},
  \texttt{tcp\_accept\_checked}, \texttt{tcp\_read\_bounded},
  \texttt{tcp\_read\_with\_timeout}, \texttt{tcp\_write\_all\_checked},
  \texttt{tcp\_read\_bytes}, \texttt{tcp\_write\_bytes},
  \texttt{udp\_open}, \texttt{udp\_send}, \texttt{udp\_recv},
  \texttt{WebSocketFrame}).
\item
  \textbf{\texttt{std.http} \& \texttt{std.http\_server}}:
  \texttt{HttpRequest}, \texttt{HttpResponse}, \texttt{http\_send},
  \texttt{http\_tls\_available()}, \texttt{http\_serve},
  \texttt{http\_serve\_mt}, routing helpers (\texttt{router\_new},
  \texttt{router\_add}, \texttt{router\_match}), and response builders
  (\texttt{json\_ok}, \texttt{text\_ok}, \texttt{bad\_request}).
  \emph{(Platform note: HTTPS via \texttt{http\_send} uses OS-native
  WinHTTP on Windows and fails closed when
  \texttt{http\_tls\_available()} is \texttt{false} on other native
  targets.)}
\item
  \textbf{\texttt{std.tls} \& \texttt{std.pcap}}: binary packet parsing
  (\texttt{pcap\_parse\_global\_header}, \texttt{pcap\_read\_packet})
  and passive TLS ClientHello inspection (\texttt{tls\_parse\_record},
  \texttt{tls\_parse\_client\_hello}, \texttt{tls\_calculate\_ja3},
  \texttt{tls\_calculate\_ja4}).
\end{itemize}

Testing failure boundaries locally is zero-dependency:

\begin{Shaded}
\begin{Highlighting}[]
\KeywordTok{bring}\NormalTok{ std}\OperatorTok{.}\NormalTok{net}

\KeywordTok{def}\NormalTok{ main() }\OperatorTok{{-}\textgreater{}} \DataTypeTok{i64}\NormalTok{ \{}
    \KeywordTok{val}\NormalTok{ invalid }\OperatorTok{=}\NormalTok{ TcpStream \{ handle: }\OperatorTok{{-}}\DecValTok{1}\NormalTok{ \}}
    \FunctionTok{assert}\NormalTok{(}\FunctionTok{is\_err}\NormalTok{(tcp\_read\_bounded(invalid, }\DecValTok{0}\NormalTok{)))}
    \FunctionTok{assert}\NormalTok{(}\FunctionTok{is\_err}\NormalTok{(tcp\_connect\_checked(}\StringTok{"127.0.0.1"}\NormalTok{, }\OperatorTok{{-}}\DecValTok{1}\NormalTok{, }\DecValTok{100}\NormalTok{, }\DecValTok{100}\NormalTok{)))}
\NormalTok{    println(}\StringTok{"checked networking: ok"}\NormalTok{)}
    \ControlFlowTok{return} \DecValTok{0}
\NormalTok{\}}
\end{Highlighting}
\end{Shaded}

\subsection{Common Mistakes}\label{common-mistakes-10}

\begin{itemize}
\tightlist
\item
  \textbf{Interpolating strings into \texttt{sql\_exec} directly}:
  always pass user or dynamic values as parameters to
  \texttt{sql\_exec\_params} / \texttt{sql\_query\_params}.
\item
  \textbf{Comparing signatures with \texttt{==} instead of
  \texttt{timing\_safe\_equal}}: use constant-time
  \texttt{timing\_safe\_equal} / \texttt{hmac\_verify} for MACs and
  secrets.
\item
  \textbf{Unbounded socket reads}: always prefer
  \texttt{tcp\_read\_bounded} / \texttt{tcp\_read\_with\_timeout} so
  malformed peers cannot exhaust memory.
\end{itemize}

\subsection{Try It Yourself}\label{try-it-yourself-10}

\begin{enumerate}
\def\labelenumi{\arabic{enumi}.}
\tightlist
\item
  Write a short program using \texttt{std.sql} that stores five sensor
  readings in \texttt{:memory:}, filters readings above a threshold with
  \texttt{sql\_query\_params}, and exports matching rows as CSV via
  \texttt{std.csv.csv\_emit\_row}.
\item
  Hash a string with \texttt{std.crypto.hash\_sha256}, verify its hex
  length (\texttt{64}), and test both matching and mismatched
  \texttt{hmac\_verify} keys.
\item
  Build a 20-line chart with \texttt{std.svg} (\texttt{svg\_new}
  -\textgreater{} \texttt{svg\_circle} -\textgreater{}
  \texttt{svg\_save}) and open the generated \texttt{.svg} file.
\end{enumerate}

\begin{center}\rule{0.5\linewidth}{0.5pt}\end{center}

\chapter{Chapter 12: Native Compilation, FFI, and Bare-Metal
Embedded}\label{chapter-12-native-compilation-ffi-and-bare-metal-embedded}

\section{12.1 Inside the Compiler Pipeline and SSA
IR}\label{inside-the-compiler-pipeline-and-ssa-ir}

IRIS compiles source text into a block-parameter SSA IR
(\texttt{IrModule}, inspired by MLIR---using basic block parameters
rather than phi nodes) and transforms it through a \textbf{20-pass
standard pipeline} (\texttt{build\_standard\_pipeline()} in
\texttt{src/pass/mod.rs}):

{\def\LTcaptype{none} % do not increment counter
\begin{longtable}[]{@{}
  >{\raggedright\arraybackslash}p{(\linewidth - 4\tabcolsep) * \real{0.3333}}
  >{\raggedright\arraybackslash}p{(\linewidth - 4\tabcolsep) * \real{0.3333}}
  >{\raggedright\arraybackslash}p{(\linewidth - 4\tabcolsep) * \real{0.3333}}@{}}
\toprule\noalign{}
\begin{minipage}[b]{\linewidth}\raggedright
Order
\end{minipage} & \begin{minipage}[b]{\linewidth}\raggedright
Pass (\texttt{-\/-dump-ir-after} name)
\end{minipage} & \begin{minipage}[b]{\linewidth}\raggedright
Responsibility
\end{minipage} \\
\midrule\noalign{}
\endhead
\bottomrule\noalign{}
\endlastfoot
1 & \texttt{ValidatePass} (\texttt{validate}) & Verifies structural SSA
invariants and terminator consistency \\
2 & \texttt{TypeInferPass} (\texttt{type-infer}) & Resolves remaining
local types across blocks and expressions \\
3 & \texttt{MonoPass} (\texttt{mono}) & Monomorphizes generic functions
and const-generic records \\
4 & \texttt{CfgBorrowCheckerPass} (\texttt{cfg\_borrow\_checker}) & Runs
CFG non-lexical lifetime borrow, move, and projection checks \\
5 & \texttt{DropElaborationPass} (\texttt{drop\_elaboration}) & Emits
explicit \texttt{EnterScope}, \texttt{ExitScope}, and \texttt{Drop}
instructions \\
6 & \texttt{ShapeCheckPass} (\texttt{shape-check}) & Verifies tensor
ranks and symbolic dimensions statically \\
7 & \texttt{InlinePass} (\texttt{inline}) & Inlines small hot
functions \\
8 & \texttt{ConstFoldPass} (\texttt{const-fold}) & Folds safe constants
(deferring overflowing ops to trap with \texttt{E0409}) \\
9 & \texttt{CopyPropPass} (\texttt{copy-prop}) & Eliminates trivial
copies \\
10 & \texttt{TailCallPass} (\texttt{tail-call}) & Lowers self-recursive
tail expressions into SSA loops before LICM \\
11 & \texttt{LicmPass} (\texttt{licm}) & Hoists loop-invariant
computations (seeing new TCO loops!) \\
12 & \texttt{LoopUnrollPass} (\texttt{loop-unroll}) & Unrolls small
fixed-trip loops \\
13 & \texttt{StrengthReducePass} (\texttt{strength-reduce}) & Simplifies
expensive integer/address operations \\
14 & \texttt{MlFusionPass} (\texttt{ml-fusion}) & Cancels double
transposes/negations and collapses redundant reshapes \\
15 & \texttt{OptPass} (\texttt{opt}) & Dead code elimination, common
subexpressions, op expansion \\
16 & \texttt{ExhaustivePass} (\texttt{exhaustive}) & Checks
pattern-matching exhaustiveness (\texttt{E0210}) \\
17 & \texttt{GcAnnotatePass} (\texttt{GcAnnotate}) & Inserts RC
\texttt{retain} / \texttt{release} calls at last use \\
18 & \texttt{FramePoolScopePass} (\texttt{frame\_pool\_scope}) &
Guarantees \texttt{iris\_ui\_pool\_destroy} runs on
\texttt{Return}/\texttt{Panic} paths \\
19 & \texttt{IrLintPass} (\texttt{lint}) & Emits structural IR
diagnostics \\
20 & \texttt{AsyncLowerPass} (\texttt{async-lower}) & Lowers
\texttt{async\ def} functions into pollable state machines after GC
annotation \\
\end{longtable}
}

For native executables, IRIS lowers SSA IR to LLVM IR, compiles
in-process via the dynamically loaded LLVM-C API into an object file,
and links against the embedded C runtime
(\texttt{src/runtime/iris\_runtime.c}) using \texttt{ld.lld} directly on
Windows MinGW (with system \texttt{clang} fallback where direct linking
is absent).

\begin{center}\rule{0.5\linewidth}{0.5pt}\end{center}

\section{12.2 Inspecting IR, LLVM, SIMD, and Profile-Guided
Optimization}\label{inspecting-ir-llvm-simd-and-profile-guided-optimization}

The \texttt{iris} compiler exposes every stage directly via
\texttt{-\/-emit} and optimization flags:

\begin{Shaded}
\begin{Highlighting}[]
\NormalTok{iris prog.iris {-}{-}emit ir              \# Optimized SSA IR (default emit kind)}
\NormalTok{iris prog.iris {-}{-}emit llvm            \# Lowered LLVM IR text for the module}
\NormalTok{iris prog.iris {-}{-}emit llvm{-}complete   \# Complete LLVM IR with runtime declarations}
\NormalTok{iris prog.iris {-}{-}emit simd            \# Vectorized LLVM IR for par for loops}
\NormalTok{iris prog.iris {-}{-}emit graph           \# DOT dependency / model graph}
\NormalTok{iris prog.iris {-}{-}emit jit             \# In{-}process LLVM ORC JIT execution}
\NormalTok{iris prog.iris {-}{-}emit binary {-}o app   \# Standalone native executable}
\NormalTok{iris prog.iris {-}{-}dump{-}ir{-}after licm   \# Dump IR right after a named pass}
\end{Highlighting}
\end{Shaded}

For maximum hot-loop throughput on native builds:

\begin{enumerate}
\def\labelenumi{\arabic{enumi}.}
\item
  \textbf{Vectorization (\texttt{-\/-simd})}: emits explicit AVX2/NEON
  vector kernels for \texttt{par\ for} and enables LLVM loop
  vectorization (\texttt{-\/-emit\ simd} inspects the LLVM IR).
\item
  \textbf{Profile-Guided Optimization (\texttt{-\/-pgo-generate} /
  \texttt{-\/-pgo-use})}:

\begin{Shaded}
\begin{Highlighting}[]
\NormalTok{iris build compute.iris {-}{-}pgo{-}generate {-}o compute\_inst}
\NormalTok{./compute\_inst                    \# Produces iris.profile (or IRIS\_PROFILE\_FILE)}
\NormalTok{iris build compute.iris {-}{-}pgo{-}use iris.profile {-}o compute\_fast}
\end{Highlighting}
\end{Shaded}
\end{enumerate}

\begin{center}\rule{0.5\linewidth}{0.5pt}\end{center}

\section{12.3 Foreign Function Interface and
Bindgen}\label{foreign-function-interface-and-bindgen}

IRIS integrates with C, C++, Rust, and Python cleanly across four
layers:

\begin{enumerate}
\def\labelenumi{\arabic{enumi}.}
\tightlist
\item
  \textbf{Direct C Declarations (\texttt{extern\ def})}: declare C
  symbols and link shared or static libraries.
\item
  \textbf{Dynamic Loading \& Out-Cells (\texttt{std.ffi})}: open shared
  libraries at runtime (\texttt{lib\_open}, \texttt{lib\_call},
  \texttt{call\_i64}, \texttt{call\_f64}, \texttt{call\_str},
  \texttt{lib\_close}) and exchange primitive pointers via manual memory
  cells (\texttt{out\_i64\_cell}, \texttt{out\_set\_i64},
  \texttt{out\_i64}, \texttt{out\_free}):
\end{enumerate}

\begin{Shaded}
\begin{Highlighting}[]
\KeywordTok{bring}\NormalTok{ std}\OperatorTok{.}\NormalTok{ffi}

\KeywordTok{def}\NormalTok{ main() }\OperatorTok{{-}\textgreater{}} \DataTypeTok{i64}\NormalTok{ \{}
    \KeywordTok{val} \FunctionTok{cell} \OperatorTok{=}\NormalTok{ out\_i64\_cell(}\DecValTok{2}\NormalTok{)}
    \FunctionTok{assert}\NormalTok{(}\FunctionTok{cell} \OperatorTok{!=} \DecValTok{0}\NormalTok{)}
\NormalTok{    out\_set\_i64(}\FunctionTok{cell}\NormalTok{, }\DecValTok{0}\NormalTok{, }\DecValTok{42}\NormalTok{)}
\NormalTok{    out\_set\_i64(}\FunctionTok{cell}\NormalTok{, }\DecValTok{1}\NormalTok{, }\DecValTok{100}\NormalTok{)}
    \FunctionTok{assert}\NormalTok{(out\_i64(}\FunctionTok{cell}\NormalTok{, }\DecValTok{0}\NormalTok{) }\OperatorTok{==} \DecValTok{42} \OperatorTok{\&\&}\NormalTok{ out\_i64(}\FunctionTok{cell}\NormalTok{, }\DecValTok{1}\NormalTok{) }\OperatorTok{==} \DecValTok{100}\NormalTok{)}
\NormalTok{    out\_free(}\FunctionTok{cell}\NormalTok{)}
\NormalTok{    println(}\StringTok{"ffi cells: ok"}\NormalTok{)}
    \ControlFlowTok{return} \DecValTok{0}
\NormalTok{\}}
\end{Highlighting}
\end{Shaded}

\begin{enumerate}
\def\labelenumi{\arabic{enumi}.}
\setcounter{enumi}{2}
\item
  \textbf{Automated Header Binding (\texttt{iris\ bindgen})}: generate
  \texttt{.iris} FFI wrappers directly from C/C++ headers (including C++
  classes/overloads via \texttt{-\/-cpp\ -\/-shim-out\ shim.cpp}):

\begin{Shaded}
\begin{Highlighting}[]
\NormalTok{iris bindgen sensor\_api.h {-}l sensor {-}o sensor\_ffi.iris}
\end{Highlighting}
\end{Shaded}
\item
  \textbf{CPython Extension Export (\texttt{-\/-emit\ python-ext})}:
  compile an IRIS module directly into a native Python module
  (\texttt{.pyd} on Windows, \texttt{.so} on Unix) importable from
  Python:

\begin{Shaded}
\begin{Highlighting}[]
\NormalTok{iris my\_kernels.iris {-}{-}emit python{-}ext {-}o my\_kernels.pyd}
\end{Highlighting}
\end{Shaded}
\end{enumerate}

\begin{center}\rule{0.5\linewidth}{0.5pt}\end{center}

\section{12.4 Freestanding Microcontroller
Bundles}\label{freestanding-microcontroller-bundles}

For deterministic firmware without an operating system or heap
allocator, \texttt{iris\ embedded} compiles a strictly verified
allocation-free scalar subset into relocatable object files and C board
headers:

\begin{Shaded}
\begin{Highlighting}[]
\NormalTok{iris embedded control\_loop.iris {-}{-}target cortex{-}m4f}
\NormalTok{iris embedded control\_loop.iris {-}{-}target cortex{-}m33}
\NormalTok{iris embedded control\_loop.iris {-}{-}target esp32{-}c3}
\NormalTok{iris embedded uno\_button\_led.iris {-}{-}target arduino{-}uno}
\end{Highlighting}
\end{Shaded}

Each invocation generates:

\begin{itemize}
\tightlist
\item
  \texttt{iris\_module.ll} and \texttt{iris\_module.o} (verified to
  contain zero heap-allocator symbols)
\item
  \texttt{iris\_embedded\_runtime.o} and
  \texttt{iris\_embedded\_harness.o}
\item
  \texttt{iris\_embedded.h} (declaring bounded GPIO, ADC, PWM, monotonic
  time, watchdog, and UART BSP hooks)
\item
  \texttt{iris\_embedded.json} (cryptographic proof fingerprint + static
  \texttt{.bss} / fixed-array footprint)
\end{itemize}

Top-level mutable scalar variables (\texttt{var\ presses:\ i64\ =\ 0},
\texttt{f64}, or \texttt{bool}) persist across hardware control steps
inside \texttt{.bss} without dynamic allocation:

\begin{Shaded}
\begin{Highlighting}[]
\CommentTok{// illustrative bare{-}metal controller pattern (compile via \textasciigrave{}iris embedded {-}{-}target cortex{-}m4f\textasciigrave{})}
\KeywordTok{extern} \KeywordTok{def}\NormalTok{ iris\_embedded\_gpio\_read(pin: }\DataTypeTok{i64}\NormalTok{) }\OperatorTok{{-}\textgreater{}} \DataTypeTok{i64}\NormalTok{ effect io}
\KeywordTok{extern} \KeywordTok{def}\NormalTok{ iris\_embedded\_gpio\_write(pin: }\DataTypeTok{i64}\NormalTok{, level: }\DataTypeTok{i64}\NormalTok{) }\OperatorTok{{-}\textgreater{}} \DataTypeTok{i64}\NormalTok{ effect io}

\KeywordTok{var}\NormalTok{ last\_state: }\DataTypeTok{i64} \OperatorTok{=} \DecValTok{0}
\KeywordTok{var}\NormalTok{ toggle\_count: }\DataTypeTok{i64} \OperatorTok{=} \DecValTok{0}

\KeywordTok{def}\NormalTok{ main() }\OperatorTok{{-}\textgreater{}} \DataTypeTok{i64}\NormalTok{ effect io \{}
    \KeywordTok{val}\NormalTok{ now }\OperatorTok{=}\NormalTok{ iris\_embedded\_gpio\_read(}\DecValTok{2}\NormalTok{)}
    \ControlFlowTok{if}\NormalTok{ now }\OperatorTok{==} \DecValTok{1} \OperatorTok{\&\&}\NormalTok{ last\_state }\OperatorTok{==} \DecValTok{0}\NormalTok{ \{}
\NormalTok{        toggle\_count }\OperatorTok{=}\NormalTok{ toggle\_count }\OperatorTok{+} \DecValTok{1}
\NormalTok{    \}}
\NormalTok{    last\_state }\OperatorTok{=}\NormalTok{ now}
\NormalTok{    iris\_embedded\_gpio\_write(}\DecValTok{13}\NormalTok{, toggle\_count }\OperatorTok{\%} \DecValTok{2}\NormalTok{)}
    \ControlFlowTok{return} \DecValTok{0}
\NormalTok{\}}
\end{Highlighting}
\end{Shaded}

\begin{quote}
\textbf{Embedded Target Limits}: - \textbf{Arduino Uno
(\texttt{avr-unknown-unknown})}: scalar control flow, top-level state,
checked arithmetic, and board hooks are physically validated; cyclic
CFGs and fixed arrays are rejected on AVR under LLVM 17 safety guards. -
\textbf{Cortex-M4F, Cortex-M33, ESP32-C3}: ELF object bundles,
fixed-array footprints, and host simulations are tested; physical
hardware board reports (\texttt{IRIS-HW/1\ ...}) remain required before
claiming target silicon validation. - \textbf{ESP32 (Xtensa)}: requires
Espressif's LLVM distribution (stock LLVM reports the missing backend
cleanly).
\end{quote}

\subsection{Common Mistakes}\label{common-mistakes-11}

\begin{itemize}
\tightlist
\item
  \textbf{Allocating strings, \texttt{list}, \texttt{record},
  \texttt{choice}, or closures inside \texttt{iris\ embedded} code}: the
  allocation-freedom proof rejects all heap constructs cleanly at
  compile time.
\item
  \textbf{Forgetting \texttt{out\_free(cell)} after
  \texttt{out\_i64\_cell} / \texttt{out\_new}}: manual FFI scratch
  buffers live outside RC tracking unless freed explicitly (though
  \texttt{ui\_native\_*} hosts reuse internal scratch cells
  automatically).
\end{itemize}

\subsection{Try It Yourself}\label{try-it-yourself-11}

\begin{enumerate}
\def\labelenumi{\arabic{enumi}.}
\tightlist
\item
  Inspect the SSA IR and LLVM IR emitted for one of your functions with
  \texttt{iris\ file.iris\ -\/-emit\ ir} vs
  \texttt{iris\ file.iris\ -\/-emit\ llvm}.
\item
  Dump IR right after tail-call optimization using
  \texttt{iris\ file.iris\ -\/-dump-ir-after\ tail-call\ -\/-emit\ ir}.
\item
  Compile \texttt{examples/embedded/control\_loop.iris} with
  \texttt{iris\ embedded\ examples/embedded/control\_loop.iris\ -\/-target\ cortex-m4f\ -o\ out\_m4f}
  and inspect \texttt{out\_m4f/iris\_embedded.json} and
  \texttt{out\_m4f/iris\_embedded.h}.
\end{enumerate}

\begin{center}\rule{0.5\linewidth}{0.5pt}\end{center}

\part{Part IV: Differentiable Programming, ML, and
Autonomy}\label{part-iv-differentiable-programming-ml-and-autonomy}

\chapter{Chapter 13: Automatic
Differentiation}\label{chapter-13-automatic-differentiation}

\section{13.1 Forward-Mode Dual Numbers and Reverse-Mode
Tapes}\label{forward-mode-dual-numbers-and-reverse-mode-tapes}

Differential programming is a core capability of IRIS rather than an
external monkey-patched runtime. IRIS exposes \textbf{two complementary
automatic differentiation modes}:

\begin{enumerate}
\def\labelenumi{\arabic{enumi}.}
\tightlist
\item
  \textbf{Forward-mode dual numbers
  (\texttt{grad\textless{}T\textgreater{}})}: ideal when differentiating
  from few inputs to many outputs, carrying primal and tangent values
  cleanly through arithmetic expressions and \texttt{@differentiable}
  functions (\texttt{grad\textless{}f64\textgreater{}}).
\item
  \textbf{Reverse-mode Wengert tapes (\texttt{tape}, \texttt{backward},
  \texttt{grad})}: ideal when optimizing a scalar loss or objective with
  respect to arbitrarily many parameters in a single backward
  sweep---recording operations dynamically across functions,
  \texttt{if}/\texttt{else} branches, and \texttt{for}/\texttt{while}
  loops.
\end{enumerate}

Use \texttt{tape(x)} when solving inverse problems, trajectory
optimization, policy optimization, or neural training; reach for
\texttt{grad\textless{}f64\textgreater{}} when computing exact
directional derivatives on small analytic formulas.

\begin{center}\rule{0.5\linewidth}{0.5pt}\end{center}

\section{13.2 Reverse-Mode Scalar Autodiff with Tape, Backward, and
Grad}\label{reverse-mode-scalar-autodiff-with-tape-backward-and-grad}

Promote any floating-point input onto the dynamic chunked tape with
\texttt{tape(value)}. All standard arithmetic (\texttt{+}, \texttt{-},
\texttt{*}, \texttt{/}, unary \texttt{-}) and scalar math intrinsics
(\texttt{sqrt}, \texttt{sin}, \texttt{cos}, \texttt{exp}, \texttt{log},
\texttt{abs}) record chain-rule nodes automatically---even across helper
calls, branches, and loop iterations:

\begin{Shaded}
\begin{Highlighting}[]
\KeywordTok{def}\NormalTok{ squared\_prediction(weight: }\DataTypeTok{f64}\NormalTok{, x: }\DataTypeTok{f64}\NormalTok{) }\OperatorTok{{-}\textgreater{}} \DataTypeTok{f64}\NormalTok{ \{}
    \KeywordTok{val}\NormalTok{ pred }\OperatorTok{=}\NormalTok{ weight }\OperatorTok{*}\NormalTok{ x}
    \ControlFlowTok{return}\NormalTok{ pred }\OperatorTok{*}\NormalTok{ pred}
\NormalTok{\}}

\KeywordTok{def}\NormalTok{ main() }\OperatorTok{{-}\textgreater{}} \DataTypeTok{i64}\NormalTok{ \{}
    \KeywordTok{val}\NormalTok{ weight }\OperatorTok{=}\NormalTok{ tape(}\FloatTok{1.0}\NormalTok{)}
    \KeywordTok{var}\NormalTok{ loss }\OperatorTok{=}\NormalTok{ tape(}\FloatTok{0.0}\NormalTok{)}

    \CommentTok{// loss = sum\_\{k=1..4\} (weight * k)\^{}2 = weight\^{}2 * (1 + 4 + 9 + 16) = 30 * weight\^{}2}
    \CommentTok{// d(loss)/d(weight) at weight = 1.0 is 60.0}
    \ControlFlowTok{for}\NormalTok{ sample }\ControlFlowTok{in} \DecValTok{1}\OperatorTok{..}\DecValTok{5}\NormalTok{ \{}
\NormalTok{        loss }\OperatorTok{=}\NormalTok{ loss }\OperatorTok{+}\NormalTok{ squared\_prediction(weight, to\_f64(sample))}
\NormalTok{    \}}

\NormalTok{    backward(loss)}
    \KeywordTok{val}\NormalTok{ slope }\OperatorTok{=} \DataTypeTok{grad}\NormalTok{(weight)}
    \FunctionTok{assert}\NormalTok{(slope }\OperatorTok{==} \FloatTok{60.0}\NormalTok{)}
\NormalTok{    println(}\StringTok{f"scalar autodiff: dL/dw = }\ControlFlowTok{\{}\NormalTok{slope}\ControlFlowTok{\}}\StringTok{"}\NormalTok{)}
    \ControlFlowTok{return} \DecValTok{0}
\NormalTok{\}}
\end{Highlighting}
\end{Shaded}

When designing differentiable physics rollouts or multi-step simulations
for reverse-mode AD, keep two numerical hygiene habits in mind:

\begin{itemize}
\tightlist
\item
  \textbf{Guard denominators against zero}: add a small softening term
  before \texttt{sqrt(...)} (for example
  \texttt{sqrt(dx\ *\ dx\ +\ dy\ *\ dy\ +\ 0.002)}) so
  \texttt{1\ /\ (2\ *\ sqrt(r2))} never emits \texttt{NaN} during
  \texttt{backward(...)}.
\item
  \textbf{Initialize step-by-step accumulators cleanly}: write
  \texttt{0.000001} or \texttt{1.0e-6} with a decimal point before
  \texttt{e} when tuning tolerances and \texttt{Adam} epsilons.
\end{itemize}

\begin{center}\rule{0.5\linewidth}{0.5pt}\end{center}

\section{13.3 Differentiable Tensors in Standard
Tensor}\label{differentiable-tensors-in-standard-tensor}

When optimizing vectors, matrices, or neural layers rather than scalar
parameters, import \texttt{std.tensor} and build graphs with
\texttt{DiffTensor}:

\begin{Shaded}
\begin{Highlighting}[]
\KeywordTok{bring}\NormalTok{ std}\OperatorTok{.}\DataTypeTok{tensor}

\KeywordTok{record}\NormalTok{ SquaredLayer \{}
\NormalTok{    marker: }\DataTypeTok{i64}\NormalTok{,}
\NormalTok{\}}

\KeywordTok{trait}\NormalTok{ Layer \{}
    \KeywordTok{def}\NormalTok{ forward(self, }\KeywordTok{input}\NormalTok{: DiffTensor) }\OperatorTok{{-}\textgreater{}}\NormalTok{ DiffTensor effect alloc}
\NormalTok{\}}

\KeywordTok{impl}\NormalTok{ Layer }\ControlFlowTok{for}\NormalTok{ SquaredLayer \{}
    \KeywordTok{def}\NormalTok{ forward(self, }\KeywordTok{input}\NormalTok{: DiffTensor) }\OperatorTok{{-}\textgreater{}}\NormalTok{ DiffTensor effect alloc \{}
        \ControlFlowTok{return}\NormalTok{ tensor\_ad\_mul(}\KeywordTok{input}\NormalTok{, }\KeywordTok{input}\NormalTok{)}
\NormalTok{    \}}
\NormalTok{\}}

\KeywordTok{def}\NormalTok{ main() }\OperatorTok{{-}\textgreater{}} \DataTypeTok{i64}\NormalTok{ \{}
    \KeywordTok{val}\NormalTok{ shape: }\DataTypeTok{list}\OperatorTok{\textless{}}\DataTypeTok{i64}\OperatorTok{\textgreater{}} \OperatorTok{=} \DataTypeTok{list}\NormalTok{(}\DecValTok{2}\NormalTok{)}
    \KeywordTok{val}\NormalTok{ weights }\OperatorTok{=}\NormalTok{ tensor\_ad\_parameter(tensor\_ad\_from\_data(}\DataTypeTok{list}\NormalTok{(}\FloatTok{1.0}\NormalTok{, }\FloatTok{2.0}\NormalTok{), shape))}
    \KeywordTok{val}\NormalTok{ factor }\OperatorTok{=}\NormalTok{ tensor\_ad\_from\_data(}\DataTypeTok{list}\NormalTok{(}\FloatTok{2.0}\NormalTok{, }\FloatTok{3.0}\NormalTok{), shape)}

    \KeywordTok{var}\NormalTok{ total }\OperatorTok{=}\NormalTok{ tensor\_ad\_from\_data(}\DataTypeTok{list}\NormalTok{(}\FloatTok{0.0}\NormalTok{, }\FloatTok{0.0}\NormalTok{), shape)}
    \ControlFlowTok{for}\NormalTok{ \_step }\ControlFlowTok{in} \DecValTok{0}\OperatorTok{..}\DecValTok{3}\NormalTok{ \{}
\NormalTok{        total }\OperatorTok{=}\NormalTok{ tensor\_ad\_add(total, tensor\_ad\_mul(weights, factor))}
\NormalTok{    \}}

    \KeywordTok{val} \KeywordTok{layer} \OperatorTok{=}\NormalTok{ SquaredLayer \{ marker: }\DecValTok{0}\NormalTok{ \}}
    \KeywordTok{val} \KeywordTok{output} \OperatorTok{=} \KeywordTok{layer}\OperatorTok{.}\NormalTok{forward(total)}
\NormalTok{    tensor\_ad\_backward(tensor\_ad\_sum(}\KeywordTok{output}\NormalTok{))}

    \KeywordTok{val}\NormalTok{ g }\OperatorTok{=}\NormalTok{ tensor\_ad\_data(tensor\_ad\_grad(weights))}
    \FunctionTok{assert}\NormalTok{(}\FunctionTok{list\_get}\NormalTok{(g, }\DecValTok{0}\NormalTok{) }\OperatorTok{==} \FloatTok{72.0}\NormalTok{)}
    \FunctionTok{assert}\NormalTok{(}\FunctionTok{list\_get}\NormalTok{(g, }\DecValTok{1}\NormalTok{) }\OperatorTok{==} \FloatTok{324.0}\NormalTok{)}
\NormalTok{    println(}\StringTok{"tensor autodiff: ok"}\NormalTok{)}
    \ControlFlowTok{return} \DecValTok{0}
\NormalTok{\}}
\end{Highlighting}
\end{Shaded}

Notice how seamlessly \texttt{DiffTensor} integrates with IRIS traits
(\texttt{trait\ Layer}) and loops---while executing identically on the
interpreter and native LLVM binaries.

\begin{center}\rule{0.5\linewidth}{0.5pt}\end{center}

\section{13.4 Optimization Loops: Gradient Descent and
Adam}\label{optimization-loops-gradient-descent-and-adam}

Combining \texttt{tape}, \texttt{backward}, and \texttt{grad} inside an
iterative loop turns any physical simulation into an inverse-problem
solver. Each optimization step starts fresh variables on the tape,
computes loss, calls \texttt{backward(loss)}, reads \texttt{grad(...)},
and updates ordinary plain-\texttt{f64} weights:

\begin{Shaded}
\begin{Highlighting}[]
\KeywordTok{def}\NormalTok{ target\_residual(x: }\DataTypeTok{f64}\NormalTok{) }\OperatorTok{{-}\textgreater{}} \DataTypeTok{f64}\NormalTok{ \{}
    \KeywordTok{val}\NormalTok{ err\_val }\OperatorTok{=}\NormalTok{ x }\OperatorTok{{-}} \FloatTok{3.0}
    \ControlFlowTok{return}\NormalTok{ err\_val }\OperatorTok{*}\NormalTok{ err\_val}
\NormalTok{\}}

\KeywordTok{def}\NormalTok{ main() }\OperatorTok{{-}\textgreater{}} \DataTypeTok{i64}\NormalTok{ \{}
    \KeywordTok{var}\NormalTok{ x }\OperatorTok{=} \FloatTok{0.0}
    \KeywordTok{val}\NormalTok{ lr }\OperatorTok{=} \FloatTok{0.25}

    \ControlFlowTok{for}\NormalTok{ \_iter }\ControlFlowTok{in} \DecValTok{0}\OperatorTok{..}\DecValTok{20}\NormalTok{ \{}
        \KeywordTok{val}\NormalTok{ xt }\OperatorTok{=}\NormalTok{ tape(x)}
        \KeywordTok{val}\NormalTok{ loss }\OperatorTok{=}\NormalTok{ target\_residual(xt)}
\NormalTok{        backward(loss)}
\NormalTok{        x }\OperatorTok{=}\NormalTok{ x }\OperatorTok{{-}}\NormalTok{ lr }\OperatorTok{*} \DataTypeTok{grad}\NormalTok{(xt)}
\NormalTok{    \}}

    \FunctionTok{assert}\NormalTok{(x }\OperatorTok{\textgreater{}} \FloatTok{2.999} \OperatorTok{\&\&}\NormalTok{ x }\OperatorTok{\textless{}} \FloatTok{3.001}\NormalTok{)}
\NormalTok{    println(}\StringTok{f"gradient descent converged: x = }\ControlFlowTok{\{}\NormalTok{x}\ControlFlowTok{\}}\StringTok{"}\NormalTok{)}
    \ControlFlowTok{return} \DecValTok{0}
\NormalTok{\}}
\end{Highlighting}
\end{Shaded}

For multi-parameter non-convex landscapes (like optimizing an orbital
slingshot velocity vector \texttt{(vx,\ vy)} around multiple gravity
wells in \texttt{projects/grav\_forge}), swap vanilla \texttt{lr\ *\ g}
for first- and second-moment Adam updates
(\texttt{m\ =\ beta1\ *\ m\ +\ (1\ -\ beta1)\ *\ g},
\texttt{v\ =\ beta2\ *\ v\ +\ (1\ -\ beta2)\ *\ g\ *\ g}) alongside
learning-rate schedules and gradient clipping.

\subsection{Common Mistakes}\label{common-mistakes-12}

\begin{itemize}
\tightlist
\item
  \textbf{Calling \texttt{grad(x)} before \texttt{backward(loss)}}:
  gradients populate only after running the reverse pass from the scalar
  \texttt{loss}.
\item
  \textbf{Trying to feed \texttt{tape(...)} values directly across outer
  epochs without re-taping}: re-wrap current \texttt{f64} parameters
  with \texttt{tape(param)} at the start of each optimization iteration
  so previous epoch graphs are reclaimed cleanly.
\item
  \textbf{Writing \texttt{sqrt(x\ *\ x\ +\ y\ *\ y)} without softening}:
  if \texttt{x\ ==\ 0\ \&\&\ y\ ==\ 0}, the derivative of
  \texttt{sqrt(0)} diverges (\texttt{Inf}/\texttt{NaN}). Add a positive
  \texttt{eps} inside \texttt{sqrt(...)}.
\end{itemize}

\subsection{Try It Yourself}\label{try-it-yourself-12}

\begin{enumerate}
\def\labelenumi{\arabic{enumi}.}
\tightlist
\item
  Minimize the Rosenbrock valley
  \(f(x, y) = (1 - x)^2 + 100 (y - x^2)^2\) starting from
  \texttt{(-1.0,\ 1.0)} using \texttt{tape} + \texttt{backward} + Adam
  and assert final distance from \texttt{(1.0,\ 1.0)} is below
  \texttt{0.01}.
\item
  Differentiate \texttt{loss\ =\ sin(w\ *\ x)\ +\ 0.1\ *\ w\ *\ w} in a
  multi-sample loop and assert the gradient against a finite-difference
  check \texttt{(f(w\ +\ eps)\ -\ f(w\ -\ eps))\ /\ (2.0\ *\ eps)}.
\item
  Train a tiny single linear layer using \texttt{std.tensor} and
  \texttt{std.nn.nn\_ad\_linear} / \texttt{nn\_ad\_squared\_error}.
\end{enumerate}

\begin{center}\rule{0.5\linewidth}{0.5pt}\end{center}

\chapter{Chapter 14: Tensors, Neural Networks, and Model
Export}\label{chapter-14-tensors-neural-networks-and-model-export}

\section{14.1 Typed Tensors and Compile-Time Shape
Checking}\label{typed-tensors-and-compile-time-shape-checking}

Static shape safety eliminates one of the most frequent sources of
runtime failures in numerical and ML code. IRIS represents dense
multi-dimensional tensors with element type and symbolic dimensions in
the type system:

\begin{Shaded}
\begin{Highlighting}[]
\CommentTok{// Compile{-}checked tensor multiplication: [M, K] @ [K, N] {-}\textgreater{} [M, N]}
\KeywordTok{def}\NormalTok{ linear\_project(}
\NormalTok{    a: }\DataTypeTok{tensor}\OperatorTok{\textless{}}\DataTypeTok{f32}\NormalTok{, [M, K]}\OperatorTok{\textgreater{}}\NormalTok{,}
\NormalTok{    b: }\DataTypeTok{tensor}\OperatorTok{\textless{}}\DataTypeTok{f32}\NormalTok{, [K, N]}\OperatorTok{\textgreater{}}\NormalTok{,}
\NormalTok{) }\OperatorTok{{-}\textgreater{}} \DataTypeTok{tensor}\OperatorTok{\textless{}}\DataTypeTok{f32}\NormalTok{, [M, N]}\OperatorTok{\textgreater{}}\NormalTok{ \{}
    \ControlFlowTok{return}\NormalTok{ a @ b}
\NormalTok{\}}
\end{Highlighting}
\end{Shaded}

During compilation, \textbf{\texttt{ShapeCheckPass}} validates rank
compatibility, contraction dimensions
(\texttt{{[}M,\ K{]}\ @\ {[}K,\ N{]}}), element-wise shape alignment,
and scalar broadcasting---producing diagnostic \texttt{E0204} before a
binary is ever emitted if two matrix dimensions disagree.

You can also move between dense arrays and sparse encodings via
\texttt{sparsify(...)} / \texttt{densify(...)}
(\texttt{sparse\textless{}T\textgreater{}}).

\begin{center}\rule{0.5\linewidth}{0.5pt}\end{center}

\section{14.2 Neural Layers, Classical ML, Datasets, and
RL}\label{neural-layers-classical-ml-datasets-and-rl}

Beyond low-level \texttt{tensor\textless{}...\textgreater{}} types, the
standard library equips you with higher-level training and evaluation
modules:

\begin{itemize}
\tightlist
\item
  \textbf{\texttt{std.tensor}}: \texttt{DiffTensor}, matrix
  multiplication, reductions (\texttt{tensor\_ad\_sum}), activations,
  and gradient extraction.
\item
  \textbf{\texttt{std.nn}}: differentiable layers
  (\texttt{nn\_ad\_linear}, \texttt{nn\_ad\_squared\_error}) as well as
  multi-layer perceptron helpers (\texttt{mlp\_create},
  \texttt{mlp\_forward}, \texttt{mlp\_backward}).
\item
  \textbf{\texttt{std.dataset}}: normalization, splitting, shuffling,
  and batching (\texttt{dataset\_mean}, \texttt{dataset\_std},
  \texttt{dataset\_normalize}, \texttt{dataset\_split},
  \texttt{dataset\_get\_batch}).
\item
  \textbf{\texttt{std.ml}}: classical ML records, tensor pairs
  (\texttt{MLTensor}, \texttt{tensor\_pair}), production health
  telemetry (\texttt{model\_health\_new},
  \texttt{model\_health\_observe}, \texttt{model\_health\_error\_rate}),
  and model versioning (\texttt{model\_registry\_new},
  \texttt{model\_registry\_add}).
\item
  \textbf{\texttt{std.rl}}: replay buffers (\texttt{Experience},
  \texttt{ReplayBuffer}, \texttt{replay\_buffer\_new},
  \texttt{replay\_buffer\_push}, \texttt{replay\_buffer\_sample}) and
  tabular Q-learning (\texttt{q\_table\_new},
  \texttt{q\_learning\_update}).
\end{itemize}

Here is a runnable differentiable linear layer built cleanly on
\texttt{std.tensor} and \texttt{std.nn}:

\begin{Shaded}
\begin{Highlighting}[]
\KeywordTok{bring}\NormalTok{ std}\OperatorTok{.}\DataTypeTok{tensor}
\KeywordTok{bring}\NormalTok{ std}\OperatorTok{.}\NormalTok{nn}

\KeywordTok{def}\NormalTok{ main() }\OperatorTok{{-}\textgreater{}} \DataTypeTok{i64}\NormalTok{ \{}
    \KeywordTok{val} \KeywordTok{input} \OperatorTok{=}\NormalTok{ tensor\_ad\_from\_data(}\DataTypeTok{list}\NormalTok{(}\FloatTok{1.0}\NormalTok{, }\FloatTok{2.0}\NormalTok{), }\DataTypeTok{list}\NormalTok{(}\DecValTok{1}\NormalTok{, }\DecValTok{2}\NormalTok{))}
    \KeywordTok{val}\NormalTok{ weights }\OperatorTok{=}\NormalTok{ tensor\_ad\_parameter(tensor\_ad\_from\_data(}\DataTypeTok{list}\NormalTok{(}\FloatTok{3.0}\NormalTok{, }\FloatTok{4.0}\NormalTok{), }\DataTypeTok{list}\NormalTok{(}\DecValTok{2}\NormalTok{, }\DecValTok{1}\NormalTok{)))}
    \KeywordTok{val}\NormalTok{ bias }\OperatorTok{=}\NormalTok{ tensor\_ad\_from\_data(}\DataTypeTok{list}\NormalTok{(}\FloatTok{1.0}\NormalTok{), }\DataTypeTok{list}\NormalTok{(}\DecValTok{1}\NormalTok{, }\DecValTok{1}\NormalTok{))}
    \KeywordTok{val}\NormalTok{ target }\OperatorTok{=}\NormalTok{ tensor\_ad\_from\_data(}\DataTypeTok{list}\NormalTok{(}\FloatTok{0.0}\NormalTok{), }\DataTypeTok{list}\NormalTok{(}\DecValTok{1}\NormalTok{, }\DecValTok{1}\NormalTok{))}

    \KeywordTok{val}\NormalTok{ pred }\OperatorTok{=}\NormalTok{ nn\_ad\_linear(}\KeywordTok{input}\NormalTok{, weights, bias)}
    \KeywordTok{val}\NormalTok{ loss }\OperatorTok{=}\NormalTok{ nn\_ad\_squared\_error(pred, target)}
\NormalTok{    tensor\_ad\_backward(loss)}

    \KeywordTok{val}\NormalTok{ g }\OperatorTok{=}\NormalTok{ tensor\_ad\_data(tensor\_ad\_grad(weights))}
    \FunctionTok{assert}\NormalTok{(}\FunctionTok{list\_get}\NormalTok{(tensor\_ad\_data(pred), }\DecValTok{0}\NormalTok{) }\OperatorTok{==} \FloatTok{12.0}\NormalTok{)}
    \FunctionTok{assert}\NormalTok{(}\FunctionTok{list\_get}\NormalTok{(g, }\DecValTok{0}\NormalTok{) }\OperatorTok{==} \FloatTok{24.0} \OperatorTok{\&\&} \FunctionTok{list\_get}\NormalTok{(g, }\DecValTok{1}\NormalTok{) }\OperatorTok{==} \FloatTok{48.0}\NormalTok{)}
\NormalTok{    println(}\StringTok{"neural layer: ok"}\NormalTok{)}
    \ControlFlowTok{return} \DecValTok{0}
\NormalTok{\}}
\end{Highlighting}
\end{Shaded}

\begin{center}\rule{0.5\linewidth}{0.5pt}\end{center}

\section{14.3 The Declarative Model DSL and Graph
Fusion}\label{the-declarative-model-dsl-and-graph-fusion}

Neural architectures can be declared directly as top-level
\texttt{model} graphs using contextual keywords (\texttt{input},
\texttt{layer}, \texttt{output}):

\begin{Shaded}
\begin{Highlighting}[]
\CommentTok{// Inspect graph: iris model.iris {-}{-}emit graph}
\CommentTok{// Export ONNX:   iris model.iris {-}{-}emit onnx{-}binary {-}o classifier.onnx}
\KeywordTok{model}\NormalTok{ Classifier \{}
    \KeywordTok{input}\NormalTok{ features: }\DataTypeTok{tensor}\OperatorTok{\textless{}}\DataTypeTok{f32}\NormalTok{, [}\DecValTok{1}\NormalTok{, }\DecValTok{4}\NormalTok{]}\OperatorTok{\textgreater{}}
    \KeywordTok{layer}\NormalTok{ hidden Dense(units }\OperatorTok{=} \DecValTok{8}\NormalTok{)}
    \KeywordTok{layer}\NormalTok{ activated ReLU(hidden)}
    \KeywordTok{output}\NormalTok{ activated}
\NormalTok{\}}
\end{Highlighting}
\end{Shaded}

When compiling tensor graphs and expressions,
\textbf{\texttt{MlFusionPass}} automatically cancels redundant identity
transposes, fuses consecutive transposes, collapses no-op reshapes, and
simplifies double negations at the IR level before code generation.

\begin{center}\rule{0.5\linewidth}{0.5pt}\end{center}

\section{14.4 ONNX, TensorRT Export, and Dynamic ML
Backends}\label{onnx-tensorrt-export-and-dynamic-ml-backends}

IRIS connects seamlessly to industry inference engines in two
directions:

\begin{enumerate}
\def\labelenumi{\arabic{enumi}.}
\tightlist
\item
  \textbf{Exporting IRIS \texttt{model} definitions}:

  \begin{itemize}
  \tightlist
  \item
    \texttt{-\/-emit\ graph} outputs an inspectable DOT/graph
    representation.
  \item
    \texttt{-\/-emit\ onnx} and
    \texttt{-\/-emit\ onnx-binary\ {[}-\/-weights\ FILE{]}} serialize
    the network directly to ONNX protobuf format (injecting trained
    weights from JSON \texttt{"layer.weight":\ {[}...{]}} or seeded
    initializers).
  \item
    \texttt{-\/-emit\ tensorrt\ {[}-\/-weights\ FILE{]}} builds a
    serialized TensorRT engine via \texttt{trtexec}. \emph{(Audit note:
    \texttt{trtexec} CLI invocation is covered by automated tests, while
    physical execution on hosts with NVIDIA TensorRT installed has not
    yet been run.)}
  \end{itemize}
\item
  \textbf{Importing trained models at runtime (\texttt{std.ml})}:

  \begin{itemize}
  \tightlist
  \item
    IRIS discovers \textbf{ONNX Runtime}, \textbf{TensorFlow C API}, and
    \textbf{LibTorch} dynamically via \texttt{dlopen} on first use
    (\texttt{onnx\_load}/\texttt{onnx\_run},
    \texttt{pytorch\_load}/\texttt{pytorch\_run},
    \texttt{torch\_load}/\texttt{torch\_run}), carrying zero
    compile-time SDK dependency when unused.
  \item
    Pass \texttt{IRIS\_ML\_STATIC\_LINK=1} if you explicitly want direct
    static linking (\texttt{MlLinkMode::StaticDirect}), and
    \texttt{IRIS\_USE\_BLAS=1} to accelerate matrix multiplication with
    OpenBLAS, Intel MKL, or macOS Accelerate.
  \end{itemize}
\end{enumerate}

\begin{center}\rule{0.5\linewidth}{0.5pt}\end{center}

\section{14.5 LLM Tool Calling and Structured
Agents}\label{llm-tool-calling-and-structured-agents}

For modern agentic services talking to local runtimes (Ollama, vLLM,
llama.cpp) or remote OpenAI-compatible endpoints, \texttt{std.llm}
provides typed message envelopes, schema-checked tool calls, embeddings,
and JSON serialization:

\begin{Shaded}
\begin{Highlighting}[]
\KeywordTok{bring}\NormalTok{ std}\OperatorTok{.}\NormalTok{llm}

\KeywordTok{def}\NormalTok{ main() }\OperatorTok{{-}\textgreater{}} \DataTypeTok{i64}\NormalTok{ \{}
    \KeywordTok{val}\NormalTok{ messages: }\DataTypeTok{list}\OperatorTok{\textless{}}\NormalTok{LlmMessage}\OperatorTok{\textgreater{}} \OperatorTok{=} \DataTypeTok{list}\NormalTok{(llm\_message(}\StringTok{"user"}\NormalTok{, }\StringTok{"Say \textbackslash{}"}\NormalTok{hello\textbackslash{}}\StringTok{""}\NormalTok{))}
    \KeywordTok{val}\NormalTok{ client }\OperatorTok{=}\NormalTok{ llm\_client(}
        \StringTok{"http://127.0.0.1:8080/v1/chat/completions"}\NormalTok{,}
        \StringTok{"http://127.0.0.1:8080/v1/embeddings"}\NormalTok{,}
        \StringTok{"local{-}model"}\NormalTok{,}
        \StringTok{""}\NormalTok{,}
        \DecValTok{5000}\NormalTok{,}
\NormalTok{    )}

    \KeywordTok{val}\NormalTok{ req }\OperatorTok{=}\NormalTok{ llm\_request\_json(client, messages, }\FloatTok{0.0}\NormalTok{, }\DecValTok{64}\NormalTok{, }\StringTok{""}\NormalTok{)}
    \FunctionTok{assert}\NormalTok{(}\FunctionTok{contains}\NormalTok{(req, }\StringTok{"\textbackslash{}"}\NormalTok{model\textbackslash{}}\StringTok{":\textbackslash{}"}\NormalTok{local}\OperatorTok{{-}}\KeywordTok{model}\NormalTok{\textbackslash{}}\StringTok{""}\NormalTok{))}

    \KeywordTok{val}\NormalTok{ resp }\OperatorTok{=}\NormalTok{ llm\_response\_from\_json(}
        \DecValTok{200}\NormalTok{,}
        \StringTok{"\{\textbackslash{}"}\NormalTok{choices\textbackslash{}}\StringTok{":[\{\textbackslash{}"}\NormalTok{message\textbackslash{}}\StringTok{":\{\textbackslash{}"}\NormalTok{content\textbackslash{}}\StringTok{":\textbackslash{}"}\NormalTok{hello\textbackslash{}}\StringTok{"\}\}]\}"}\NormalTok{,}
\NormalTok{    )}
    \FunctionTok{assert}\NormalTok{(}\FunctionTok{is\_ok}\NormalTok{(resp) }\OperatorTok{\&\&}\NormalTok{ unwrap(resp)}\OperatorTok{.}\NormalTok{content }\OperatorTok{==} \StringTok{"hello"}\NormalTok{)}
    \FunctionTok{assert}\NormalTok{(}\FunctionTok{is\_err}\NormalTok{(llm\_response\_from\_json(}\DecValTok{503}\NormalTok{, }\StringTok{"unavailable"}\NormalTok{)))}
\NormalTok{    println(}\StringTok{"llm protocol: ok"}\NormalTok{)}
    \ControlFlowTok{return} \DecValTok{0}
\NormalTok{\}}
\end{Highlighting}
\end{Shaded}

To make agent tool execution deterministic and safe, decode incoming
JSON tool arguments into validated typed IRIS \texttt{record} structs
before running domain logic (\texttt{projects/llm\_agent}).

\subsection{Common Mistakes}\label{common-mistakes-13}

\begin{itemize}
\tightlist
\item
  \textbf{Assuming \texttt{model} DSL blocks define an imperative
  \texttt{main()} entry point}: inspect or export pure
  \texttt{model\ Name\ \{\ ...\ \}} files with \texttt{-\/-emit\ graph},
  \texttt{-\/-emit\ onnx}, or \texttt{-\/-emit\ onnx-binary}.
\item
  \textbf{Enabling HTTPS LLM endpoints on non-Windows native builds
  without checking TLS}: inspect \texttt{http\_tls\_available()} or
  target a local plaintext HTTP sidecar when deploying Linux/macOS
  agents natively until non-Windows TLS is locked.
\end{itemize}

\subsection{Try It Yourself}\label{try-it-yourself-13}

\begin{enumerate}
\def\labelenumi{\arabic{enumi}.}
\tightlist
\item
  Save the \texttt{Classifier} \texttt{model} above to
  \texttt{classifier.iris} and run
  \texttt{iris\ classifier.iris\ -\/-emit\ graph} and
  \texttt{iris\ classifier.iris\ -\/-emit\ onnx} to inspect the
  generated computation graph.
\item
  Normalize a synthetic 1D measurement list with
  \texttt{std.dataset.dataset\_normalize} and verify unit-scale
  statistics with \texttt{std.math.mean} and \texttt{std.math.std\_dev}.
\item
  Explore \texttt{projects/learning\_service/main.iris} to see
  deterministic training, model versioning, and health degradation
  tracking working together.
\end{enumerate}

\begin{center}\rule{0.5\linewidth}{0.5pt}\end{center}

\chapter{Chapter 15: Autonomous Intelligent Systems, Robotics, and
Governed
Evolution}\label{chapter-15-autonomous-intelligent-systems-robotics-and-governed-evolution}

\section{15.1 Active Inference, Uncertainty, Adaptive Control, and
Lifecycle}\label{active-inference-uncertainty-adaptive-control-and-lifecycle}

IRIS was engineered for long-lived autonomous and homeostatic software.
Four domain modules capture online statistics, epistemic uncertainty,
adaptation, and degradation handling cleanly:

\begin{itemize}
\tightlist
\item
  \textbf{\texttt{std.ais}}: running Welford statistics
  (\texttt{RunningStats}), generative world models
  (\texttt{WorldModel}), viable homeostatic envelopes, genomic
  representation (\texttt{Genome}), Elastic Weight Consolidation
  (\texttt{ewc\_penalty\_grad}), and agent health states
  (\texttt{agent\_runtime\_new}, \texttt{agent\_record\_outcome}).
  Specific tested active-inference capabilities are demonstrated in
  \texttt{projects/gridworld\_inference} (tabular, fully observed,
  seeded runs of Dirichlet exploration + expected free energy).
\item
  \textbf{\texttt{std.uncertainty} \& \texttt{std.adaptive}}:
  closed-form Normal-Inverse-Gamma conjugate Bayesian posteriors
  (\texttt{nig\_prior}, \texttt{nig\_update},
  \texttt{nig\_predictive\_variance}, \texttt{nig\_credible\_interval})
  and online risk monitors (\texttt{RiskMetrics}, \texttt{Uncertainty}).
\item
  \textbf{\texttt{std.meta\_learning}}: fast inner-loop MAML and Reptile
  adaptation helpers (\texttt{maml\_inner\_update}, \texttt{maml\_eval},
  \texttt{reptile\_step}).
\end{itemize}

\begin{Shaded}
\begin{Highlighting}[]
\KeywordTok{bring}\NormalTok{ std}\OperatorTok{.}\NormalTok{ml}
\KeywordTok{bring}\NormalTok{ std}\OperatorTok{.}\NormalTok{ais}
\KeywordTok{bring}\NormalTok{ std}\OperatorTok{.}\NormalTok{uncertainty}

\KeywordTok{def}\NormalTok{ main() }\OperatorTok{{-}\textgreater{}} \DataTypeTok{i64}\NormalTok{ \{}
    \CommentTok{// Bayesian uncertainty update over noisy telemetry}
    \KeywordTok{var}\NormalTok{ post }\OperatorTok{=}\NormalTok{ nig\_default()}
\NormalTok{    post }\OperatorTok{=}\NormalTok{ nig\_update\_batch(post, }\DataTypeTok{list}\NormalTok{(}\FloatTok{10.0}\NormalTok{, }\FloatTok{10.2}\NormalTok{, }\FloatTok{9.8}\NormalTok{, }\FloatTok{10.1}\NormalTok{))}
    \FunctionTok{assert}\NormalTok{(nig\_mean(post) }\OperatorTok{\textgreater{}} \FloatTok{9.9} \OperatorTok{\&\&}\NormalTok{ nig\_mean(post) }\OperatorTok{\textless{}} \FloatTok{10.1}\NormalTok{)}

    \CommentTok{// Production agent degradation monitoring}
    \KeywordTok{var}\NormalTok{ agent }\OperatorTok{=}\NormalTok{ agent\_runtime\_new(}\DecValTok{2}\NormalTok{)}
\NormalTok{    agent }\OperatorTok{=}\NormalTok{ agent\_record\_outcome(agent, }\FloatTok{1.0}\NormalTok{, }\OperatorTok{{-}}\FloatTok{1.0}\NormalTok{, }\ConstantTok{false}\NormalTok{)}
\NormalTok{    agent }\OperatorTok{=}\NormalTok{ agent\_record\_outcome(agent, }\FloatTok{1.0}\NormalTok{, }\OperatorTok{{-}}\FloatTok{1.0}\NormalTok{, }\ConstantTok{false}\NormalTok{)}
    \FunctionTok{assert}\NormalTok{(agent}\OperatorTok{.}\NormalTok{state }\OperatorTok{==} \StringTok{"degraded"}\NormalTok{)}

\NormalTok{    println(}\StringTok{"ais and uncertainty: ok"}\NormalTok{)}
    \ControlFlowTok{return} \DecValTok{0}
\NormalTok{\}}
\end{Highlighting}
\end{Shaded}

\begin{center}\rule{0.5\linewidth}{0.5pt}\end{center}

\section{15.2 ROS 2 Robotics
Integration}\label{ros-2-robotics-integration}

\texttt{std.ros2} provides direct robotics primitives---geometry
transforms (\texttt{Vector3}, \texttt{Quaternion}, \texttt{Pose},
\texttt{Transform}, \texttt{transform\_point},
\texttt{transform\_inverse}), QoS profiles (\texttt{qos\_sensor\_data},
\texttt{qos\_control}), and live middleware bindings to
\textbf{\texttt{rclcpp}}, \textbf{\texttt{rclcpp\_action}},
\textbf{\texttt{rclcpp\_lifecycle}}, and \textbf{\texttt{tf2\_ros}} via
a thin C-ABI bridge (\texttt{iris\_ros2.dll} /
\texttt{libiris\_ros2.so}).

Offline geometry and QoS helpers work anywhere with zero ROS
installation:

\begin{Shaded}
\begin{Highlighting}[]
\KeywordTok{bring}\NormalTok{ std}\OperatorTok{.}\NormalTok{ros2}

\KeywordTok{def}\NormalTok{ main() }\OperatorTok{{-}\textgreater{}} \DataTypeTok{i64}\NormalTok{ \{}
    \KeywordTok{val}\NormalTok{ tf }\OperatorTok{=}\NormalTok{ Transform \{}
\NormalTok{        translation: Vector3 \{ x: }\FloatTok{1.0}\NormalTok{, y: }\FloatTok{2.0}\NormalTok{, z: }\FloatTok{3.0}\NormalTok{ \},}
\NormalTok{        rotation: Quaternion \{ x: }\FloatTok{0.0}\NormalTok{, y: }\FloatTok{0.0}\NormalTok{, z: }\FloatTok{0.0}\NormalTok{, w: }\FloatTok{1.0}\NormalTok{ \},}
\NormalTok{    \}}
    \KeywordTok{val}\NormalTok{ pt }\OperatorTok{=}\NormalTok{ transform\_point(tf, Point \{ x: }\FloatTok{1.0}\NormalTok{, y: }\FloatTok{1.0}\NormalTok{, z: }\FloatTok{1.0}\NormalTok{ \})}
    \KeywordTok{val}\NormalTok{ orig }\OperatorTok{=}\NormalTok{ transform\_point(transform\_inverse(tf), pt)}

    \FunctionTok{assert}\NormalTok{(pt}\OperatorTok{.}\NormalTok{x }\OperatorTok{==} \FloatTok{2.0} \OperatorTok{\&\&}\NormalTok{ pt}\OperatorTok{.}\NormalTok{y }\OperatorTok{==} \FloatTok{3.0} \OperatorTok{\&\&}\NormalTok{ pt}\OperatorTok{.}\NormalTok{z }\OperatorTok{==} \FloatTok{4.0}\NormalTok{)}
    \FunctionTok{assert}\NormalTok{(orig}\OperatorTok{.}\NormalTok{x }\OperatorTok{==} \FloatTok{1.0} \OperatorTok{\&\&}\NormalTok{ qos\_sensor\_data()}\OperatorTok{.}\NormalTok{depth }\OperatorTok{==} \DecValTok{5}\NormalTok{)}
\NormalTok{    println(}\StringTok{"ros2 geometry: ok"}\NormalTok{)}
    \ControlFlowTok{return} \DecValTok{0}
\NormalTok{\}}
\end{Highlighting}
\end{Shaded}

\begin{quote}
\textbf{ROS 2 Verification Boundary}: Live \texttt{rclcpp} topics,
services, actions, parameters, lifecycle nodes, and \texttt{tf2} are
verified on Windows Humble (Cyclone DDS) and Linux Humble (Fast DDS),
including \texttt{projects/turtlesim\_autopilot},
\texttt{projects/lifecycle\_sensor\_fusion}, and
\texttt{projects/evolving\_turtle}. macOS builds remain untested, each
context uses one single-threaded executor, and tests run against
simulators rather than physical robot hardware.
\end{quote}

\begin{center}\rule{0.5\linewidth}{0.5pt}\end{center}

\section{15.3 Typed Metaprogramming and
Reflection}\label{typed-metaprogramming-and-reflection}

Programs running under the compiler/interpreter host can inspect AST/IR
metadata, type-check candidate patches through the real Rust compiler
pipeline, and evaluate modified fragments dynamically via
\textbf{\texttt{std.meta}} and \textbf{\texttt{std.reflect}} (plus the
\texttt{iris\ meta} subcommand):

\begin{Shaded}
\begin{Highlighting}[]
\CommentTok{// Run under compiler host (\textasciigrave{}IRIS\_FORCE\_INTERP=1 iris {-}{-}emit eval ...\textasciigrave{})}
\KeywordTok{bring}\NormalTok{ std}\OperatorTok{.}\NormalTok{meta}
\KeywordTok{bring}\NormalTok{ std}\OperatorTok{.}\NormalTok{reflect}

\KeywordTok{def}\NormalTok{ main() }\OperatorTok{{-}\textgreater{}} \DataTypeTok{i64}\NormalTok{ \{}
    \FunctionTok{assert}\NormalTok{(meta\_available())}
    \KeywordTok{val}\NormalTok{ source }\OperatorTok{=} \StringTok{"def main() {-}\textgreater{} i64 \{ return 0 \}"}
    \KeywordTok{val}\NormalTok{ prog }\OperatorTok{=}\NormalTok{ meta\_program(source)}
    \FunctionTok{assert}\NormalTok{(}\FunctionTok{is\_ok}\NormalTok{(meta\_analyze(prog)))}

    \KeywordTok{val}\NormalTok{ edit }\OperatorTok{=}\NormalTok{ MetaEdit \{ start\_byte: }\DecValTok{27}\NormalTok{, end\_byte: }\DecValTok{28}\NormalTok{, replacement: }\StringTok{"7"}\NormalTok{ \}}
    \KeywordTok{val}\NormalTok{ edited }\OperatorTok{=}\NormalTok{ unwrap(meta\_apply(prog, edit))}
    \FunctionTok{assert}\NormalTok{(unwrap(reflect\_eval\_i64(edited)) }\OperatorTok{==} \DecValTok{7}\NormalTok{)}
\NormalTok{    println(}\StringTok{"metaprogramming: ok"}\NormalTok{)}
    \ControlFlowTok{return} \DecValTok{0}
\NormalTok{\}}
\end{Highlighting}
\end{Shaded}

In standalone ahead-of-time native binaries, \texttt{meta\_available()}
and \texttt{reflect\_available()} return \texttt{false} cleanly (unless
hosted via compiler/FFI evolution APIs) rather than silently shipping an
unverified embedded compiler.

\begin{center}\rule{0.5\linewidth}{0.5pt}\end{center}

\section{15.4 Seven-Gate Governed
Self-Evolution}\label{seven-gate-governed-self-evolution}

A system that updates its own logic in production cannot rely on
guesswork or unguarded eval loops. IRIS's \textbf{\texttt{iris\ evolve}}
command enforces a mandatory \textbf{seven-gate verification protocol}
for policies (\texttt{i64}, \texttt{f64:N}, or \texttt{record} ABIs) and
whole stateful programs
(\texttt{-\/-manifest\ program.json\ -\/-scenarios\ scenarios.json}):

\begin{enumerate}
\def\labelenumi{\arabic{enumi}.}
\tightlist
\item
  \textbf{Gate 1 --- Language safety}: full parser, type, NLL borrow,
  strict effect, and RC-safety validation passes cleanly.
\item
  \textbf{Gate 2 --- Resources \& ABI}: source size, function count, IR
  block/instruction budgets, and exact binary ABI signatures match the
  active generation.
\item
  \textbf{Gate 3 --- Deterministic tests}: candidate tests execute
  cleanly inside the bounded interpreter/sandbox without hangs or
  assertion traps.
\item
  \textbf{Gate 4 --- Transactional rollback}:
  \texttt{test\_transaction\_rollback()} passes, proving speculative
  mutation restoration works cleanly.
\item
  \textbf{Gate 5 --- Safety constitution}: SHA-256-pinned immutable
  safety rules (\texttt{-\/-constitution} +
  \texttt{-\/-constitution-sha256}) and \texttt{-\/-min-output} /
  \texttt{-\/-max-output} / \texttt{-\/-bound} invariants evaluate every
  candidate input/output pair outside candidate code.
\item
  \textbf{Gate 6 --- Shadow canary}: baseline and candidate run
  identical canary cases; promotion requires strictly improved pass
  counts and lower aggregate error.
\item
  \textbf{Gate 7 --- Verified activation \& audit chain}: LLVM ORC JIT
  compiles the verified candidate, atomically activates a lease-guarded
  generation (preserving 8 rollback generations by default via
  \texttt{HotSwapLease}), and appends a canonical hash-chained JSONL
  entry verified against an independently stored \texttt{-\/-audit-head}
  checkpoint (with optional HMAC-SHA256 authentication).
\end{enumerate}

Related subcommands extend governed evolution across continuous loops
(\texttt{iris\ evolve-live}), live state-carrying HTTP services
(\texttt{iris\ evolve-serve}), genetic AST policy search
(\texttt{iris\ evolve-search}), SLA-monitoring autonomic daemons
(\texttt{iris\ service-daemon}), and explicit opt-in unrestricted host
replacement (\texttt{iris\ evolve-unrestricted}, documented separately
for self-modifying host engines).

\subsection{Common Mistakes}\label{common-mistakes-14}

\begin{itemize}
\tightlist
\item
  \textbf{Expecting \texttt{std.meta} inside a standalone
  \texttt{iris\ build} binary}: compiler-hosted metaprogramming returns
  \texttt{meta\_available()\ ==\ false} in AOT binaries; run hosted
  edits with \texttt{IRIS\_FORCE\_INTERP=1} or drive evolution via
  \texttt{iris\ evolve} / \texttt{crates/iris\_ffi}.
\item
  \textbf{Storing \texttt{-\/-audit} and \texttt{-\/-audit-head} in the
  same untrusted directory}: for production threat models, store
  \texttt{-\/-audit-head} checkpoints and HMAC keys in an isolated trust
  domain.
\end{itemize}

\subsection{Try It Yourself}\label{try-it-yourself-14}

\begin{enumerate}
\def\labelenumi{\arabic{enumi}.}
\tightlist
\item
  Compute conjugate posterior intervals for a noisy synthetic telemetry
  stream using \texttt{std.uncertainty.nig\_update\_batch} and
  \texttt{nig\_credible\_interval}.
\item
  Inspect \texttt{examples/09\_evolution/typed\_edit.iris} under
  \texttt{IRIS\_FORCE\_INTERP=1\ iris\ -\/-emit\ eval\ examples/09\_evolution/typed\_edit.iris}.
\item
  Walk through \texttt{projects/autonomous\_evolution\_lab/} to watch
  all seven evolution gates approve a superior controller and reject an
  unsafe candidate.
\end{enumerate}

\begin{center}\rule{0.5\linewidth}{0.5pt}\end{center}

\part{Part V: Native Desktop
Applications}\label{part-v-native-desktop-applications}

\chapter{Chapter 16: Building Native Desktop UIs with Standard
UI}\label{chapter-16-building-native-desktop-uis-with-standard-ui}

\section{16.1 Architecture: IRIS Immediate-Mode Trees and Retained
Native
Host}\label{architecture-iris-immediate-mode-trees-and-retained-native-host}

IRIS provides a built-in native GUI stack through
\textbf{\texttt{std.ui}} backed by
\textbf{\texttt{crates/iris\_ui\_host}}, compiled as a shared library
(\texttt{iris\_ui\_host.dll} on Windows, \texttt{libiris\_ui\_host.so}
on Linux, \texttt{libiris\_ui\_host.dylib} on macOS). Rather than
embedding a browser engine or wrapping bloated external frameworks, IRIS
splits responsibilities cleanly:

\begin{itemize}
\tightlist
\item
  \textbf{Pure IRIS layer (\texttt{std.ui*})}: constructs immediate-mode
  widget trees (\texttt{UiTree}), runs iterative flexbox/grid
  box-constraint layout (\texttt{BoxConstraints}), hit-tests
  pointer/keyboard events into \texttt{UiState}, animates springs
  (\texttt{std.ui\_motion}), and encodes compact binary vector commands
  (\texttt{DrawBuffer}).
\item
  \textbf{Native Rust host (\texttt{iris\_ui\_host})}: manages OS
  windows and DPI scaling (\texttt{winit}), presents pixels with
  zero-allocation framebuffer reuse (\texttt{softbuffer}), rasterizes 2D
  anti-aliased paths, rounded rectangles, gradients, shadows, SVG icons,
  arcs, and live camera frames (\texttt{tiny-skia} SIMD), shapes text
  cleanly (\texttt{cosmic-text} layout cache), and synchronizes
  screen-reader accessibility trees (\texttt{accesskit}).
\end{itemize}

Importing \texttt{bring\ std.ui} transitively re-exports all six UI
submodules (\texttt{std.ui\_theme}, \texttt{std.ui\_protocol},
\texttt{std.ui\_motion}, \texttt{std.ui\_icons},
\texttt{std.ui\_layout}, and \texttt{std.ui\_widgets}). Because layout
and rendering are separated cleanly from OS event pumping, you can
render any UI tree straight to an \textbf{SVG string
(\texttt{ui\_render\_svg})}, an \textbf{HTML mirror}, or an
\textbf{offscreen PNG snapshot (\texttt{-\/-headless-snapshot})} for
automated tests without opening a visible desktop window.

\begin{center}\rule{0.5\linewidth}{0.5pt}\end{center}

\section{16.2 Scaffolding a Desktop App and Running the Frame
Loop}\label{scaffolding-a-desktop-app-and-running-the-frame-loop}

Start a desktop application cleanly using the \texttt{ui} project
template:

\begin{Shaded}
\begin{Highlighting}[]
\NormalTok{iris new ui mission\_console}
\NormalTok{cd mission\_console}
\NormalTok{cargo build {-}{-}release {-}p iris\_ui\_host   \# once in workspace root if building from source}
\NormalTok{iris build {-}{-}gui main.iris {-}o mission\_console}
\end{Highlighting}
\end{Shaded}

Even without opening a window, \texttt{std.ui} trees construct, lay out,
and serialize to vector SVG cleanly:

\begin{Shaded}
\begin{Highlighting}[]
\KeywordTok{bring}\NormalTok{ std}\OperatorTok{.}\NormalTok{ui}

\KeywordTok{def}\NormalTok{ main() }\OperatorTok{{-}\textgreater{}} \DataTypeTok{i64}\NormalTok{ \{}
    \KeywordTok{val}\NormalTok{ tree }\OperatorTok{=}\NormalTok{ ui\_new()}
    \KeywordTok{val}\NormalTok{ header }\OperatorTok{=}\NormalTok{ ui\_heading(tree, }\StringTok{"hdr"}\NormalTok{, }\StringTok{"IRIS Mission Console"}\NormalTok{, UiTone}\OperatorTok{.}\NormalTok{Accent)}
    \KeywordTok{val}\NormalTok{ btn }\OperatorTok{=}\NormalTok{ ui\_button(tree, }\StringTok{"btn\_arm"}\NormalTok{, }\StringTok{"Arm Telemetry"}\NormalTok{, }\StringTok{"act:arm"}\NormalTok{)}
    \KeywordTok{val}\NormalTok{ card }\OperatorTok{=}\NormalTok{ ui\_card(tree, }\StringTok{"main\_card"}\NormalTok{, UiStyle \{ gap: }\FloatTok{10.0}\NormalTok{ \}, }\DataTypeTok{list}\NormalTok{(header, btn))}
    \KeywordTok{val}\NormalTok{ laid\_out }\OperatorTok{=}\NormalTok{ ui\_layout(ui\_set\_root(tree, card), }\FloatTok{640.0}\NormalTok{, }\FloatTok{360.0}\NormalTok{)}

    \KeywordTok{val}\NormalTok{ svg }\OperatorTok{=}\NormalTok{ ui\_render\_svg(laid\_out)}
    \FunctionTok{assert}\NormalTok{(}\FunctionTok{contains}\NormalTok{(svg, }\StringTok{"\textless{}svg"}\NormalTok{))}
    \FunctionTok{assert}\NormalTok{(}\FunctionTok{contains}\NormalTok{(svg, }\StringTok{"IRIS Mission Console"}\NormalTok{))}
\NormalTok{    println(}\StringTok{"std.ui headless tree and svg: ok"}\NormalTok{)}
    \ControlFlowTok{return} \DecValTok{0}
\NormalTok{\}}
\end{Highlighting}
\end{Shaded}

For full interactive desktop windows, simple state machines can use
\texttt{UiAppConfig} + \texttt{ui\_run\_app} with explicit closure
parameter annotations
(\texttt{\textbar{}st:\ UiState,\ tree:\ UiTree\textbar{}\ render\_fn(st,\ tree)}).
Multi-panel applications with custom application records drive an
explicit \texttt{while\ ui\_native\_is\_open(lib)} frame loop paired
with \texttt{ui\_frame\_pool\_begin()} (detailed in Section 16.6).

\begin{center}\rule{0.5\linewidth}{0.5pt}\end{center}

\section{16.3 Responsive Layout Containers and
Breakpoints}\label{responsive-layout-containers-and-breakpoints}

IRIS organizes widgets into nested layout containers driven by
\texttt{UiStyle} constraints (\texttt{padding}, \texttt{gap},
\texttt{flex}, \texttt{align}, \texttt{min\_w}, \texttt{min\_h},
\texttt{max\_w}, \texttt{max\_h}):

{\def\LTcaptype{none} % do not increment counter
\begin{longtable}[]{@{}
  >{\raggedright\arraybackslash}p{(\linewidth - 2\tabcolsep) * \real{0.5000}}
  >{\raggedright\arraybackslash}p{(\linewidth - 2\tabcolsep) * \real{0.5000}}@{}}
\toprule\noalign{}
\begin{minipage}[b]{\linewidth}\raggedright
Container
\end{minipage} & \begin{minipage}[b]{\linewidth}\raggedright
Role
\end{minipage} \\
\midrule\noalign{}
\endhead
\bottomrule\noalign{}
\endlastfoot
\texttt{ui\_col(tree,\ key,\ style,\ children)} & Vertical flex stack \\
\texttt{ui\_row(tree,\ key,\ style,\ children)} & Horizontal flex
container \\
\texttt{ui\_stack(tree,\ key,\ style,\ children)} & Shared-bounds
layered overlay container \\
\texttt{ui\_card(tree,\ key,\ style,\ children)} & Bordered surface
container using active theme tokens \\
\texttt{ui\_wrap(tree,\ key,\ style,\ children)} & Responsive wrapping
row \\
\texttt{ui\_grid(tree,\ key,\ cols,\ gap,\ children)} & Fixed/adaptive
column grid \\
\texttt{ui\_split\_pane(tree,\ key,\ ratio,\ left,\ right)} &
Proportional two-pane workspace divider \\
\texttt{ui\_spacer(tree,\ flex,\ min\_w,\ min\_h)} & Elastic layout
spacer \\
\texttt{ui\_divider(tree,\ label)} & Section divider with optional
inline caption \\
\end{longtable}
}

Adapt layouts automatically across window widths with
\textbf{\texttt{UiBreakpoint}} (\texttt{Compact}, \texttt{Medium},
\texttt{Expanded}, \texttt{Wide}), \texttt{ui\_breakpoint(width)}, and
\texttt{ui\_responsive\_cols(width,\ compact,\ medium,\ expanded)}:

\begin{Shaded}
\begin{Highlighting}[]
\KeywordTok{bring}\NormalTok{ std}\OperatorTok{.}\NormalTok{ui}

\KeywordTok{def}\NormalTok{ main() }\OperatorTok{{-}\textgreater{}} \DataTypeTok{i64}\NormalTok{ \{}
    \KeywordTok{val}\NormalTok{ narrow\_cols }\OperatorTok{=}\NormalTok{ ui\_responsive\_cols(}\FloatTok{520.0}\NormalTok{, }\DecValTok{1}\NormalTok{, }\DecValTok{2}\NormalTok{, }\DecValTok{4}\NormalTok{)}
    \KeywordTok{val}\NormalTok{ wide\_cols }\OperatorTok{=}\NormalTok{ ui\_responsive\_cols(}\FloatTok{1280.0}\NormalTok{, }\DecValTok{1}\NormalTok{, }\DecValTok{2}\NormalTok{, }\DecValTok{4}\NormalTok{)}
    \FunctionTok{assert}\NormalTok{(narrow\_cols }\OperatorTok{==} \DecValTok{1}\NormalTok{)}
    \FunctionTok{assert}\NormalTok{(wide\_cols }\OperatorTok{==} \DecValTok{4}\NormalTok{)}
\NormalTok{    println(}\StringTok{"responsive layout helpers: ok"}\NormalTok{)}
    \ControlFlowTok{return} \DecValTok{0}
\NormalTok{\}}
\end{Highlighting}
\end{Shaded}

\begin{center}\rule{0.5\linewidth}{0.5pt}\end{center}

\section{16.4 Built-In Themes, Vector SVG Icons, and Core
Widgets}\label{built-in-themes-vector-svg-icons-and-core-widgets}

\texttt{std.ui\_theme} ships \textbf{six curated color and elevation
palettes} out of the box---every widget reads semantic theme colors
(\texttt{bg}, \texttt{panel}, \texttt{border}, \texttt{text},
\texttt{muted}, \texttt{accent}, \texttt{teal}, \texttt{warn},
\texttt{danger}) so runtime theme switching requires one function swap:

\begin{itemize}
\tightlist
\item
  \texttt{obsidian\_theme()} --- canonical dark IRIS studio palette
\item
  \texttt{daylight\_theme()} --- crisp high-readability light palette
\item
  \texttt{midnight\_theme()} --- deep slate navy studio theme
\item
  \texttt{nord\_theme()} --- cool Arctic blue-gray palette
\item
  \texttt{emerald\_studio\_theme()} --- high-contrast emerald command
  theme
\item
  \texttt{high\_contrast\_theme()} --- accessible high-contrast palette
\end{itemize}

You can also derive custom accents
(\texttt{theme\_with\_accent(theme,\ accent)}), parse hex colors
(\texttt{hex\_rgb}), or blend palette stops (\texttt{mix\_color},
\texttt{lighten}, \texttt{darken}, \texttt{contrast\_fg}).

\subsection{\texorpdfstring{Vector SVG Icons
(\texttt{std.ui\_icons})}{Vector SVG Icons (std.ui\_icons)}}\label{vector-svg-icons-std.ui_icons}

\texttt{std.ui\_icons} defines \textbf{99 crisp 24×24 vector
\texttt{UiIcon} variants} (\texttt{Search}, \texttt{Plus},
\texttt{Play}, \texttt{Pause}, \texttt{Settings}, \texttt{Brain},
\texttt{Network}, \texttt{Chart}, \texttt{Bolt}, \texttt{Shield},
\texttt{Terminal}, \texttt{Rocket}, \texttt{Cpu}, \texttt{Camera},
\texttt{Sun}, \texttt{Moon}, and more) rasterized directly via
\texttt{draw\_cmd\_svg\_path} at arbitrary DPI:

\begin{itemize}
\tightlist
\item
  Inside widget trees:
  \texttt{ui\_icon(tree,\ key,\ icon,\ size,\ tone)} and
  \texttt{ui\_icon\_button(tree,\ key,\ icon,\ text,\ tone,\ action)}.
\item
  Directly onto custom canvases: \texttt{draw\_icon},
  \texttt{draw\_icon\_styled}, \texttt{draw\_svg\_icon}, or lookup by
  name with \texttt{ui\_icon\_from\_name("camera")}.
\end{itemize}

\begin{center}\rule{0.5\linewidth}{0.5pt}\end{center}

\section{16.5 Interactive State, Forms, Modals, Drawers, and
Toasts}\label{interactive-state-forms-modals-drawers-and-toasts}

Every interactive control carries a stable string \texttt{key} and emits
typed action events folded into \texttt{UiState} via
\texttt{ui\_process\_events}. Watch the exact signatures for common
widgets:

\begin{itemize}
\tightlist
\item
  \textbf{Buttons}: \texttt{ui\_button(tree,\ key,\ text,\ action)} (4
  args) and \texttt{ui\_button\_tone(tree,\ key,\ text,\ action,\ tone)}
  (5 args).
\item
  \textbf{Switches \& Choices}: \texttt{ui\_checkbox},
  \texttt{ui\_toggle}, \texttt{ui\_radio}, \texttt{ui\_select},
  \texttt{ui\_segmented}, \texttt{ui\_tabs}, \texttt{ui\_nav\_item},
  \texttt{ui\_accordion}, \texttt{ui\_chip}, \texttt{ui\_menu},
  \texttt{ui\_stepper}.
\item
  \textbf{Sliders}: \texttt{ui\_slider} and \texttt{ui\_slider\_range}.
\item
  \textbf{Text Entry}:

  \begin{itemize}
  \tightlist
  \item
    \texttt{ui\_text\_input(tree,\ key,\ text,\ placeholder,\ style)}
    --- 5th parameter is \texttt{UiStyle}!
  \item
    \texttt{ui\_text\_area(tree,\ key,\ text,\ placeholder,\ min\_h)}
    --- 5th parameter is \texttt{min\_h:\ f64}!
  \end{itemize}
\item
  \textbf{Telemetry \& Visualization}: \texttt{ui\_stat\_card}
  (\texttt{StatSpec}), \texttt{ui\_bar\_chart} (\texttt{BarChartSpec}),
  \texttt{ui\_line\_chart} (\texttt{ChartSeries}),
  \texttt{ui\_sparkline}, \texttt{ui\_progress\_bar},
  \texttt{ui\_progress\_ring}, \texttt{ui\_spinner},
  \texttt{ui\_avatar}, \texttt{ui\_breadcrumb}, \texttt{ui\_image}
  (\texttt{ImageSpec}), \texttt{ui\_radar\_scope}.
\item
  \textbf{Overlays \& Transient Notifications}:
  \texttt{ui\_modal(tree,\ key,\ title,\ open,\ action,\ children)},
  \texttt{ui\_drawer(...)}, and floating toast notifications via
  \texttt{ui\_state\_push\_toast(...)}.
\end{itemize}

\begin{center}\rule{0.5\linewidth}{0.5pt}\end{center}

\section{16.6 Smooth Springs, Animations, and Frame-Pool Memory
Safety}\label{smooth-springs-animations-and-frame-pool-memory-safety}

Animate panels, numeric badges, glow opacities, and camera crosshairs
smoothly using \textbf{\texttt{std.ui\_motion}} (\texttt{SpringConfig}
presets: \texttt{spring\_snappy()}, \texttt{spring\_smooth()},
\texttt{spring\_bouncy()}, \texttt{spring\_gentle()}, plus
\texttt{UiEase}, \texttt{motion\_cubic\_bezier}, and
\texttt{motion\_lerp\_rect}) integrated right into \texttt{UiState}:

\begin{itemize}
\tightlist
\item
  Trigger spring motion via \texttt{ui\_state\_spring\_from} (after
  event folding).
\item
  Advance clocks each frame via
  \texttt{val\ now\ =\ ui\_native\_clock\_secs(lib)} -\textgreater{}
  \texttt{ui\_state\_tick} -\textgreater{}
  \texttt{ui\_state\_spring\_value}.
\item
  Check \texttt{ui\_state\_has\_active\_transitions} to decide whether
  to poll with a short frame timer (\textasciitilde16 ms for 60 FPS
  motion) or sleep longer when idle.
\end{itemize}

\subsection{\texorpdfstring{The Mandatory \texttt{ui\_frame\_pool}
Lifecycle}{The Mandatory ui\_frame\_pool Lifecycle}}\label{the-mandatory-ui_frame_pool-lifecycle}

As Section 8.5 established, record field assignments
(\texttt{app.score\ =\ v}) allocate on whichever allocator is currently
active. Follow this five-step sequence in every interactive
\texttt{std.ui} host loop:

\begin{Shaded}
\begin{Highlighting}[]
\CommentTok{// illustrative desktop host frame loop (requires built \textasciigrave{}iris\_ui\_host\textasciigrave{} library)}
\KeywordTok{bring}\NormalTok{ std}\OperatorTok{.}\NormalTok{ui}

\KeywordTok{def}\NormalTok{ run\_window(lib: NativeUiLib) }\OperatorTok{{-}\textgreater{}} \DataTypeTok{i64}\NormalTok{ \{}
    \KeywordTok{var}\NormalTok{ buf }\OperatorTok{=}\NormalTok{ draw\_buf\_new(}\DecValTok{131072}\NormalTok{)}
    \KeywordTok{var}\NormalTok{ ui\_st }\OperatorTok{=}\NormalTok{ ui\_state\_new()}
    \KeywordTok{var}\NormalTok{ clicks: }\DataTypeTok{i64} \OperatorTok{=} \DecValTok{0}

    \ControlFlowTok{while}\NormalTok{ ui\_native\_is\_open(lib) \{}
        \KeywordTok{val}\NormalTok{ pool }\OperatorTok{=}\NormalTok{ ui\_frame\_pool\_begin()}
        \KeywordTok{val}\NormalTok{ m }\OperatorTok{=}\NormalTok{ ui\_native\_metrics(lib)}
        \KeywordTok{val}\NormalTok{ tree }\OperatorTok{=}\NormalTok{ ui\_new()}
        \KeywordTok{val}\NormalTok{ btn }\OperatorTok{=}\NormalTok{ ui\_button(tree, }\StringTok{"inc"}\NormalTok{, }\StringTok{f"Clicks: }\ControlFlowTok{\{}\NormalTok{clicks}\ControlFlowTok{\}}\StringTok{"}\NormalTok{, }\StringTok{"act:inc"}\NormalTok{)}
        \KeywordTok{val}\NormalTok{ root }\OperatorTok{=}\NormalTok{ ui\_card(tree, }\StringTok{"root"}\NormalTok{, UiStyle \{ padding: }\FloatTok{16.0}\NormalTok{ \}, }\DataTypeTok{list}\NormalTok{(btn))}
        \KeywordTok{val}\NormalTok{ laid }\OperatorTok{=}\NormalTok{ ui\_layout\_with\_state(ui\_set\_root(tree, root), ui\_st, m}\OperatorTok{.}\NormalTok{width, m}\OperatorTok{.}\NormalTok{height)}

\NormalTok{        buf }\OperatorTok{=}\NormalTok{ ui\_paint(laid, ui\_st, buf)}
        \KeywordTok{val}\NormalTok{ \_sub }\OperatorTok{=}\NormalTok{ ui\_native\_submit(lib, buf)}
        \KeywordTok{val}\NormalTok{ events }\OperatorTok{=}\NormalTok{ ui\_native\_poll(lib, }\DecValTok{16}\NormalTok{)}
        \KeywordTok{val}\NormalTok{ next\_st }\OperatorTok{=}\NormalTok{ ui\_process\_events(laid, ui\_st, events)}
        \KeywordTok{val}\NormalTok{ \_cur }\OperatorTok{=}\NormalTok{ ui\_native\_set\_cursor(lib, next\_st}\OperatorTok{.}\NormalTok{active\_cursor)}

        \CommentTok{// Pause pool BEFORE detaching UiState or updating persistent app state!}
        \KeywordTok{val}\NormalTok{ paused }\OperatorTok{=}\NormalTok{ ui\_frame\_pool\_pause()}
\NormalTok{        ui\_st }\OperatorTok{=}\NormalTok{ ui\_state\_detach(next\_st)}
        \KeywordTok{val}\NormalTok{ activated }\OperatorTok{=}\NormalTok{ ui\_st}\OperatorTok{.}\NormalTok{activated\_action}
        \KeywordTok{val}\NormalTok{ \_res }\OperatorTok{=}\NormalTok{ ui\_frame\_pool\_resume(paused)}
        \KeywordTok{val}\NormalTok{ \_dst }\OperatorTok{=}\NormalTok{ ui\_frame\_pool\_destroy(pool)}

        \ControlFlowTok{if}\NormalTok{ activated }\OperatorTok{==} \StringTok{"act:inc"}\NormalTok{ \{}
\NormalTok{            clicks }\OperatorTok{=}\NormalTok{ clicks }\OperatorTok{+} \DecValTok{1}
\NormalTok{        \}}
\NormalTok{    \}}

    \KeywordTok{val}\NormalTok{ \_f }\OperatorTok{=}\NormalTok{ draw\_buf\_free(buf)}
    \ControlFlowTok{return} \DecValTok{0}
\NormalTok{\}}
\end{Highlighting}
\end{Shaded}

Pause -\textgreater{} \texttt{ui\_state\_detach} -\textgreater{} extract
surviving values -\textgreater{} resume -\textgreater{} destroy
-\textgreater{} update app state. Keeping that order prevents
use-after-free bugs completely even across thousands of live frames.

\subsection{Common Mistakes}\label{common-mistakes-15}

\begin{itemize}
\tightlist
\item
  \textbf{Mixing up the 5th argument of \texttt{ui\_text\_input} vs
  \texttt{ui\_text\_area}}: \texttt{ui\_text\_input} expects
  \texttt{UiStyle} whereas \texttt{ui\_text\_area} expects
  \texttt{min\_h:\ f64}.
\item
  \textbf{Re-allocating \texttt{draw\_buf\_new} every frame inside the
  loop}: allocate \texttt{var\ buf\ =\ draw\_buf\_new(...)} \emph{once}
  before the loop and release it \emph{once} after the loop
  (\texttt{draw\_buf\_free(buf)}).
\item
  \textbf{Passing un-annotated closure parameters to
  \texttt{ui\_run\_app}}: write
  \texttt{\textbar{}s:\ UiState,\ t:\ UiTree\textbar{}\ render(s,\ t)}
  rather than untyped \texttt{\textbar{}s,\ t\textbar{}}.
\end{itemize}

\subsection{Try It Yourself}\label{try-it-yourself-15}

\begin{enumerate}
\def\labelenumi{\arabic{enumi}.}
\tightlist
\item
  Build a responsive 2-card SVG telemetry layout in pure IRIS
  (\texttt{ui\_render\_svg}) that switches between 1 column and 2
  columns depending on \texttt{width}.
\item
  Scaffold a starter app via \texttt{iris\ new\ ui\ hello\_gui}, build
  \texttt{iris\_ui\_host} in release mode, and run it in both desktop
  mode and \texttt{-\/-headless-snapshot}.
\item
  Add a modal dialog (\texttt{ui\_modal}) and a theme switcher cycling
  between \texttt{obsidian\_theme()}, \texttt{nord\_theme()}, and
  \texttt{daylight\_theme()}.
\end{enumerate}

\begin{center}\rule{0.5\linewidth}{0.5pt}\end{center}

\chapter{Chapter 17: Custom Vector Painting, Live Camera Capture, and
Standalone
Packaging}\label{chapter-17-custom-vector-painting-live-camera-capture-and-standalone-packaging}

\section{17.1 Custom Vector Canvas Painting with
DrawBuffer}\label{custom-vector-canvas-painting-with-drawbuffer}

Standard widgets compose cleanly with hardware-fast, high-precision
\textbf{2D vector canvases} emitted straight into \texttt{DrawBuffer}
commands and rasterized by \texttt{iris\_ui\_host}:

{\def\LTcaptype{none} % do not increment counter
\begin{longtable}[]{@{}
  >{\raggedright\arraybackslash}p{(\linewidth - 2\tabcolsep) * \real{0.5000}}
  >{\raggedright\arraybackslash}p{(\linewidth - 2\tabcolsep) * \real{0.5000}}@{}}
\toprule\noalign{}
\begin{minipage}[b]{\linewidth}\raggedright
Command
\end{minipage} & \begin{minipage}[b]{\linewidth}\raggedright
Purpose
\end{minipage} \\
\midrule\noalign{}
\endhead
\bottomrule\noalign{}
\endlastfoot
\texttt{draw\_cmd\_round\_rect} /
\texttt{draw\_cmd\_round\_rect\_corners} & Filled/bordered rounded rects
with uniform or per-corner \texttt{UiRadii} \\
\texttt{draw\_cmd\_circle} / \texttt{draw\_cmd\_arc} & Solid rings,
nodes, orbital arcs, sweep radars \\
\texttt{draw\_cmd\_line} / \texttt{draw\_cmd\_polyline} & Solid or
dashed lines, waveforms, orbital trajectory trails \\
\texttt{draw\_cmd\_linear\_gradient} /
\texttt{draw\_cmd\_radial\_gradient} & Multi-stop smooth shader fills
(\texttt{UiGradientStop}) \\
\texttt{draw\_cmd\_box\_shadow} & Soft elevation drop shadows \\
\texttt{draw\_cmd\_path} / \texttt{draw\_cmd\_svg\_path} & Arbitrary
vector paths and SVG \texttt{d} strings \\
\texttt{draw\_cmd\_text} & Font-shaped \texttt{cosmic-text} labels
inside custom canvases \\
\texttt{draw\_cmd\_image\_rgba} & Blit CPU RGBA pixel grids (heatmaps,
neural activations) \\
\texttt{draw\_cmd\_push\_clip} / \texttt{draw\_cmd\_push\_round\_clip} /
\texttt{draw\_cmd\_pop\_clip} & Rectangular and rounded scissor masks \\
\texttt{draw\_cmd\_push\_transform} / \texttt{draw\_cmd\_pop\_transform}
& Affine translation, scaling, and rotation matrices \\
\texttt{draw\_cmd\_push\_opacity} / \texttt{draw\_cmd\_pop\_opacity} &
Grouped alpha blending layers \\
\end{longtable}
}

Because \texttt{ui\_paint} writes all widget commands into a
\texttt{DrawBuffer}, you can paint custom interactive simulation
viewports right inside a widget card's layout rectangle---or compose
full-window canvas layers beneath crisp HUD cards!

\begin{center}\rule{0.5\linewidth}{0.5pt}\end{center}

\section{17.2 Live Native Camera Capture and Synthetic Sensor
Simulation}\label{live-native-camera-capture-and-synthetic-sensor-simulation}

IRIS desktop applications can open hardware webcams (and seamlessly fall
back to a deterministic synthetic sensor feed in CI, VMs, or
\texttt{-\/-sim} mode) directly through
\textbf{\texttt{std.ui\_protocol}} and
\textbf{\texttt{crates/iris\_ui\_host/src/camera.rs}}:

\begin{itemize}
\tightlist
\item
  \textbf{\texttt{ui\_native\_camera\_open(lib,\ device\_index,\ req\_w,\ req\_h,\ synthetic\_only)}}:
  opens device \texttt{device\_index} using Windows Media Foundation
  dynamically on Windows (or activates synthetic sensor stream when
  \texttt{synthetic\_only\ =\ true} or no hardware camera is present),
  returning a positive 1-based \texttt{cam\_id}.
\item
  \textbf{\texttt{ui\_native\_camera\_info(lib,\ cam\_id)\ -\textgreater{}\ UiCameraInfo}}:
  reports live frame \texttt{width}, \texttt{height}, \texttt{mode}
  (\texttt{0=Offline}, \texttt{1=Live}, \texttt{2=Synthetic}), monotonic
  sequence \texttt{seq}, and OS \texttt{device\_name} (for example
  \texttt{"Lenovo\ EasyCamera"}).
\item
  \textbf{\texttt{draw\_cmd\_camera\_frame(buf,\ rect,\ cam\_id,\ mirror,\ tint\_mode)}}:
  blits the latest decoded RGBA camera frame straight onto any
  \texttt{UiRect} inside \texttt{DrawBuffer} with optional mirroring and
  optical tint modes (\texttt{0=RGB}, \texttt{1=NVG}, \texttt{2=FLIR},
  \texttt{3=Slate}).
\item
  \textbf{\texttt{ui\_native\_camera\_sample\_grid(lib,\ cam\_id,\ grid\_w,\ grid\_h,\ mirror,\ out\_cell)}}:
  downsamples the latest frame into an FFI \texttt{out\_cell} byte
  buffer (\texttt{4..=512} resolution) for real-time IRIS
  computer-vision loops (frame differencing, centroid saliency, optical
  motion scoring) without marshaling full-HD pixels through script
  objects every tick!
\item
  \textbf{\texttt{ui\_native\_camera\_set\_sim\_target(...)}} \&
  \textbf{\texttt{ui\_native\_camera\_close(lib,\ cam\_id)}}: inject
  deterministic moving synthetic targets for headless integration tests
  and tear down capture cleanly.
\end{itemize}

\begin{center}\rule{0.5\linewidth}{0.5pt}\end{center}

\section{17.3 Engineering High-Frame-Rate Native
Interfaces}\label{engineering-high-frame-rate-native-interfaces}

Getting smooth 60 FPS responsiveness on a CPU vector rasterizer comes
down to avoiding wasted work in both IRIS script logic and
\texttt{iris\_ui\_host}:

\begin{enumerate}
\def\labelenumi{\arabic{enumi}.}
\tightlist
\item
  \textbf{Always stage the release \texttt{iris\_ui\_host}}: debug
  builds of \texttt{tiny-skia} / \texttt{cosmic-text} run roughly 10×
  slower than release builds. Build
  \texttt{cargo\ build\ -\/-release\ -p\ iris\_ui\_host} and pass
  \texttt{-\/-release} to \texttt{iris\ bundle} (note: modern IRIS
  prefers release \texttt{iris\_ui\_host} automatically during release
  bundles).
\item
  \textbf{Avoid redundant rounded-clip masks}: full-screen
  \texttt{PushRoundClip} masks allocate and composite pixel alpha masks
  across the entire viewport. Keep \texttt{PushClip} rectangular for
  uncurved viewports and paint vector grids/trails in batched
  \texttt{draw\_cmd\_polyline} commands rather than hundreds of tiny
  separate rect/line primitives.
\item
  \textbf{Sample camera grids only when
  \texttt{info.seq\ !=\ last\_cam\_seq}}: skip re-running CV feature
  extraction when the camera background thread has not produced a newer
  frame sequence number.
\item
  \textbf{Poll events instead of blocking unconditionally during active
  motion}: use \texttt{ui\_native\_poll(lib,\ 0)} or short timeouts
  while \texttt{ui\_state\_has\_active\_transitions(ui\_st)} or live
  simulation/camera capture is active, and longer waits when purely
  idle.
\item
  \textbf{Re-use \texttt{DrawBuffer} and detach cleanly}: pre-allocate
  \texttt{draw\_buf\_new(...)} outside the loop and never let persistent
  data leak onto the frame pool.
\end{enumerate}

\emph{(Reference benchmark from this codebase: after clip-mask dirty
bounding, fast scanline camera blitting, and release host staging,
\texttt{projects/target\_takedown} renders 36 headless frames including
camera decode, CV grid analysis, HUD widgets, and PNG encode in
\textasciitilde1.37 s / \textasciitilde38 ms per frame, down from 14.67
s debug/unoptimized).}

\begin{center}\rule{0.5\linewidth}{0.5pt}\end{center}

\section{17.4 Desktop Manifest Configuration and GUI
Builds}\label{desktop-manifest-configuration-and-gui-builds}

Configure standalone desktop properties in the \texttt{{[}app{]}}
section of \texttt{iris.toml}:

\begin{Shaded}
\begin{Highlighting}[]
\KeywordTok{[package]}
\DataTypeTok{name} \OperatorTok{=} \StringTok{"grav\_forge"}
\DataTypeTok{version} \OperatorTok{=} \StringTok{"1.0.0"}
\DataTypeTok{entry} \OperatorTok{=} \StringTok{"main.iris"}

\KeywordTok{[dependencies]}

\KeywordTok{[app]}
\DataTypeTok{title} \OperatorTok{=} \StringTok{"IRIS GravForge"}
\DataTypeTok{identifier} \OperatorTok{=} \StringTok{"io.iris.grav\_forge"}
\DataTypeTok{publisher} \OperatorTok{=} \StringTok{"Moon9t"}
\DataTypeTok{icon} \OperatorTok{=} \StringTok{"assets/icon.png"}
\DataTypeTok{gui} \OperatorTok{=} \ConstantTok{true}
\DataTypeTok{bundle\_ui\_host} \OperatorTok{=} \ConstantTok{true}
\DataTypeTok{window\_width} \OperatorTok{=} \DecValTok{1400}
\DataTypeTok{window\_height} \OperatorTok{=} \DecValTok{860}
\DataTypeTok{min\_width} \OperatorTok{=} \DecValTok{920}
\DataTypeTok{min\_height} \OperatorTok{=} \DecValTok{580}
\DataTypeTok{resizable} \OperatorTok{=} \ConstantTok{true}
\DataTypeTok{categories} \OperatorTok{=} \StringTok{"Science;Graphics;"}
\DataTypeTok{assets} \OperatorTok{=} \OperatorTok{[}\StringTok{"assets/presets.json"}\OperatorTok{]}
\end{Highlighting}
\end{Shaded}

To compile a standalone GUI binary directly on Windows without spawning
a console window behind it, use:

\begin{Shaded}
\begin{Highlighting}[]
\NormalTok{iris build {-}{-}gui main.iris {-}o MyApp.exe}
\end{Highlighting}
\end{Shaded}

Under the hood, \texttt{-\/-gui} links the executable against the
Windows GUI subsystem (\texttt{SUBSYSTEM:WINDOWS} /
\texttt{mainCRTStartup}), while \texttt{-\/-console} explicitly keeps a
console attached during diagnostics.

\begin{center}\rule{0.5\linewidth}{0.5pt}\end{center}

\section{17.5 Bundling, Installers, and Operating System
Registration}\label{bundling-installers-and-operating-system-registration}

When you are ready to ship an installable desktop application---complete
with icon conversion (\texttt{.png}, \texttt{.ico}, \texttt{.icns}),
asset staging, bundled \texttt{iris\_ui\_host} shared library, and OS
metadata---use \textbf{\texttt{iris\ bundle}}:

\begin{Shaded}
\begin{Highlighting}[]
\NormalTok{\# Stage a self{-}contained portable bundle directory (or use {-}{-}format zip)}
\NormalTok{iris bundle . {-}{-}format dir {-}{-}release {-}o dist/MyApp}

\NormalTok{\# Generate platform native distributions}
\NormalTok{iris bundle . {-}{-}format zip {-}o dist/MyApp.zip         \# Portable zip archive}
\NormalTok{iris bundle . {-}{-}format inno {-}o dist/MyApp{-}setup.iss  \# Windows Inno Setup installer script}
\NormalTok{iris bundle . {-}{-}format app {-}o dist/MyApp.app         \# macOS .app bundle}
\NormalTok{iris bundle . {-}{-}format appimage {-}o dist/MyApp.AppDir \# Linux AppDir / AppImage structure}

\NormalTok{\# Install onto the host OS (copies files + creates Start Menu / Registry / .desktop entries)}
\NormalTok{iris bundle . {-}{-}install {-}{-}release}

\NormalTok{\# Cleanly remove an installed app (pass {-}{-}no{-}register to skip shortcut/registry registration)}
\NormalTok{iris bundle . {-}{-}uninstall}
\end{Highlighting}
\end{Shaded}

You can override manifest values from the command line with
\texttt{-\/-title}, \texttt{-\/-identifier}, \texttt{-\/-icon},
\texttt{-\/-gui} / \texttt{-\/-console}, and
\texttt{-\/-install-dir\ \textless{}DIR\textgreater{}}. Combined with
\texttt{-\/-headless-snapshot}, a CI pipeline can build the GUI binary,
verify rendered pixels into PNG, package an archive or installer script,
and smoke-test install/uninstall automatically.

\subsection{Common Mistakes}\label{common-mistakes-16}

\begin{itemize}
\tightlist
\item
  \textbf{Shipping an application executable without
  \texttt{iris\_ui\_host.dll} / \texttt{.so} / \texttt{.dylib} beside
  it}: either set \texttt{bundle\_ui\_host\ =\ true} in
  \texttt{{[}app{]}} and distribute with \texttt{iris\ bundle}, or copy
  the release \texttt{iris\_ui\_host} library next to your binary.
\item
  \textbf{Debugging performance issues against debug
  \texttt{iris\_ui\_host}}: always build
  \texttt{cargo\ build\ -\/-release\ -p\ iris\_ui\_host} and pass
  \texttt{-\/-release} to \texttt{iris\ bundle}.
\item
  \textbf{Adding \texttt{-\/-headless-snapshot} without closing cleanly
  after a deterministic frame count}: support CLI flags
  (\texttt{-\/-headless-snapshot\ {[}path{]}}, \texttt{-\/-sim}) in
  desktop apps so CI and regression scripts can verify startup and
  rendering without manual clicks.
\end{itemize}

\subsection{Try It Yourself}\label{try-it-yourself-16}

\begin{enumerate}
\def\labelenumi{\arabic{enumi}.}
\tightlist
\item
  Paint a custom canvas scene combining a radial gradient background, a
  multi-point \texttt{draw\_cmd\_polyline} sine wave, a circular sweep
  \texttt{draw\_cmd\_arc}, and centered text using \texttt{DrawBuffer}.
\item
  Add a \texttt{-\/-headless-snapshot\ {[}path{]}} flag to
  \texttt{hello\_gui} that opens \texttt{UiHostMode.Headless}, submits 3
  frames, calls \texttt{ui\_native\_save\_png(lib,\ path)}, and exits
  cleanly with \texttt{0}.
\item
  Package a project into a portable bundle with
  \texttt{iris\ bundle\ .\ -\/-format\ dir\ -\/-release\ -o\ target/my\_app\_bundle}
  and launch the standalone binary directly inside that folder.
\end{enumerate}

\begin{center}\rule{0.5\linewidth}{0.5pt}\end{center}

\part{Part VI: Engineering in IRIS}\label{part-vi-engineering-in-iris}

\chapter{Chapter 18: Packages, Testing, Benchmarking, and
Debugging}\label{chapter-18-packages-testing-benchmarking-and-debugging}

\section{18.1 Manifests, Lockfiles, and the Package
Manager}\label{manifests-lockfiles-and-the-package-manager}

Every multi-file IRIS project is anchored by an
\textbf{\texttt{iris.toml}} manifest declaring the package name,
version, entry point, optional registry URL, dependencies, and optional
\texttt{{[}app{]}} desktop metadata:

\begin{Shaded}
\begin{Highlighting}[]
\KeywordTok{[package]}
\DataTypeTok{name}    \OperatorTok{=} \StringTok{"telemetry\_service"}
\DataTypeTok{version} \OperatorTok{=} \StringTok{"0.2.0"}
\DataTypeTok{entry}   \OperatorTok{=} \StringTok{"main.iris"}

\KeywordTok{[dependencies]}
\DataTypeTok{utils}  \OperatorTok{=} \OperatorTok{\{ }\DataTypeTok{path}\OperatorTok{ =} \StringTok{"../shared/utils"}\OperatorTok{ \}}
\DataTypeTok{router} \OperatorTok{=} \OperatorTok{\{ }\DataTypeTok{git}\OperatorTok{ =} \StringTok{"https://github.com/Eclipse{-}Softworks/iris{-}router.git"}\OperatorTok{, }\DataTypeTok{tag}\OperatorTok{ =} \StringTok{"v1.0.0"}\OperatorTok{ \}}
\DataTypeTok{json}   \OperatorTok{=} \OperatorTok{\{ }\DataTypeTok{registry}\OperatorTok{ =} \StringTok{"json"}\OperatorTok{, }\DataTypeTok{version}\OperatorTok{ =} \StringTok{"\^{}1.0.0"}\OperatorTok{ \}}
\end{Highlighting}
\end{Shaded}

Manage dependencies and reproducible builds via \texttt{iris\ pkg} (and
\texttt{iris\ install} shorthand):

{\def\LTcaptype{none} % do not increment counter
\begin{longtable}[]{@{}
  >{\raggedright\arraybackslash}p{(\linewidth - 2\tabcolsep) * \real{0.5000}}
  >{\raggedright\arraybackslash}p{(\linewidth - 2\tabcolsep) * \real{0.5000}}@{}}
\toprule\noalign{}
\begin{minipage}[b]{\linewidth}\raggedright
Command
\end{minipage} & \begin{minipage}[b]{\linewidth}\raggedright
Action
\end{minipage} \\
\midrule\noalign{}
\endhead
\bottomrule\noalign{}
\endlastfoot
\texttt{iris\ pkg\ init} & Create a new \texttt{iris.toml} in the
current directory \\
\texttt{iris\ pkg\ add\ \textless{}name\textgreater{}\ -\/-path\ \textless{}p\textgreater{}}
& Add a local path dependency \\
\texttt{iris\ pkg\ add\ \textless{}name\textgreater{}\ -\/-git\ \textless{}url\textgreater{}}
& Add a Git dependency (\texttt{-\/-branch}, \texttt{-\/-tag}, or
\texttt{-\/-rev}) \\
\texttt{iris\ pkg\ add\ \textless{}name\textgreater{}\ -\/-registry\ \textless{}pkg\textgreater{}\ -\/-version\ \textless{}req\textgreater{}}
& Add a registry semver dependency \\
\texttt{iris\ pkg\ install\ {[}-\/-offline{]}} & Fetch dependencies and
generate/sync \texttt{iris.lock} \\
\texttt{iris\ pkg\ update\ {[}name{]}} & Refresh dependencies within
version constraints \\
\texttt{iris\ pkg\ vendor} & Mirror dependencies locally into
\texttt{.iris/vendor} for air-gapped builds \\
\texttt{iris\ pkg\ check} / \texttt{list} & Inspect and verify
dependency resolution \\
\texttt{iris\ pkg\ build\ {[}-\/-offline{]}} / \texttt{run} & Ensure
dependencies, then build/run entry binary \\
\texttt{iris\ pkg\ publish\ {[}-\/-registry\ R{]}\ {[}-\/-dry-run{]}} &
Publish an immutable version archive to a registry \\
\texttt{iris\ pkg\ search\ \textless{}text\textgreater{}} /
\texttt{serve\ \textless{}dir\textgreater{}} & Search registries or
serve a local HTTP package index \\
\end{longtable}
}

Commit \texttt{iris.lock} alongside \texttt{iris.toml} so team builds
pin exact commit SHAs and SHA-256 content checksums.

\begin{center}\rule{0.5\linewidth}{0.5pt}\end{center}

\section{18.2 Assertion-Backed Testing}\label{assertion-backed-testing}

IRIS discovers and executes every zero-argument \texttt{test\_*}
function across your files when you run \texttt{iris\ test}. Every test
function should return \texttt{0} after asserting real behavioral
invariants---or pair built-in \texttt{assert(...)} with rich failure
messages from \textbf{\texttt{std.testing}} (\texttt{assert\_eq},
\texttt{assert\_approx\_eq}, \texttt{assert\_str\_eq},
\texttt{assert\_true}, \texttt{assert\_some}, \texttt{assert\_none},
\texttt{assert\_ok}, \texttt{assert\_err}, \texttt{assert\_list\_eq},
\texttt{fail}):

\begin{Shaded}
\begin{Highlighting}[]
\KeywordTok{bring}\NormalTok{ std}\OperatorTok{.}\NormalTok{testing}

\KeywordTok{def}\NormalTok{ normalize\_gain(raw: }\DataTypeTok{f64}\NormalTok{) }\OperatorTok{{-}\textgreater{}} \DataTypeTok{f64}\NormalTok{ \{}
    \ControlFlowTok{if}\NormalTok{ raw }\OperatorTok{\textless{}} \FloatTok{0.0}\NormalTok{ \{}
        \ControlFlowTok{return} \FloatTok{0.0}
\NormalTok{    \}}
    \ControlFlowTok{if}\NormalTok{ raw }\OperatorTok{\textgreater{}} \FloatTok{10.0}\NormalTok{ \{}
        \ControlFlowTok{return} \FloatTok{10.0}
\NormalTok{    \}}
    \ControlFlowTok{return}\NormalTok{ raw}
\NormalTok{\}}

\KeywordTok{def}\NormalTok{ test\_gain\_clamping() }\OperatorTok{{-}\textgreater{}} \DataTypeTok{i64}\NormalTok{ \{}
\NormalTok{    assert\_approx\_eq(normalize\_gain(}\OperatorTok{{-}}\FloatTok{3.0}\NormalTok{), }\FloatTok{0.0}\NormalTok{, }\FloatTok{0.000001}\NormalTok{, }\StringTok{"clamps negative input"}\NormalTok{)}
\NormalTok{    assert\_approx\_eq(normalize\_gain(}\FloatTok{4.5}\NormalTok{), }\FloatTok{4.5}\NormalTok{, }\FloatTok{0.000001}\NormalTok{, }\StringTok{"preserves in{-}range gain"}\NormalTok{)}
\NormalTok{    assert\_approx\_eq(normalize\_gain(}\FloatTok{22.0}\NormalTok{), }\FloatTok{10.0}\NormalTok{, }\FloatTok{0.000001}\NormalTok{, }\StringTok{"clamps upper ceiling"}\NormalTok{)}
    \ControlFlowTok{return} \DecValTok{0}
\NormalTok{\}}

\KeywordTok{def}\NormalTok{ main() }\OperatorTok{{-}\textgreater{}} \DataTypeTok{i64}\NormalTok{ \{}
    \FunctionTok{assert}\NormalTok{(test\_gain\_clamping() }\OperatorTok{==} \DecValTok{0}\NormalTok{)}
\NormalTok{    println(}\StringTok{"testing: ok"}\NormalTok{)}
    \ControlFlowTok{return} \DecValTok{0}
\NormalTok{\}}
\end{Highlighting}
\end{Shaded}

Run tests with optional filtering or clean CI output:

\begin{Shaded}
\begin{Highlighting}[]
\NormalTok{iris test}
\NormalTok{iris test projects/ledger/main.iris {-}{-}no{-}color}
\NormalTok{iris watch main.iris {-}{-}test}
\end{Highlighting}
\end{Shaded}

Remember: avoid tests that merely print output without asserting. Every
test should verify exact outputs, boundaries, and negative error paths
(\texttt{assert(is\_err(...))}).

\begin{center}\rule{0.5\linewidth}{0.5pt}\end{center}

\section{18.3 Benchmarking, Profiling, and
Optimization}\label{benchmarking-profiling-and-optimization}

Measure compilation and execution timings objectively with
\textbf{\texttt{iris\ bench}} and \textbf{\texttt{iris\ profile}}:

\begin{itemize}
\tightlist
\item
  \textbf{\texttt{iris\ bench\ \textless{}file.iris\textgreater{}\ {[}-n\ ITERATIONS{]}}}:
  runs 3 warmup passes followed by \texttt{ITERATIONS} (default
  \texttt{10}) measured runs, reporting minimum, maximum, mean, median,
  and standard deviation across parse time, lower + pass-pipeline
  compile time, and native execution time.
\item
  \textbf{\texttt{iris\ profile\ \textless{}file.iris\textgreater{}}}:
  instruments execution to collect per-function inclusive/exclusive
  timings (\texttt{total\_us}, \texttt{self\_us}), call counts, IR
  instruction counts, and folded stacks for flame graphs.
\item
  \textbf{Native acceleration flags}: combine \texttt{-\/-simd},
  \texttt{-\/-pgo-generate} + \texttt{-\/-pgo-use\ FILE}, and (on
  release builds) \texttt{-\/-release} UI host libraries after measuring
  baseline hotspots.
\end{itemize}

\begin{center}\rule{0.5\linewidth}{0.5pt}\end{center}

\section{18.4 Formatting, Diagnostics, JSON Output, and Bug
Reports}\label{formatting-diagnostics-json-output-and-bug-reports}

Keep codebases uniform and CI automation clean using IRIS's built-in
developer tools:

\begin{itemize}
\tightlist
\item
  \textbf{\texttt{iris\ fmt\ {[}file{]}} /
  \texttt{iris\ fmt\ -\/-check\ {[}file{]}}}: formats a single file (or
  all \texttt{*.iris} files in the current directory) cleanly with exit
  code \texttt{1} on formatting drift in \texttt{-\/-check} mode.
\item
  \textbf{\texttt{iris\ check\ {[}paths...{]}}}: runs lexical, parse,
  type, NLL borrow, exhaustiveness, lint, and optional
  \texttt{-\/-strict-effects} checks without linking. Add
  \texttt{-\/-error-format\ json} when integrating external CI bots or
  editors.
\item
  \textbf{\texttt{iris\ explain\ \textless{}CODE\textgreater{}}}: prints
  an extended explanation, root cause, and suggested fix for any
  compiler diagnostic (\texttt{E0001}--\texttt{E0902}).
\item
  \textbf{\texttt{iris\ bug-report\ \textless{}file.iris\textgreater{}}}:
  packages the minimal source file, compiler diagnostics, and toolchain
  environment metadata into a local diagnostic folder if you ever
  encounter an internal compiler error (nothing is uploaded
  automatically).
\end{itemize}

\begin{center}\rule{0.5\linewidth}{0.5pt}\end{center}

\section{18.5 Dual-Mode Debugging: Time-Travel Interpreter and Native
DWARF}\label{dual-mode-debugging-time-travel-interpreter-and-native-dwarf}

\texttt{iris\ dap} speaks the standard Debug Adapter Protocol
(\texttt{vscode-iris} configures it automatically) in two complementary
modes:

\begin{enumerate}
\def\labelenumi{\arabic{enumi}.}
\tightlist
\item
  \textbf{Trace-based interpreter DAP (\texttt{"native":\ false},
  default)}: records deterministic execution traces supporting
  line/conditional/hit-count/logpoint breakpoints, expression watches,
  variable/local inspection, named \texttt{test\_*} launches, and
  \textbf{time-travel debugging (\texttt{stepBack},
  \texttt{reverseContinue})}.
\item
  \textbf{Native DWARF DAP (\texttt{"native":\ true} or
  \texttt{"attach"} with
  \texttt{processId}/\texttt{pid}/\texttt{target})}: compile with
  \textbf{\texttt{iris\ build\ -g\ main.iris}} to emit DWARF debug
  metadata mapping machine instructions to IRIS source lines
  (\texttt{-O0}) across scalars, \texttt{str}, \texttt{record} structs,
  \texttt{choice} enums, fixed arrays, \texttt{option}, \texttt{result},
  tuples, lists, and maps across brought modules. \texttt{iris\ dap}
  relays sessions to GNU GDB (\texttt{gdb\ -i\ dap}, tested with GDB
  17.1 on Windows), \texttt{lldb-dap}, or \texttt{IRIS\_NATIVE\_DAP},
  exposing \texttt{supportsStepBack} automatically whenever the
  underlying native debugger provides reverse execution (e.g.,
  \texttt{rr}).
\end{enumerate}

\subsection{Common Mistakes}\label{common-mistakes-17}

\begin{itemize}
\tightlist
\item
  \textbf{Debugging a native release binary without \texttt{-g}}:
  compile with \texttt{iris\ build\ -g} when launching native GDB/LLDB
  DAP sessions.
\item
  \textbf{Confusing \texttt{-\/-emit\ eval} with trace DAP}: VS Code DAP
  defaults to trace mode for instant
  \texttt{stepBack}/\texttt{reverseContinue}; switch to
  \texttt{"native":\ true} when stepping compiled LLVM binaries or
  attaching by PID.
\end{itemize}

\subsection{Try It Yourself}\label{try-it-yourself-17}

\begin{enumerate}
\def\labelenumi{\arabic{enumi}.}
\tightlist
\item
  Initialize a package with \texttt{iris\ pkg\ init}, inspect
  \texttt{iris.toml}, and run \texttt{iris\ fmt\ -\/-check}.
\item
  Benchmark a computational loop using
  \texttt{iris\ bench\ main.iris\ -n\ 20} and compare timings against
  \texttt{iris\ profile\ main.iris}.
\item
  Build a small program with
  \texttt{iris\ check\ -\/-error-format\ json} and inspect the
  structured JSON diagnostic fields.
\end{enumerate}

\begin{center}\rule{0.5\linewidth}{0.5pt}\end{center}

\chapter{Chapter 19: Architecture Walkthroughs of Real IRIS
Applications}\label{chapter-19-architecture-walkthroughs-of-real-iris-applications}

Seeing how IRIS constructs compose at scale cements every concept
covered in this book. This chapter tours four concrete, verified
applications included directly in the repository under
\texttt{projects/}.

\begin{center}\rule{0.5\linewidth}{0.5pt}\end{center}

\section{19.1 Checked Ledger and Command-Line
Tools}\label{checked-ledger-and-command-line-tools}

\subsection{\texorpdfstring{\texttt{projects/ledger} --- Multi-File
Domain Modeling and Integer
Money}{projects/ledger --- Multi-File Domain Modeling and Integer Money}}\label{projectsledger-multi-file-domain-modeling-and-integer-money}

A compact example of strict domain isolation across
\texttt{"domain.iris"} and \texttt{main.iris}:

\begin{itemize}
\tightlist
\item
  Monetary balances are tracked exclusively in integer
  \texttt{cents:\ i64} rather than floating-point numbers, taking
  advantage of IRIS's checked arithmetic (\texttt{E0409}) so overflowing
  accumulations cannot silently corrupt balances.
\item
  Overdrafts and negative amounts return typed
  \texttt{result\textless{}i64,\ str\textgreater{}} values rather than
  panicking, keeping \texttt{test\_ledger()} fast and verifiable on both
  interpreter and native backends
  (\texttt{iris\ test\ projects/ledger/main.iris}).
\end{itemize}

\subsection{\texorpdfstring{\texttt{projects/cli\_tools} --- Unix-Style
Binary \& Text
Pipelines}{projects/cli\_tools --- Unix-Style Binary \& Text Pipelines}}\label{projectscli_tools-unix-style-binary-text-pipelines}

Four self-contained UNIX utilities (\texttt{csvq}, \texttt{logstat},
\texttt{pcapinspect}, \texttt{tlsfp}) reading arguments from
\texttt{args()} (\texttt{iris\ run\ tool.iris\ -\/-\ ...}):

\begin{itemize}
\tightlist
\item
  Demonstrates \texttt{std.csv}, \texttt{std.json},
  \texttt{std.buffer.ByteBuffer}, \texttt{std.pcap}, and
  \texttt{std.tls} JA3/JA4 TLS ClientHello fingerprinting running at
  native LLVM speed while verified against reference Python scripts
  across both backends.
\end{itemize}

\begin{center}\rule{0.5\linewidth}{0.5pt}\end{center}

\section{19.2 Concurrent Gateways and Self-Upgrading
Services}\label{concurrent-gateways-and-self-upgrading-services}

\subsection{\texorpdfstring{\texttt{projects/async\_gateway} ---
High-Concurrency Cooperative
Fan-Out}{projects/async\_gateway --- High-Concurrency Cooperative Fan-Out}}\label{projectsasync_gateway-high-concurrency-cooperative-fan-out}

Implements an HTTP/1.1 service where every connection and supplier
fan-out is an \texttt{async\ def} task executing on a single cooperative
thread:

\begin{itemize}
\tightlist
\item
  \texttt{/quote/\textless{}sku\textgreater{}} spawns three concurrent
  supplier tasks bounded by a 150 ms deadline (\texttt{timeout\_after}),
  cancelling tardy suppliers cooperatively while returning best pricing
  and cancellation counts.
\item
  \texttt{/fastest/\textless{}sku\textgreater{}} races tasks via
  \texttt{async\_any} + \texttt{cancel\_all}.
\item
  Verified natively under 64 simultaneous HTTP client threads and 300
  concurrent self-test quotes in \textasciitilde85--95 ms natively
  (\texttt{tests/async\_gateway.rs}).
\end{itemize}

\subsection{\texorpdfstring{\texttt{projects/self\_upgrading\_service}
--- Live State-Preserving
Upgrades}{projects/self\_upgrading\_service --- Live State-Preserving Upgrades}}\label{projectsself_upgrading_service-live-state-preserving-upgrades}

Driven by \texttt{iris\ evolve-serve} + \texttt{client.iris}: promotes
upgraded inventory server versions over HTTP while carrying live stock
counters cleanly across generations---rejecting unsafe/retyped versions
at the seven verification gates and preserving rollback safety.

\begin{center}\rule{0.5\linewidth}{0.5pt}\end{center}

\section{19.3 Differentiable Orbital Simulation: IRIS
GravForge}\label{differentiable-orbital-simulation-iris-gravforge}

Located at \texttt{projects/grav\_forge}, \textbf{IRIS GravForge} brings
reverse-mode AD, ADTs, traits, algebraic effects, JSON snapshots, and 60
FPS native UI together into a single interactive physics sandbox:

\begin{Shaded}
\begin{Highlighting}[]
\NormalTok{projects/grav\_forge/}
\NormalTok{├── iris.toml      \# [app] window\_width=1400, gui=true, bundle\_ui\_host=true}
\NormalTok{├── physics.iris   \# BodyKind/ScenarioId/SimEvent choices, CelestialMetric trait,}
\NormalTok{│                  \# algebraic effect telemetry (\textasciigrave{}grav\_audit\_hook\textasciigrave{}), symplectic RK{-}style steppers,}
\NormalTok{│                  \# JSON scene codec}
\NormalTok{├── autodiff.iris  \# Reverse{-}mode burn solver (\textasciigrave{}tape\textasciigrave{}/\textasciigrave{}backward\textasciigrave{}/\textasciigrave{}grad\textasciigrave{} + Adam)}
\NormalTok{├── canvas.iris    \# Batched spacetime metric grid \& multi{-}stop star/black{-}hole rendering}
\NormalTok{└── main.iris      \# Interactive frame loop with \textasciigrave{}ui\_frame\_pool\_pause\textasciigrave{} + 6 themes}
\end{Highlighting}
\end{Shaded}

Key architectural patterns worth adopting in your own scientific or
creative software:

\begin{enumerate}
\def\labelenumi{\arabic{enumi}.}
\tightlist
\item
  \textbf{Reverse-mode AD in the loop}: holding Shift or launching a
  Probe engages \texttt{autodiff.iris}, rolling out smoothed N-body
  gravitation
  (\texttt{soft2\ =\ cfg.softening\ *\ cfg.softening\ +\ 0.0012}
  preventing \texttt{sqrt(0)} singularities), computing miss loss with
  \texttt{tape} + \texttt{backward} + \texttt{grad}, and steering
  \texttt{(vx,\ vy)} live with Adam momentum so trajectories thread
  curved gravity wells toward the beacon.
\item
  \textbf{Polymorphism across heterogeneous entities}:
  \texttt{trait\ CelestialMetric} inspects bodies uniformly, while
  \texttt{choice\ BodyKind} dispatches gravity laws and rendering styles
  via exhaustive \texttt{when}.
\item
  \textbf{Algebraic effect telemetry}: simulation milestones emit
  \texttt{grav\_audit\_hook} intercepted cleanly via
  \texttt{handle\ \{\ ...\ \}\ with\ \{\ ...\ \}} in UI event handlers.
\item
  \textbf{Pure IRIS scene import/export}: snapshots encode/decode JSON
  via \texttt{std.json} + OS clipboard
  (\texttt{ui\_native\_clipboard\_set} / \texttt{get}), tested cleanly
  even without a GPU or display via
  \texttt{test\_grav\_forge\_ad\_and\_physics}.
\end{enumerate}

\begin{center}\rule{0.5\linewidth}{0.5pt}\end{center}

\section{19.4 Live Camera Vision and HUD: Aegis Target
Takedown}\label{live-camera-vision-and-hud-aegis-target-takedown}

Located at \texttt{projects/target\_takedown}, \textbf{IRIS
Aegis-Vision} pairs live native webcam capture with a real-time
computer-vision tactical HUD:

\begin{itemize}
\tightlist
\item
  \textbf{Zero-copy capture + fast sensor grid}:
  \texttt{ui\_native\_camera\_open} captures frames via Windows Media
  Foundation on a background thread (or switches cleanly to
  deterministic synthetic bogeys via \texttt{-\/-sim} or automatic
  fallback when no webcam is present). The main thread only re-runs CV
  detection when \texttt{info.seq\ !=\ last\_cam\_seq}, sampling a
  compact \texttt{64×36} downsampled grid
  (\texttt{ui\_native\_camera\_sample\_grid}) into pre-allocated
  \texttt{VisionWorkspace} FFI cells \emph{outside} the per-frame pool.
\item
  \textbf{Vision \& lock-on state machine}: computes adaptive running
  background delta + chromatic target salience across \texttt{8×6}
  regions, filters track centroids with an alpha-beta velocity
  estimator, projects lead reticles, and drives \texttt{TargetStatus}
  transitions (\texttt{Acquired} -\textgreater{} \texttt{Tracking}
  -\textgreater{} \texttt{Locked} -\textgreater{} \texttt{Engaging}
  -\textgreater{} \texttt{Neutralized}).
\item
  \textbf{Tactical HUD composition}: draws live tinted video
  (\texttt{RGB} / \texttt{NVG} / \texttt{FLIR} / \texttt{SLATE} via
  \texttt{draw\_cmd\_camera\_frame}), spring-animated reticles, lock
  progress rings, telemetry sparklines, and interactive controls with
  strict \texttt{ui\_frame\_pool} pause/detach safety---achieving smooth
  \textasciitilde38 ms/frame headless throughput and real-time
  interactive desktop frame rates when compiled with release
  \texttt{iris\_ui\_host}.
\end{itemize}

\begin{center}\rule{0.5\linewidth}{0.5pt}\end{center}

\section{19.5 Navigating the Companion Examples and
Projects}\label{navigating-the-companion-examples-and-projects}

To continue exploring by domain after finishing this book, inspect the
curated index files in the repository:

\begin{itemize}
\tightlist
\item
  \textbf{\texttt{examples/README.md} (\texttt{examples/catalog.json})}:
  progressive standalone lessons (\texttt{01\_basics/} through
  \texttt{10\_integration/}, plus \texttt{embedded/}) verified
  automatically across native and interpreter backends
  (\texttt{tools/verify\_learning.py},
  \texttt{cargo\ test\ -\/-test\ examples\_showcase}).
\item
  \textbf{\texttt{projects/README.md}}: 30+ full multi-file applications
  spanning active inference (\texttt{gridworld\_inference}), ROS 2
  (\texttt{turtlesim\_autopilot}, \texttt{lifecycle\_sensor\_fusion},
  \texttt{evolving\_turtle}), autonomous control
  (\texttt{quadrotor\_autopilot},
  \texttt{autonomous\_deep\_space\_probe}), networking
  (\texttt{intrusion\_detection\_system},
  \texttt{cyber\_operations\_center}), UI galleries
  (\texttt{ui\_studio}, \texttt{ui\_gallery}, \texttt{terra\_pulse},
  \texttt{ais\_sentinel}), and embedded hardware controls
  (\texttt{uno\_controls}).
\end{itemize}

\begin{center}\rule{0.5\linewidth}{0.5pt}\end{center}

\part{Appendices}\label{appendices}

\chapter{Appendix A: Complete Language Grammar and Reserved
Words}\label{appendix-a-complete-language-grammar-and-reserved-words}

\section{Reserved and Contextual
Keywords}\label{reserved-and-contextual-keywords}

\textbf{Reserved keywords:}

\begin{Shaded}
\begin{Highlighting}[]
\NormalTok{def      val      var      let      return   if       else     while}
\NormalTok{break    continue bring    when     match    choice   for      in}
\NormalTok{async    await    const    type     trait    impl     pub      extern}
\NormalTok{with     dyn      resume   try      catch    raise    move     unsafe}
\NormalTok{defmacro true     false}
\end{Highlighting}
\end{Shaded}

\textbf{Words reserved for future use:}

\begin{Shaded}
\begin{Highlighting}[]
\NormalTok{match  enum  struct  fn  let  mut  use  self  Self  super  crate  where  as  ref  box  macro}
\end{Highlighting}
\end{Shaded}

\textbf{Contextual keywords} (valid as ordinary identifiers everywhere
outside their dedicated syntactic position):

\begin{itemize}
\tightlist
\item
  \texttt{record}, \texttt{model}, \texttt{mod} (when beginning an item
  declaration)
\item
  \texttt{effect} (an \texttt{effect} declaration or
  \texttt{-\textgreater{}\ T\ effect\ ...} function/lambda clause)
\item
  \texttt{loop}, \texttt{defer}, \texttt{yield}, \texttt{select} (when
  beginning a statement in their syntactic form)
\item
  \texttt{spawn} (statement followed by \texttt{\{} or
  \texttt{(group)\ \{})
\item
  \texttt{par} (before \texttt{for})
\item
  \texttt{to} (binary conversion expression \texttt{x\ to\ f64})
\item
  \texttt{by} (range step modifier \texttt{0..n\ by\ 2})
\item
  \texttt{input}, \texttt{layer}, \texttt{output} (inside
  \texttt{model\ \{\ ...\ \}} blocks)
\end{itemize}

\section{Core Grammar Summary (BNF)}\label{core-grammar-summary-bnf}

\begin{Shaded}
\begin{Highlighting}[]
\NormalTok{module       ::= \{ top\_level \}}
\NormalTok{top\_level    ::= function\_def | record\_def | enum\_def | const\_def}
\NormalTok{               | type\_alias | trait\_def | impl\_def | bring\_decl}
\NormalTok{               | extern\_def | model\_def | macro\_def}

\NormalTok{bring\_decl   ::= [ "pub" ] "bring" ( IDENT \{ "." IDENT \} | STRING\_LIT )}
\NormalTok{function\_def ::= [ "pub" ] [ "async" | "const" ] "def" IDENT [ type\_params ]}
\NormalTok{                 "(" params ")" "{-}\textgreater{}" type [ effect\_clause ] block}
\NormalTok{type\_params  ::= "[" type\_param \{ "," type\_param \} "]"}
\NormalTok{params       ::= [ param \{ "," param \} ]}
\NormalTok{param        ::= IDENT ":" type [ "=" expr ]}

\NormalTok{record\_def   ::= [ "pub" ] "record" IDENT [ type\_params ] "\{" field\_defs "\}"}
\NormalTok{field\_defs   ::= field\_def \{ "," field\_def \} [ "," ]}
\NormalTok{field\_def    ::= IDENT ":" type [ "=" expr ]}

\NormalTok{enum\_def     ::= [ "pub" ] "choice" IDENT "\{" variant\_defs "\}"}
\NormalTok{variant\_defs ::= variant\_def \{ "," variant\_def \} [ "," ]}
\NormalTok{variant\_def  ::= IDENT [ "(" type \{ "," type \} ")" ]}

\NormalTok{const\_def    ::= [ "pub" ] "const" IDENT [ ":" type ] "=" expr}
\NormalTok{type\_alias   ::= [ "pub" ] "type" IDENT "=" type}
\NormalTok{trait\_def    ::= "trait" IDENT [ "[" ... "]" ] "\{" \{ trait\_method \} "\}"}
\NormalTok{impl\_def     ::= "impl" [ "[" ... "]" ] [ IDENT "for" ] type "\{" \{ function\_def \} "\}"}
\NormalTok{extern\_def   ::= "extern" "def" IDENT "(" params ")" "{-}\textgreater{}" type [ effect\_clause ]}
\NormalTok{macro\_def    ::= "defmacro" IDENT "(" [ idents ] ")" "=\textgreater{}" expr}

\NormalTok{model\_def    ::= "model" IDENT "\{" \{ model\_item \} "\}"}
\NormalTok{model\_item   ::= "input" IDENT ":" type}
\NormalTok{               | "layer" IDENT IDENT [ "(" layer\_args ")" ]}
\NormalTok{               | "output" IDENT}

\NormalTok{block        ::= "\{" \{ stmt \} [ expr ] "\}"}
\NormalTok{stmt         ::= "val" ( IDENT [ ":" type ] | "(" IDENT \{ "," IDENT \} ")" ) "=" expr [ ";" ]}
\NormalTok{               | "var" IDENT [ ":" type ] "=" expr [ ";" ]}
\NormalTok{               | expr ( "=" | "+=" | "{-}=" | "*=" | "/=" ) expr [ ";" ]}
\NormalTok{               | "while" expr block [ ";" ]}
\NormalTok{               | "loop" block [ ";" ]}
\NormalTok{               | "for" IDENT "in" expr [ ".." expr [ "by" expr ] ] block [ ";" ]}
\NormalTok{               | "par" "for" IDENT "in" expr ".." expr block [ ";" ]}
\NormalTok{               | "spawn" [ "(" expr ")" ] block [ ";" ]}
\NormalTok{               | "return" [ expr ] [ ";" ]}
\NormalTok{               | "break" [ ";" ]}
\NormalTok{               | "continue" [ ";" ]}
\NormalTok{               | expr [ ";" ]}
\end{Highlighting}
\end{Shaded}

\begin{center}\rule{0.5\linewidth}{0.5pt}\end{center}

\chapter{Appendix B: Built-in Functions and Standard Library
Map}\label{appendix-b-built-in-functions-and-standard-library-map}

\section{Global Built-in Functions (No Import
Required)}\label{global-built-in-functions-no-import-required}

\begin{itemize}
\tightlist
\item
  \textbf{Core I/O \& Assertions}: \texttt{print}, \texttt{println},
  \texttt{read\_line}, \texttt{read\_i64}, \texttt{read\_f64},
  \texttt{assert}, \texttt{panic}, \texttt{type\_of}, \texttt{args}
\item
  \textbf{Conversions}: \texttt{to\_str}, \texttt{to\_i64},
  \texttt{to\_f64}, \texttt{format}, \texttt{parse\_i64},
  \texttt{parse\_f64}
\item
  \textbf{Math \& Checked/Wrapping Ops}: \texttt{sin}, \texttt{cos},
  \texttt{tan}, \texttt{exp}, \texttt{log}, \texttt{log2},
  \texttt{sqrt}, \texttt{abs}, \texttt{floor}, \texttt{ceil},
  \texttt{round}, \texttt{sign}, \texttt{pow}, \texttt{min},
  \texttt{max}, \texttt{clamp}, \texttt{math\_pi}, \texttt{math\_e},
  \texttt{math\_inf}, \texttt{is\_nan}, \texttt{is\_inf},
  \texttt{wrapping\_add}, \texttt{wrapping\_sub}, \texttt{wrapping\_mul}
\item
  \textbf{Bitwise}: \texttt{band}, \texttt{bor}, \texttt{bxor},
  \texttt{shl}, \texttt{shr}, \texttt{bitnot}
\item
  \textbf{Strings \& Regex}: \texttt{len}, \texttt{concat},
  \texttt{contains}, \texttt{starts\_with}, \texttt{ends\_with},
  \texttt{to\_upper}, \texttt{to\_lower}, \texttt{trim},
  \texttt{repeat}, \texttt{split}, \texttt{join}, \texttt{find},
  \texttt{slice}, \texttt{str\_index}, \texttt{str\_replace},
  \texttt{str\_reverse}, \texttt{char\_at}, \texttt{str\_pad\_left},
  \texttt{str\_pad\_right}, \texttt{str\_chars}, \texttt{str\_bytes},
  \texttt{str\_count}, \texttt{regex\_match}, \texttt{regex\_find\_all},
  \texttt{regex\_replace\_all}
\item
  \textbf{Collections}: \texttt{list}, \texttt{list\_get},
  \texttt{list\_set}, \texttt{list\_len}, \texttt{list\_push},
  \texttt{list\_pop}, \texttt{list\_map}, \texttt{list\_filter},
  \texttt{list\_reduce}, \texttt{list\_any}, \texttt{list\_all},
  \texttt{list\_zip}, \texttt{list\_enumerate}, \texttt{list\_flatten},
  \texttt{list\_unique}, \texttt{list\_reverse}, \texttt{list\_sorted},
  \texttt{list\_sum}, \texttt{list\_min}, \texttt{list\_max},
  \texttt{map}, \texttt{map\_get}, \texttt{map\_set},
  \texttt{map\_contains}, \texttt{map\_remove}, \texttt{map\_keys},
  \texttt{map\_values}, \texttt{map\_len}, \texttt{cell},
  \texttt{cell\_get}, \texttt{cell\_set}
\item
  \textbf{Option / Result / Autodiff / Time}: \texttt{some},
  \texttt{none}, \texttt{is\_some}, \texttt{unwrap}, \texttt{ok},
  \texttt{err}, \texttt{is\_ok}, \texttt{is\_err}, \texttt{tape},
  \texttt{backward}, \texttt{grad}, \texttt{sparsify}, \texttt{densify},
  \texttt{time\_now\_ms}, \texttt{sleep\_ms}
\item
  \textbf{Concurrency \& Atomics}: \texttt{channel}, \texttt{send},
  \texttt{recv}, \texttt{spawn}, \texttt{chan\_try\_recv},
  \texttt{chan\_len}, \texttt{select}, \texttt{timeout},
  \texttt{thread\_count}, \texttt{atomic}, \texttt{atomic\_load},
  \texttt{atomic\_store}, \texttt{atomic\_add}, \texttt{mutex},
  \texttt{mutex\_lock}, \texttt{mutex\_unlock}
\end{itemize}

\section{\texorpdfstring{Embedded Standard Library
(\texttt{bring\ std.\textless{}module\textgreater{}})}{Embedded Standard Library (bring std.\textless module\textgreater)}}\label{embedded-standard-library-bring-std.module}

IRIS ships \textbf{58 embedded \texttt{.iris} modules} (\textbf{59}
import names in \texttt{STDLIB\_MODULE\_NAMES}, where
\texttt{std.tensorx} aliases \texttt{std.tensor}). See
\href{stdlib-reference.md}{Standard Library Reference} for every
function signature.

{\def\LTcaptype{none} % do not increment counter
\begin{longtable}[]{@{}
  >{\raggedright\arraybackslash}p{(\linewidth - 2\tabcolsep) * \real{0.5000}}
  >{\raggedright\arraybackslash}p{(\linewidth - 2\tabcolsep) * \real{0.5000}}@{}}
\toprule\noalign{}
\begin{minipage}[b]{\linewidth}\raggedright
Domain
\end{minipage} & \begin{minipage}[b]{\linewidth}\raggedright
Modules
\end{minipage} \\
\midrule\noalign{}
\endhead
\bottomrule\noalign{}
\endlastfoot
\textbf{Desktop Native UI \& Media} & \texttt{ui}, \texttt{ui\_theme},
\texttt{ui\_protocol}, \texttt{ui\_layout}, \texttt{ui\_motion},
\texttt{ui\_icons}, \texttt{ui\_widgets} \\
\textbf{AI, ML, AD \& Autonomy} & \texttt{adaptive}, \texttt{ais},
\texttt{dataset}, \texttt{llm}, \texttt{meta\_learning}, \texttt{ml},
\texttt{nn}, \texttt{rl}, \texttt{stochastic}, \texttt{tensor}
(\texttt{tensorx}), \texttt{uncertainty} \\
\textbf{Data, Storage \& Visualization} & \texttt{csv},
\texttt{dataframe}, \texttt{json}, \texttt{kv}, \texttt{sql},
\texttt{svg}, \texttt{table}, \texttt{termplot} \\
\textbf{Collections, Strings \& Core} & \texttt{bitset},
\texttt{buffer}, \texttt{collections}, \texttt{container},
\texttt{deque}, \texttt{fmt}, \texttt{heap}, \texttt{iter},
\texttt{log}, \texttt{math}, \texttt{queue}, \texttt{set},
\texttt{string}, \texttt{testing}, \texttt{time}, \texttt{unicode} \\
\textbf{Systems, Networking \& Hardware} & \texttt{async},
\texttt{crypto}, \texttt{ffi}, \texttt{fs}, \texttt{http},
\texttt{http\_server}, \texttt{net}, \texttt{os}, \texttt{path},
\texttt{pcap}, \texttt{ros2}, \texttt{serial}, \texttt{tls} \\
\textbf{Safety, Evolution \& Reflection} & \texttt{meta},
\texttt{reflect}, \texttt{speculation} \\
\end{longtable}
}

\begin{center}\rule{0.5\linewidth}{0.5pt}\end{center}

\chapter{Appendix C: Type System and Operator Precedence
Reference}\label{appendix-c-type-system-and-operator-precedence-reference}

\section{Built-in Types}\label{built-in-types}

\begin{itemize}
\tightlist
\item
  \textbf{Scalars}: \texttt{i8}, \texttt{u8}, \texttt{i32},
  \texttt{u32}, \texttt{i64}, \texttt{u64}, \texttt{usize},
  \texttt{f32}, \texttt{f64}, \texttt{bool}, \texttt{str}, \texttt{unit}
\item
  \textbf{Composites \& Containers}:

  \begin{itemize}
  \tightlist
  \item
    Tuples \texttt{(T1,\ T2,\ ...)} and fixed stack arrays
    \texttt{{[}T;\ N{]}}
  \item
    \texttt{list\textless{}T\textgreater{}},
    \texttt{map\textless{}K,\ V\textgreater{}},
    \texttt{tensor\textless{}DTYPE,\ {[}DIMS{]}\textgreater{}},
    \texttt{sparse\textless{}T\textgreater{}},
    \texttt{grad\textless{}T\textgreater{}}
  \item
    \texttt{option\textless{}T\textgreater{}},
    \texttt{result\textless{}T,\ E\textgreater{}}
  \item
    \texttt{chan\textless{}T\textgreater{}},
    \texttt{atomic\textless{}T\textgreater{}},
    \texttt{mutex\textless{}T\textgreater{}}, \texttt{task\_group},
    \texttt{cell\textless{}T\textgreater{}}
  \item
    References \texttt{\&T}, \texttt{\&mut\ T}, dynamic traits
    \texttt{dyn\ Trait}, functions
    \texttt{(Args)\ -\textgreater{}\ Ret\ {[}effect\ ...{]}}
  \item
    Named \texttt{record} structs and \texttt{choice} algebraic enums
  \end{itemize}
\end{itemize}

\section{Operator Precedence (Highest to
Lowest)}\label{operator-precedence-highest-to-lowest}

{\def\LTcaptype{none} % do not increment counter
\begin{longtable}[]{@{}
  >{\raggedright\arraybackslash}p{(\linewidth - 6\tabcolsep) * \real{0.2500}}
  >{\raggedright\arraybackslash}p{(\linewidth - 6\tabcolsep) * \real{0.2500}}
  >{\raggedright\arraybackslash}p{(\linewidth - 6\tabcolsep) * \real{0.2500}}
  >{\raggedright\arraybackslash}p{(\linewidth - 6\tabcolsep) * \real{0.2500}}@{}}
\toprule\noalign{}
\begin{minipage}[b]{\linewidth}\raggedright
Precedence
\end{minipage} & \begin{minipage}[b]{\linewidth}\raggedright
Category
\end{minipage} & \begin{minipage}[b]{\linewidth}\raggedright
Operators
\end{minipage} & \begin{minipage}[b]{\linewidth}\raggedright
Notes
\end{minipage} \\
\midrule\noalign{}
\endhead
\bottomrule\noalign{}
\endlastfoot
1 (highest) & Postfix & \texttt{.field}, \texttt{.method()},
\texttt{.0}, \texttt{{[}idx{]}}, \texttt{?}, \texttt{!} (macro) &
Left-to-right \\
2 & Prefix unary & \texttt{-}, \texttt{!}, \texttt{not}, \texttt{\&},
\texttt{\&mut}, \texttt{*} & Right-to-left \\
3 & Matrix Contraction & \texttt{@} & Left-to-right tensor matmul \\
4 & Multiplicative & \texttt{*}, \texttt{/}, \texttt{\%} &
Left-to-right \\
5 & Additive & \texttt{+}, \texttt{-} & Left-to-right \\
6 & Coalescing / Cast & \texttt{to}, \texttt{??} & Explicit conversion /
default \\
7 & Comparison & \texttt{==}, \texttt{!=}, \texttt{\textless{}},
\texttt{\textless{}=}, \texttt{\textgreater{}}, \texttt{\textgreater{}=}
& Non-chaining \\
8 & Logical AND & \texttt{\&\&}, \texttt{and} & Short-circuiting \\
9 (lowest) & Logical OR &
\texttt{\textbackslash{}\textbar{}\textbackslash{}\textbar{}},
\texttt{or} & Short-circuiting \\
\end{longtable}
}

\begin{center}\rule{0.5\linewidth}{0.5pt}\end{center}

\chapter{Appendix D: Complete CLI
Reference}\label{appendix-d-complete-cli-reference}

The unified \texttt{iris} binary organizes \textbf{31 subcommands},
\textbf{16 \texttt{-\/-emit} modes}, and global compiler/execution
controls:

\section{Subcommands (31)}\label{subcommands-31}

\begin{itemize}
\tightlist
\item
  \textbf{Build \& Packaging}: \texttt{build}, \texttt{bundle},
  \texttt{run}, \texttt{embedded}, \texttt{new}, \texttt{pkg},
  \texttt{install}, \texttt{bindgen}
\item
  \textbf{Evolution \& Services}: \texttt{evolve},
  \texttt{evolve-unrestricted}, \texttt{evolve-search},
  \texttt{evolve-live}, \texttt{evolve-serve}, \texttt{service-daemon},
  \texttt{meta}
\item
  \textbf{Development, Quality \& Debugging}: \texttt{check},
  \texttt{test}, \texttt{watch}, \texttt{fmt}, \texttt{docs},
  \texttt{bench}, \texttt{profile}, \texttt{explain}, \texttt{doctor},
  \texttt{bug-report}
\item
  \textbf{Interactive \& Toolchain}: \texttt{repl}, \texttt{lsp},
  \texttt{dap}, \texttt{setup}, \texttt{upgrade}, \texttt{help}
\end{itemize}

\section{\texorpdfstring{Emit Kinds
(\texttt{-\/-emit\ \textless{}EMIT\textgreater{}})}{Emit Kinds (-\/-emit \textless EMIT\textgreater)}}\label{emit-kinds---emit-emit}

{\def\LTcaptype{none} % do not increment counter
\begin{longtable}[]{@{}
  >{\raggedright\arraybackslash}p{(\linewidth - 2\tabcolsep) * \real{0.5000}}
  >{\raggedright\arraybackslash}p{(\linewidth - 2\tabcolsep) * \real{0.5000}}@{}}
\toprule\noalign{}
\begin{minipage}[b]{\linewidth}\raggedright
Mode
\end{minipage} & \begin{minipage}[b]{\linewidth}\raggedright
Purpose
\end{minipage} \\
\midrule\noalign{}
\endhead
\bottomrule\noalign{}
\endlastfoot
\texttt{ir} (default) & Print SSA Intermediate Representation \\
\texttt{llvm} & Print lowered LLVM IR for the module \\
\texttt{llvm-complete} & Print complete LLVM IR including linked runtime
declarations \\
\texttt{eval} & Execute program natively (set
\texttt{IRIS\_FORCE\_INTERP=1} for interpreter execution) \\
\texttt{jit} & Execute in-process via LLVM ORC JIT \\
\texttt{binary} & Emit standalone native executable \\
\texttt{simd} & Emit LLVM IR with AVX2/NEON vectorized kernels \\
\texttt{pgo-instrument} / \texttt{pgo-optimize} & Emit instrumented or
PGO-optimized LLVM IR (\texttt{-\/-pgo-generate} /
\texttt{-\/-pgo-use}) \\
\texttt{cuda} / \texttt{cuda-ptx} & Emit NVIDIA CUDA LLVM IR or PTX
device code (\texttt{-\/-gpu\ cuda} runs numeric/tensor
\texttt{par\ for}; PTX \& CPU fallback tested, not yet run on NVIDIA GPU
hardware) \\
\texttt{graph} & Export DOT dependency / model topology graph \\
\texttt{onnx} / \texttt{onnx-binary} & Export declarative \texttt{model}
graph as text or ONNX protobuf (\texttt{-\/-weights\ FILE}) \\
\texttt{tensorrt} & Build serialized TensorRT engine via
\texttt{trtexec} (CLI invocation tested, not yet run where TensorRT is
installed) \\
\texttt{python-ext} & Compile module as a CPython extension
(\texttt{.pyd} / \texttt{.so}) \\
\end{longtable}
}

\section{Key Global \& Build Options}\label{key-global-build-options}

\texttt{-\/-output\ \textless{}FILE\textgreater{}} (\texttt{-o}),
\texttt{-\/-error-format\ human\textbar{}json},
\texttt{-\/-target\ \textless{}TRIPLE\textgreater{}},
\texttt{-\/-dump-ir-after\ \textless{}PASS\textgreater{}},
\texttt{-\/-max-steps\ \textless{}N\textgreater{}},
\texttt{-\/-max-depth\ \textless{}N\textgreater{}},
\texttt{-\/-no-cache}, \texttt{-\/-sandbox},
\texttt{-\/-strict-effects}, \texttt{-\/-simd},
\texttt{-\/-pgo-generate},
\texttt{-\/-pgo-use\ \textless{}FILE\textgreater{}},
\texttt{-\/-weights\ \textless{}FILE\textgreater{}},
\texttt{-\/-gpu\ cuda}, \texttt{-\/-debug} (\texttt{-g}),
\texttt{-\/-define\ NAME{[}=VALUE{]}} (\texttt{-D}),
\texttt{iris\ build\ -\/-gui}, and
\texttt{iris\ bundle\ -\/-format\ auto\textbar{}dir\textbar{}zip\textbar{}inno\textbar{}app\textbar{}appimage\ -\/-release\ -\/-install\ -\/-uninstall}.

\begin{center}\rule{0.5\linewidth}{0.5pt}\end{center}

\chapter{Appendix E: Compiler Diagnostic and Warning
Reference}\label{appendix-e-compiler-diagnostic-and-warning-reference}

Run \texttt{iris\ explain\ \textless{}CODE\textgreater{}} for detailed
inline help and fixes on any code:

{\def\LTcaptype{none} % do not increment counter
\begin{longtable}[]{@{}
  >{\raggedright\arraybackslash}p{(\linewidth - 4\tabcolsep) * \real{0.3333}}
  >{\raggedright\arraybackslash}p{(\linewidth - 4\tabcolsep) * \real{0.3333}}
  >{\raggedright\arraybackslash}p{(\linewidth - 4\tabcolsep) * \real{0.3333}}@{}}
\toprule\noalign{}
\begin{minipage}[b]{\linewidth}\raggedright
Range
\end{minipage} & \begin{minipage}[b]{\linewidth}\raggedright
Code
\end{minipage} & \begin{minipage}[b]{\linewidth}\raggedright
Diagnostic Title
\end{minipage} \\
\midrule\noalign{}
\endhead
\bottomrule\noalign{}
\endlastfoot
\textbf{Lexer \& Parser} & \texttt{E0001}--\texttt{E0007} &
\texttt{E0001} Unexpected Character, \texttt{E0002} Unterminated String,
\texttt{E0003} Invalid Escape Sequence, \texttt{E0004} Invalid Literal,
\texttt{E0005} Unexpected Token, \texttt{E0006} Unclosed Delimiter,
\texttt{E0007} Nesting Too Deep \\
\textbf{Lowering \& Type Resolution} & \texttt{E0100}--\texttt{E0129} &
\texttt{E0100} Undefined Variable/Function, \texttt{E0101} Type
Mismatch, \texttt{E0102} Duplicate Definition, \texttt{E0103}
Unsupported Construct, \texttt{E0104}--\texttt{E0107} Model DSL Layer/Op
Errors, \texttt{E0108} Rejected Construct, \texttt{E0109} Wrong Arg
Count, \texttt{E0110} Arg Type Mismatch, \texttt{E0111} Unknown Type,
\texttt{E0112} Assign to Immutable \texttt{val},
\texttt{E0113}--\texttt{E0115} Field/Module Errors, \texttt{E0116}
Missing Return, \texttt{E0117}--\texttt{E0119}
Method/Callable/\texttt{if}-without-\texttt{else},
\texttt{E0122}--\texttt{E0129} Variant/Loop/Format/\texttt{when}/Trait
Impl Errors \\
\textbf{IR Passes \& Borrow Safety} & \texttt{E0200}--\texttt{E0210},
\texttt{E0382} & \texttt{E0200} Use Before Definition, \texttt{E0201}
Multiple Definition, \texttt{E0202} Pass Type Error, \texttt{E0203}
Missing Terminator, \texttt{E0204} Shape Mismatch, \texttt{E0205}
Unresolved Type, \texttt{E0206} Const Eval Error, \texttt{E0207} Type
Annotation Needed, \texttt{E0208} Branch Type Mismatch, \texttt{E0209}
Wrong Return Type, \texttt{E0210} Non-exhaustive \texttt{when},
\texttt{E0382} Use After Move \\
\textbf{Codegen \& Effects} & \texttt{E0300}--\texttt{E0304} &
\texttt{E0300} Codegen Failed, \texttt{E0301} No Native Toolchain,
\texttt{E0302} Program Failed, \texttt{E0303} Unavailable on Target,
\texttt{E0304} Effect Bound Exceeded (\texttt{-\/-strict-effects}) \\
\textbf{Runtime Traps} & \texttt{E0400}--\texttt{E0410} & \texttt{E0400}
Runtime Error, \texttt{E0401} Undefined Value, \texttt{E0402} Division
By Zero, \texttt{E0403} Index Out of Bounds, \texttt{E0404} Runtime Type
Error, \texttt{E0405} Runtime Unsupported, \texttt{E0406}
Panic/\texttt{assert}/\texttt{unwrap}, \texttt{E0407} Security
Violation, \texttt{E0408} Stack Overflow, \texttt{E0409} Integer
Overflow, \texttt{E0410} GPU Error \\
\textbf{System \& Internal} & \texttt{E0500}--\texttt{E0902} &
\texttt{E0500} I/O Error, \texttt{E0501} Preprocessor Error,
\texttt{E0900}--\texttt{E0902} Internal Compiler Errors
(\texttt{iris\ bug-report}) \\
\end{longtable}
}

The compiler and LSP (\texttt{find\_source\_warnings}) also flag unused
variables, parameters, private unused functions, unreachable code,
unsafe blocks, and loop conditions whose variables never mutate in-body.

\begin{center}\rule{0.5\linewidth}{0.5pt}\end{center}

\textbf{Edition}: IRIS \texttt{1.0.0-rc2} \textbf{Toolchain}: Tested on
Windows (LLVM 23.1.1 + MinGW UCRT64 \texttt{ld.lld}), Linux
(\texttt{x86\_64}), and macOS (\texttt{aarch64}) \textbf{License}: GNU
General Public License v2.0 or later --- see \href{../LICENSE}{LICENSE}
\textbf{Source}:
\href{https://github.com/Eclipse-Softworks/IRIS}{github.com/Eclipse-Softworks/IRIS}

\backmatter
\end{document}
